% Contraction de texte : reformulation des idées d’un texte à résumer — Français Tle Tech — Cours signé OnBuch+
% Programme officiel MINESEC. Compiler avec : tectonic N09-contraction-texte-reformulation-idees-texte-resumer.tex
\newcommand{\DOCMATIERE}{Français}
\newcommand{\DOCNIVEAU}{Tle Tech}
\newcommand{\DOCMODULE}{Français – classe de Terminale technique}
\newcommand{\DOCLECON}{N09}
\newcommand{\DOCTITRE}{Contraction de texte : reformulation des idées d’un texte à résumer}
% OnBuch+ — préambule « cours signés » ENRICHI (compilé avec tectonic / XeLaTeX)
% Système visuel « Pop » d'OnBuch+ : fond crème (jamais blanc), bordures encre
% épaisses, ombres « dures » décalées, titres Archivo Black en capitales.
% Version MATHÉMATIQUES : graphes (pgfplots/tikz), boîtes pédagogiques
% (définition, propriété, méthode, attention, le savais-tu, exercice résolu,
% exercice, TP, bilan), page de garde et signature OnBuch+ ancrée en bas.
%
% Macros attendues dans main.tex AVANT \input{preamble} :
%   \DOCMATIERE  \DOCNIVEAU  \DOCMODULE  \DOCLECON (numéro)  \DOCTITRE
\documentclass[11pt]{article}

\usepackage{fontspec}
\usepackage[french]{babel}
\usepackage[a4paper,margin=2cm,top=3.4cm,bottom=2.8cm,headheight=1.4cm,footskip=1.3cm]{geometry}
\usepackage{amsmath,amssymb,amsfonts}
\usepackage{mathtools} % \xmapsto, \coloneqq... souvent utilisés par les modèles
\usepackage{cancel} % simplifications barrées (bilans de forces)
% \abs{z} (module, valeur absolue) et \norm{u} (norme), utilisés par les modèles
\providecommand{\abs}{}\let\abs\relax\DeclarePairedDelimiter{\abs}{\lvert}{\rvert}
\providecommand{\norm}{}\let\norm\relax\DeclarePairedDelimiter{\norm}{\lVert}{\rVert}
\usepackage[table]{xcolor} % [table] : fournit aussi \rowcolors (alternance de lignes)
\usepackage{pagecolor}
\usepackage{tikz}
\usetikzlibrary{arrows.meta,positioning,calc,shapes.geometric,shapes.misc,shapes.arrows,decorations.pathmorphing,decorations.pathreplacing,decorations.markings,patterns,shadows,mindmap,trees,fit,backgrounds,matrix,intersections,angles,quotes}
\usepackage{pgfplots}
\usepackage[version=4]{mhchem} % équations nucléaires \ce{^{235}_{92}U}
\usepackage[european,siunitx]{circuitikz} % schémas électriques
\pgfplotsset{compat=1.18}
\usepgfplotslibrary{fillbetween,groupplots} % aires entre courbes (calcul intégral)
\usepackage{enumitem}
\usepackage{fancyhdr}
% [most] charge toutes les bibliothèques tcolorbox (listings, minted, raster,
% documentation...) et gonfle inutilement la mémoire TeX sur un document long
% (~300 boîtes) : on ne charge que ce qui est réellement utilisé.
\usepackage[skins,breakable]{tcolorbox}
\usepackage{graphicx}
\usepackage{colortbl}
\usepackage{booktabs}
\usepackage{tabularx}
\usepackage{multirow}
\usepackage{diagbox} % case d'en-tête diagonale des tableaux à double entrée
% Colonnes extensibles centrée / gauche / droite (C, L, R), souvent utilisées par les modèles
\newcolumntype{C}{>{\centering\arraybackslash}X}
\newcolumntype{L}{>{\raggedright\arraybackslash}X}
\newcolumntype{R}{>{\raggedleft\arraybackslash}X}
\usepackage{array}
\usepackage{multicol}
\usepackage{siunitx}
% parse-numbers=false : accepte aussi \SI{2,5 \times 10^{-3}}{...}
\sisetup{locale=FR,output-decimal-marker={,},group-minimum-digits=4,parse-numbers=false}
\DeclareSIUnit\ohm{\ensuremath{\symup{\Omega}}} % U+2126 absent de la police
\DeclareSIUnit\division{div} % réglages d'oscilloscope (ms/div, V/div)
\DeclareSIUnit\tr{tr} % tours (vitesse de rotation, ex. \SI{1500}{\tr\per\minute})
\DeclareSIUnit\molar{M} % concentration molaire (ex. \SI{3}{\molar}, \SI{1}{\micro\molar})
\DeclareSIUnit\an{an} % année (le modèle confond parfois avec \year, un compteur TeX) — ex. \SI{5}{\kilo\meter\per\an}
\DeclareSIUnit\FCFA{FCFA} % franc CFA (ex. \SI{500}{\FCFA\per\kilo\gram})
\DeclareSIUnit\degreeBrix{\textdegree Brix} % teneur en sucre d'un sirop/jus (réfractométrie)
\DeclareSIUnit\month{mois} % le modèle utilise parfois \month, un compteur TeX, comme unité de durée
\DeclareSIUnit\microliter{\micro\liter} % le modèle écrit parfois \microliter en un seul token
\DeclareSIUnit\million{millions} % dénombrement (ex. \SI{15}{\million\per\mL} spermatozoïdes)
\usepackage{qrcode}
% \degree : commande fréquemment utilisée par les modèles pour un angle ou une
% température en dehors de \SI{}{} (non fournie par les paquets chargés).
\providecommand{\degree}{^{\circ}}
% \qed (fin de démonstration), utilisé par les modèles sans amsthm
\providecommand{\qed}{\hfill$\square$}
% \barre : double barre des valeurs interdites dans un tableau de variations (tkz-tab)
\providecommand{\barre}{\ensuremath{\|}}
% environnement proof (amsthm non chargé) utilisé par les modèles
\ifdefined\proof\else\newenvironment{proof}[1][Démonstration]{\par\noindent\textit{#1.}\ }{\qed\par}\fi
\usepackage{lastpage}
\usepackage{hyperref}
\hypersetup{hidelinks}

% ============================================================
% Palette « Pop » OnBuch+ — mêmes hex que la classe P (plus_ui.dart)
% ============================================================
\definecolor{poporange}{HTML}{F59321}
\definecolor{popdark}{HTML}{E07A0C}
\definecolor{popcream}{HTML}{FFF6E8}
\definecolor{popink}{HTML}{1C1714}
\definecolor{poppurple}{HTML}{7A5AE0}
\definecolor{popgreen}{HTML}{2E9E6A}
\definecolor{poppink}{HTML}{E0457B}
\definecolor{popblue}{HTML}{2F7DE1}
\definecolor{popgold}{HTML}{C98A12}
\definecolor{popmuted}{HTML}{8A7F76}
\definecolor{popline}{HTML}{EADFD2}
\definecolor{popgray}{HTML}{8A7F76}
\colorlet{popgrey}{popgray}
% Alias tolérants (le modèle nomme parfois les couleurs différemment)
\colorlet{popwhite}{white}
\colorlet{popblack}{popink}
\colorlet{popred}{poppink}
\colorlet{popyellow}{popgold}
% Teintes claires (fonds de tableaux/encadrés) — le modèle nomme parfois
% les couleurs avec un suffixe L (ex. popblueL) sans les avoir définies.
\colorlet{poporangeL}{poporange!20}
\colorlet{popdarkL}{popdark!20}
\colorlet{poppurpleL}{poppurple!20}
\colorlet{popgreenL}{popgreen!20}
\colorlet{poppinkL}{poppink!20}
\colorlet{popblueL}{popblue!20}
\colorlet{popgoldL}{popgold!20}
\colorlet{popredL}{popred!20}
\colorlet{popgrayL}{popgray!20}
% Teintes claires pour fonds de tableaux / zones
\colorlet{popblueL}{popblue!10!white}
\colorlet{poporangeL}{poporange!14!white}
\colorlet{popgreenL}{popgreen!12!white}
\colorlet{poppurpleL}{poppurple!12!white}
\colorlet{poppinkL}{poppink!12!white}
\colorlet{popgoldL}{popgold!14!white}

\pagecolor{popcream}

% ============================================================
% Typographie
% ============================================================
\defaultfontfeatures{Ligatures=TeX}
\setmainfont{PlusJakartaSans-Medium.ttf}[Path=../../../fonts/,BoldFont=PlusJakartaSans-Bold.ttf]
% Maths en sans-serif (Fira Math) avec chiffres et lettres droites de Plus
% Jakarta, pour que les formules se fondent dans le texte.
\usepackage[math-style=ISO,bold-style=ISO]{unicode-math}
\setmathfont{FiraMath-Regular.otf}[Path=../../../fonts/]
\setmathfont{PlusJakartaSans-Medium.ttf}[Path=../../../fonts/,range={up/num,up/{Latin,latin},\mathup}]
\usepackage{icomma} % virgule décimale sans espace en mode math
% Caractères mathématiques Unicode absents des polices de texte (Plus Jakarta,
% Archivo Black) : rendus par leur équivalent mathématique, en texte comme en
% formule — sinon ils s'affichent en carré vide (ex. « Arithmétique dans ℤ »).
\usepackage{newunicodechar}
\newunicodechar{ℝ}{\ensuremath{\mathbb{R}}}
\newunicodechar{ℤ}{\ensuremath{\mathbb{Z}}}
\newunicodechar{ℂ}{\ensuremath{\mathbb{C}}}
\newunicodechar{ℕ}{\ensuremath{\mathbb{N}}}
\newunicodechar{ℚ}{\ensuremath{\mathbb{Q}}}
\newunicodechar{𝔻}{\ensuremath{\mathbb{D}}}
\newunicodechar{α}{\ensuremath{\alpha}}
\newunicodechar{β}{\ensuremath{\beta}}
\newunicodechar{γ}{\ensuremath{\gamma}}
\newunicodechar{δ}{\ensuremath{\delta}}
\newunicodechar{ε}{\ensuremath{\varepsilon}}
\newunicodechar{θ}{\ensuremath{\theta}}
\newunicodechar{λ}{\ensuremath{\lambda}}
\newunicodechar{μ}{\ensuremath{\mu}}
\newunicodechar{σ}{\ensuremath{\sigma}}
\newunicodechar{τ}{\ensuremath{\tau}}
\newunicodechar{φ}{\ensuremath{\varphi}}
\newunicodechar{ψ}{\ensuremath{\psi}}
\newunicodechar{ω}{\ensuremath{\omega}}
\newunicodechar{Σ}{\ensuremath{\Sigma}}
% majuscules produites par \MakeUppercase dans les titres (α-aminés -> Α)
\newunicodechar{Α}{\ensuremath{\alpha}}
\newunicodechar{Β}{\ensuremath{\beta}}
\newunicodechar{Γ}{\ensuremath{\Gamma}}
\newunicodechar{Δ}{\ensuremath{\Delta}}
\newunicodechar{Ε}{\ensuremath{\varepsilon}}
\newunicodechar{Θ}{\ensuremath{\Theta}}
\newunicodechar{Λ}{\ensuremath{\Lambda}}
\newunicodechar{Μ}{\ensuremath{\mu}}
\newunicodechar{Τ}{\ensuremath{\tau}}
\newunicodechar{Φ}{\ensuremath{\Phi}}
\newunicodechar{Ψ}{\ensuremath{\Psi}}
\newunicodechar{Ω}{\ensuremath{\Omega}}
\newunicodechar{↦}{\ensuremath{\mapsto}}
\newunicodechar{∈}{\ensuremath{\in}}
\newunicodechar{∉}{\ensuremath{\notin}}
\newunicodechar{∧}{\ensuremath{\wedge}}
\newunicodechar{⇔}{\ensuremath{\Leftrightarrow}}
\newunicodechar{⟺}{\ensuremath{\Longleftrightarrow}}
\newunicodechar{⇒}{\ensuremath{\Rightarrow}}
\newunicodechar{⟹}{\ensuremath{\Longrightarrow}}
\newunicodechar{∘}{\ensuremath{\circ}}
\newunicodechar{∪}{\ensuremath{\cup}}
\newunicodechar{∩}{\ensuremath{\cap}}
\newunicodechar{≡}{\ensuremath{\equiv}}
\newunicodechar{⊂}{\ensuremath{\subset}}
\newunicodechar{⊥}{\ensuremath{\perp}}
\newunicodechar{∃}{\ensuremath{\exists}}
\newunicodechar{∀}{\ensuremath{\forall}}
\newunicodechar{∛}{\ensuremath{\sqrt[3]{\,}}}
\newunicodechar{∜}{\ensuremath{\sqrt[4]{\,}}}
\newunicodechar{⃗}{\ensuremath{{}^{\to}}}
\newunicodechar{ⁿ}{\ifmmode^{n}\else\textsuperscript{\ensuremath{n}}\fi}
\newunicodechar{ˣ}{\ifmmode^{x}\else\textsuperscript{\ensuremath{x}}\fi}
\newunicodechar{ᵇ}{\ifmmode^{b}\else\textsuperscript{\ensuremath{b}}\fi}
\newunicodechar{ʳ}{\ifmmode^{r}\else\textsuperscript{\ensuremath{r}}\fi}
\newunicodechar{ᵘ}{\ifmmode^{u}\else\textsuperscript{\ensuremath{u}}\fi}
\newunicodechar{ʸ}{\ifmmode^{y}\else\textsuperscript{\ensuremath{y}}\fi}
\newunicodechar{ᵃ}{\ifmmode^{a}\else\textsuperscript{\ensuremath{a}}\fi}
\newunicodechar{ᵉ}{\ifmmode^{e}\else\textsuperscript{\ensuremath{e}}\fi}
\newunicodechar{ᵏ}{\ifmmode^{k}\else\textsuperscript{\ensuremath{k}}\fi}
\newunicodechar{ᵖ}{\ifmmode^{p}\else\textsuperscript{\ensuremath{p}}\fi}
\newunicodechar{ᴺ}{\ifmmode^{N}\else\textsuperscript{\ensuremath{N}}\fi}
\newunicodechar{ᶜ}{\ifmmode^{c}\else\textsuperscript{\ensuremath{c}}\fi}
\newunicodechar{ʰ}{\ifmmode^{h}\else\textsuperscript{\ensuremath{h}}\fi}
\newunicodechar{ᵠ}{\ifmmode^{q}\else\textsuperscript{\ensuremath{q}}\fi}
\newunicodechar{ₙ}{\ifmmode_{n}\else\textsubscript{\ensuremath{n}}\fi}
\newunicodechar{ₐ}{\ifmmode_{a}\else\textsubscript{\ensuremath{a}}\fi}
\newunicodechar{ᵢ}{\ifmmode_{i}\else\textsubscript{\ensuremath{i}}\fi}
\newunicodechar{ᵣ}{\ifmmode_{r}\else\textsubscript{\ensuremath{r}}\fi}
\newunicodechar{ₖ}{\ifmmode_{k}\else\textsubscript{\ensuremath{k}}\fi}
\newunicodechar{ᵦ}{\ifmmode_{\beta}\else\textsubscript{\ensuremath{\beta}}\fi}
\newunicodechar{ₓ}{\ifmmode_{x}\else\textsubscript{\ensuremath{x}}\fi}
\newfontfamily\popdisplay{ArchivoBlack-Regular.ttf}[Path=../../../fonts/]
\newfontfamily\poplogo{SpaceGrotesk-Bold.ttf}[Path=../../../fonts/]
\newfontfamily\popsemi{PlusJakartaSans-SemiBold.ttf}[Path=../../../fonts/]
\newfontfamily\popxbold{PlusJakartaSans-ExtraBold.ttf}[Path=../../../fonts/]
\newfontfamily\popxboldls{PlusJakartaSans-ExtraBold.ttf}[Path=../../../fonts/,LetterSpace=3.0]

\setlist[enumerate]{leftmargin=1.7em,itemsep=5pt,topsep=4pt}
\setlist[itemize]{leftmargin=1.7em,itemsep=4pt,topsep=4pt,label={\color{poporange}\textbullet}}
\setlength{\parindent}{0pt}
\setlength{\parskip}{5pt}
\renewcommand{\arraystretch}{1.25}

% pgfplots : style Pop par défaut
\pgfplotsset{
  every axis/.append style={axis line style={popink,line width=1pt},tick style={popink},
    grid style={popline,line width=0.5pt},label style={font=\small},tick label style={font=\footnotesize}},
  popcurve/.style={poporange,line width=1.8pt,no markers,smooth},
}
% Même style disponible dans un \draw TikZ simple (hors pgfplots)
\tikzset{popcurve/.style={poporange,line width=1.8pt}}
\tikzset{popgrid/.style={popline,very thin}}
\colorlet{popcurve}{poporange} % le nom du style est parfois utilisé comme couleur

% ============================================================
% Pill / squiggle
% ============================================================
\newcommand{\poppill}[3]{%
  \tikz[baseline=-0.55ex]\node[draw=popink,line width=1.2pt,rounded corners=2.6pt,%
    fill=#2,inner xsep=6.5pt,inner ysep=3pt]{%
    \popxboldls\fontsize{7}{7}\selectfont\color{#3}\MakeUppercase{#1}};%
}
\newcommand{\popsquiggle}[1]{%
  \tikz[baseline=-0.4ex]{
    \draw[#1,line width=2.6pt,line cap=round]
      plot [smooth] coordinates {(0,0.09) (0.34,0.01) (0.68,0.21) (1.02,0.01) (1.36,0.10)};
  }%
}
% Mot-clé surligné (marqueur orange sous le texte)
\newcommand{\cle}[1]{{\popxbold\color{popdark}#1}}

% ============================================================
% En-tête / pied
% ============================================================
\pagestyle{fancy}
\fancyhf{}
\renewcommand{\headrulewidth}{0pt}
\renewcommand{\footrulewidth}{0pt}
\lhead{\raisebox{-0.15ex}{\poplogo\fontsize{13}{13}\selectfont\color{popink}OnBuch\color{poporange}+}}
\rhead{\poppill{\DOCMATIERE\ \textbullet\ \DOCNIVEAU\ \textbullet\ Leçon \DOCLECON}{white}{popink}}
\lfoot{\popxbold\fontsize{7.6}{7.6}\selectfont\color{popmuted}\MakeUppercase{Cours signé OnBuch+}\ \textbullet\ {\popsemi\DOCTITRE}}
\rfoot{\tikz[baseline=-0.6ex]\node[rounded corners=8.5pt,fill=popink,minimum height=17pt,minimum width=17pt,inner xsep=5pt,inner sep=0]{\popxbold\fontsize{8}{8}\selectfont\color{popcream}\thepage\,/\,\pageref*{LastPage}};}

% ============================================================
% Titres
% ============================================================
\newcommand{\chaptitre}[2]{%
  \begin{tcolorbox}[enhanced,colback=poporange,colframe=popink,boxrule=2.6pt,arc=11pt,
      drop shadow=popink,left=15pt,right=15pt,top=12pt,bottom=11pt,before skip=3pt,after skip=18pt]
    \poppill{#2}{popink}{popcream}\par\vspace{9pt}
    {\raggedright\hyphenpenalty=10000\exhyphenpenalty=10000\popdisplay\fontsize{22}{24}\selectfont\color{popcream}\MakeUppercase{#1}\par}
  \end{tcolorbox}
}
% Section numérotée : pastille orange avec le numéro + titre Archivo
\newcounter{popsec}
\newcommand{\coursec}[1]{%
  \stepcounter{popsec}\setcounter{popsub}{0}%
  \par\vspace{16pt}\noindent
  \begin{minipage}{\linewidth}
  \tikz[baseline=-0.7ex]\node[draw=popink,line width=1.6pt,fill=poporange,rounded corners=4pt,minimum size=22pt,inner sep=0]{\popdisplay\fontsize{12}{12}\selectfont\color{popcream}\thepopsec};%
  \hspace{8pt}{\popdisplay\fontsize{15}{17}\selectfont\color{popink}\MakeUppercase{#1}}\par
  \vspace{4pt}\noindent{\color{poporange}\rule{36pt}{3.6pt}}
  \end{minipage}\par\nopagebreak\vspace{8pt}
}
\newcounter{popsub}
\newcommand{\courssub}[1]{%
  \stepcounter{popsub}%
  \par\vspace{9pt}\noindent{\popxbold\fontsize{11.5}{13}\selectfont\color{popink}{\color{poporange}\thepopsec.\thepopsub}\hspace{6pt}#1}\par\nopagebreak\vspace{4pt}
}

% ============================================================
% Boîtes pédagogiques — même recette : fond blanc, bordure épaisse colorée,
% ombre dure encre, badge plein. Titre optionnel [#1].
% ============================================================
\tcbset{popbase/.style={breakable,enhanced,colback=white,boxrule=2.3pt,arc=9pt,
  shadow={3pt}{-3pt}{0pt}{popink},left=14pt,right=14pt,top=11pt,bottom=11pt,before skip=11pt,after skip=15pt,
  fontupper=\popsemi\fontsize{10.6}{14.5}\selectfont\color{popink},
  fontlower=\popsemi\fontsize{10.6}{14.5}\selectfont\color{popink}}}
% Coche dessinée (le glyphe ✓ n'existe pas dans les polices du document)
\newcommand{\popcheck}[1]{\tikz[baseline=-0.5ex]\draw[#1,line width=1.6pt,line cap=round,line join=round](-0.13,0.0)--(-0.04,-0.09)--(0.13,0.1);}
\providecommand{\coche}{\popcheck{popgreen}}
\providecommand{\resultat}[1]{\par\smallskip\noindent\fcolorbox{popgreen}{popgreenL}{\parbox{\dimexpr\linewidth-2\fboxsep-2\fboxrule\relax}{\textbf{Résultat.} #1}}\par\smallskip}
\newcommand{\popboxhead}[3]{\poppill{#1}{#2}{white}\ifx\relax#3\relax\else\hspace{6pt}{\popxbold\fontsize{10.6}{12}\selectfont\color{popink}#3}\fi\par\vspace{7pt}}

\newtcolorbox{aretenir}[1][]{popbase,colframe=popgreen,before upper={\popboxhead{À retenir}{popgreen}{#1}}}
\newtcolorbox{exemplebox}[1][]{popbase,colframe=poporange,before upper={\popboxhead{Exemple}{poporange}{#1}}}
\newtcolorbox{definition}[1][]{popbase,colframe=popblue,before upper={\popboxhead{Définition}{popblue}{#1}}}
\newtcolorbox{propriete}[1][]{popbase,colframe=poppurple,before upper={\popboxhead{Propriété}{poppurple}{#1}}}
\newtcolorbox{methode}[1][]{popbase,colframe=popgold,before upper={\popboxhead{Méthode}{popgold}{#1}}}
\newtcolorbox{attention}[1][]{popbase,colframe=poppink,before upper={\popboxhead{Attention}{poppink}{#1}}}
\newtcolorbox{experience}[1][]{popbase,colframe=popblue,colback=popblueL,before upper={\popboxhead{Expérience}{popblue}{#1}}}
% « Le savais-tu ? » : carte violette pleine (contexte, culture, Cameroun)
\newtcolorbox{savaistu}[1][]{popbase,colback=poppurple,colframe=popink,
  fontupper=\popsemi\fontsize{10.4}{14.2}\selectfont\color{white},
  before upper={\poppill{Le savais-tu\,?}{white}{poppurple}\ifx\relax#1\relax\else\hspace{6pt}{\popxbold\color{white}#1}\fi\par\vspace{7pt}}}
% Exercice résolu : énoncé en haut, \tcblower puis la solution sur fond crème
\newcounter{popexo}\newcounter{popexr}
\newtcolorbox{exoresolu}[1][]{popbase,colframe=poporange,colbacklower=popcream,
  segmentation style={popink,line width=1.2pt,dashed},
  before upper={\stepcounter{popexr}\popboxhead{Exercice résolu \thepopexr}{poporange}{#1}},
  before lower={\poppill{Solution}{popink}{popcream}\par\vspace{6pt}}}
% Exercice (non résolu en place) — la correction est dans la partie Corrigés
\newtcolorbox{exercice}[1][]{popbase,colframe=popink,
  before upper={\stepcounter{popexo}\popboxhead{Exercice \thepopexo}{popink}{#1}}}
% Corrigé
\newtcolorbox{corrige}[1][]{popbase,colframe=popgreen,colback=popgreenL,
  before upper={\popboxhead{Corrigé}{popgreen}{#1}}}
% Objectifs / prérequis / bilan
\newtcolorbox{objectifs}{popbase,colframe=popink,colback=poporangeL,
  before upper={\popboxhead{Objectifs de la leçon}{popink}{}}}
\newtcolorbox{prerequis}{popbase,colframe=popink,colback=popblueL,
  before upper={\popboxhead{Prérequis}{popblue}{}}}
\newtcolorbox{bilan}{popbase,colframe=popink,colback=white,boxrule=2.8pt,
  before upper={\popboxhead{Fiche bilan}{poporange}{L'essentiel à connaître pour l'examen}}}
% Figure encadrée avec légende
\newtcolorbox{popfigure}[1][]{enhanced,breakable=false,colback=white,colframe=popline,boxrule=1.4pt,arc=8pt,
  left=10pt,right=10pt,top=10pt,bottom=8pt,before skip=10pt,after skip=12pt,halign=center,
  lower separated=false,halign lower=center,
  fontlower=\popsemi\fontsize{9}{11.5}\selectfont\color{popmuted},
  #1}
\newcommand{\legende}[1]{\par\vspace{6pt}{\popsemi\fontsize{9}{11.5}\selectfont\color{popmuted}#1}}

% ============================================================
% Signature OnBuch+
% ============================================================
\newcommand{\poplogoinline}[1]{{\poplogo\fontsize{#1}{#1}\selectfont\color{popink}OnBuch\color{poporange}+}}
\newcommand{\signatureonbuch}{%
  \begin{tcolorbox}[enhanced,colback=popink,colframe=popink,boxrule=2.2pt,arc=10pt,
      drop shadow=poporange,left=14pt,right=14pt,top=10pt,bottom=10pt,before skip=0pt,after skip=0pt]
    \tikz[baseline=-0.7ex]\node[circle,fill=popgreen,draw=popcream,line width=1.3pt,minimum size=22pt,inner sep=0]{\popcheck{white}};%
    \hspace{9pt}\begin{minipage}[c]{0.62\linewidth}
      {\popxbold\fontsize{12}{13}\selectfont\color{popcream}Signé\ \textbullet\ {\poplogo OnBuch\color{poporange}+}}\par\vspace{2pt}
      {\raggedright\popsemi\fontsize{8}{10.5}\selectfont\color{popline}\MakeUppercase{Équipe pédagogique OnBuch+}\\\MakeUppercase{Conforme au programme officiel MINESEC}\par}
    \end{minipage}\hfill
    {\poplogo\fontsize{22}{22}\selectfont\color{popcream}OnBuch\color{poporange}+}
  \end{tcolorbox}%
}

% ============================================================
% Page de garde
% #1 titre · #2 matière · #3 niveau · #4 module · #5 n° leçon · #6 durée ·
% #7 description / objectifs (texte ou itemize)
% ============================================================
\newcommand{\pagegarde}[7]{%
  \thispagestyle{empty}
  \newgeometry{a4paper,margin=1.7cm,top=1.6cm,bottom=1.4cm}%
  \begin{tikzpicture}[remember picture,overlay]
    \fill[poporange!18] ([xshift=-3.2cm,yshift=-3.6cm]current page.north east) circle (4.6cm);
    \fill[poppurple!14] ([xshift=2.4cm,yshift=7.6cm]current page.south west) circle (3.8cm);
  \end{tikzpicture}%
  {\poplogo\fontsize{30}{30}\selectfont\color{popink}OnBuch\color{poporange}+}\hfill
  \raisebox{0.6ex}{\poppill{Cours signé \textbullet\ Premium}{popink}{popcream}}\par
  \vspace{1pt}\popsquiggle{poppurple}\par
  \vspace{8pt}
  \begin{tcolorbox}[enhanced,colback=poporange,colframe=popink,boxrule=2.8pt,arc=13pt,
      drop shadow=popink,left=18pt,right=18pt,top=12pt,bottom=13pt,after skip=10pt]
    \poppill{#2 \textbullet\ #3}{popink}{popcream}\hspace{5pt}\poppill{#4}{white}{popink}\par\vspace{8pt}
    {\popdisplay\fontsize{14}{14}\selectfont\color{popink}LEÇON #5}\par\vspace{4pt}
    {\raggedright\hyphenpenalty=10000\exhyphenpenalty=10000\popdisplay\fontsize{25}{27}\selectfont\color{popcream}\MakeUppercase{#1}\par}
    \vspace{8pt}
    {\popxbold\fontsize{9.5}{9.5}\selectfont\color{popink}\MakeUppercase{Durée conseillée : #6}}
  \end{tcolorbox}
  \begin{tcolorbox}[enhanced,colback=white,colframe=popink,boxrule=2.2pt,arc=10pt,
      drop shadow=popink,left=16pt,right=16pt,top=10pt,bottom=10pt,after skip=8pt]
    \poppill{Dans cette leçon}{poppurple}{white}\par\vspace{5pt}
    \popsemi\fontsize{9.3}{12.6}\selectfont\color{popink}#7
  \end{tcolorbox}
  \noindent\begin{minipage}[t]{0.487\linewidth}
    \begin{tcolorbox}[enhanced,colback=popgreen,colframe=popink,boxrule=2.2pt,arc=10pt,
        drop shadow=popink,left=11pt,right=11pt,top=10pt,bottom=10pt,height=3.1cm,valign=center]
      \tcbox[colback=white,colframe=popink,boxrule=1.3pt,arc=5pt,boxsep=0pt,nobeforeafter,
        left=3pt,right=3pt,top=3pt,bottom=3pt,valign=center]{\qrcode[height=1.9cm]{https://play.google.com/store/apps/details?id=com.luvvix.onbuch&hl=fr}}%
      \hspace{8pt}\begin{minipage}[c]{0.5\linewidth}
        {\popdisplay\fontsize{11}{12.5}\selectfont\color{white}\MakeUppercase{Télécharge OnBuch}}\par\vspace{4pt}
        {\raggedright\popsemi\fontsize{8.4}{11}\selectfont\color{white}Gratuit \textbullet\ Google Play\\Tuteur IA, annales, concours, résultats\par}
      \end{minipage}
    \end{tcolorbox}
  \end{minipage}\hfill
  \begin{minipage}[t]{0.487\linewidth}
    \begin{tcolorbox}[enhanced,colback=poppurple,colframe=popink,boxrule=2.2pt,arc=10pt,
        drop shadow=popink,left=11pt,right=11pt,top=10pt,bottom=10pt,height=3.1cm,valign=center]
      \tikz[baseline=-0.7ex]\node[circle,fill=white,draw=popink,line width=1.3pt,minimum size=1.6cm,inner sep=0]{\popdisplay\fontsize{24}{24}\selectfont\color{poppurple}+};%
      \hspace{8pt}\begin{minipage}[c]{0.6\linewidth}
        {\popdisplay\fontsize{11}{12.5}\selectfont\color{white}\MakeUppercase{Passe à OnBuch+}}\par\vspace{4pt}
        {\raggedright\popsemi\fontsize{8.4}{11}\selectfont\color{white}Tous les cours signés, Tuteur IA illimité, fiches PDF — onglet OnBuch+ de l'app\par}
      \end{minipage}
    \end{tcolorbox}
  \end{minipage}
  \vspace{8pt}
  \popcontenu
  \vfill
  \signatureonbuch
  \restoregeometry
  \newpage
}

% Rangée « Ce cours contient » (page de garde)
\newcommand{\poptile}[3]{% #1 couleur #2 gros texte #3 légende
  \begin{tcolorbox}[enhanced,colback=#1,colframe=popink,boxrule=2pt,arc=9pt,shadow={2.5pt}{-2.5pt}{0pt}{popink},
      left=6pt,right=6pt,top=9pt,bottom=9pt,halign=center,valign=center,height=1.75cm,width=\linewidth,nobeforeafter]
    {\popdisplay\fontsize{12.5}{13}\selectfont\color{white}\MakeUppercase{#2}}\par\vspace{3pt}
    {\centering\popsemi\fontsize{7.8}{9.5}\selectfont\color{white}#3\par}
  \end{tcolorbox}}
\newcommand{\popcontenu}{%
  \par\noindent{\popxboldls\fontsize{7.5}{7.5}\selectfont\color{popmuted}CE COURS SIGNÉ CONTIENT}\par\vspace{7pt}
  \noindent\begin{minipage}[t]{0.235\linewidth}\poptile{popblue}{Cours}{complet, illustré, pas à pas}\end{minipage}\hfill
  \begin{minipage}[t]{0.235\linewidth}\poptile{poporange}{Méthodes}{et exercices résolus}\end{minipage}\hfill
  \begin{minipage}[t]{0.235\linewidth}\poptile{poppink}{Exercices}{type examen + corrigés}\end{minipage}\hfill
  \begin{minipage}[t]{0.235\linewidth}\poptile{popgreen}{Fiche bilan}{l'essentiel à retenir}\end{minipage}\par}

% Sommaire
\newtcolorbox{sommaire}{enhanced,colback=white,colframe=popink,boxrule=2.2pt,arc=10pt,
  drop shadow=popink,left=18pt,right=18pt,top=13pt,bottom=13pt,before skip=6pt,after skip=20pt,
  before upper={\poppill{Sommaire}{poppurple}{white}\par\vspace{9pt}\popsemi\fontsize{10.6}{14.5}\selectfont\color{popink}}}

% Signature de fin de cours — poussée tout en bas de la dernière page
\newcommand{\findecours}{%
  \par\vfill\nopagebreak
  \begin{center}\popsquiggle{poporange}\end{center}\vspace{4pt}
  \signatureonbuch
}
\AtBeginDocument{\def\llbracket{\mathopen{[\mkern-3.5mu[}}\def\rrbracket{\mathclose{]\mkern-3.5mu]}}} % intervalles d'entiers (glyphe ⟦ absent de la police)
% Glyphes absents de Fira Math : redéfinis à partir de glyphes disponibles
\AtBeginDocument{%
  \def\setminus{\mathbin{\backslash}}\let\smallsetminus\setminus
  \def\vdots{\mathord{\vbox{\baselineskip3.2pt\lineskiplimit0pt\kern4pt\hbox{.}\hbox{.}\hbox{.}}}}%
  \def\ddots{\mathinner{\mkern1mu\raise6.5pt\hbox{.}\mkern2mu\raise3.6pt\hbox{.}\mkern2mu\raise0.7pt\hbox{.}\mkern1mu}}%
  \def\triangle{\mathord{\scalebox{0.9}{$\symup{\Delta}$}}}%
  \def\nabla{\mathord{\raisebox{\depth}{\scalebox{1}[-1]{$\symup{\Delta}$}}}}%
  \def\checkmark{\popcheck{popdark}}%
  \def\star{\mathbin{\ast}}\def\bigstar{\mathord{\ast}}%
  \def\diamond{\mathbin{\scalebox{0.6}{$\Diamond$}}}%
  \def\bigcup{\mathop{\mathchoice{\raisebox{-0.3em}{\scalebox{1.7}{$\cup$}}}{\scalebox{1.25}{$\cup$}}{\cup}{\cup}}\displaylimits}%
  \def\bigcap{\mathop{\mathchoice{\raisebox{-0.3em}{\scalebox{1.7}{$\cap$}}}{\scalebox{1.25}{$\cap$}}{\cap}{\cap}}\displaylimits}%
  \def\lesseqqgtr{\mathrel{\lessgtr}}\def\bigoplus{\mathop{\oplus}\displaylimits}%
}
\providecommand{\ssi}{si et seulement si} % abréviation fréquente
\AtBeginDocument{\providecommand{\wideparen}{\overparen}} % arc de cercle

%% FIGFIX-BEGIN v4 — correctifs globaux des figures (docs/onbuch-generation/tools/figfix)
\usepackage{adjustbox}\usepackage{etoolbox}
% toute figure TikZ/pgfplots plus large que la ligne est réduite à la largeur disponible
\BeforeBeginEnvironment{tikzpicture}{\begin{adjustbox}{max width=\linewidth}}
\AfterEndEnvironment{tikzpicture}{\end{adjustbox}}
% légendes pgfplots lisibles par-dessus les courbes
\pgfplotsset{every axis/.append style={legend style={fill=white,fill opacity=0.92,text opacity=1,draw=black!15,inner sep=3pt}}}
% étiquettes d'arêtes (syntaxe edge["..."]) avec fond blanc
\tikzset{every edge quotes/.append style={fill=white,inner sep=1.2pt}}
%% FIGFIX-END
\begin{document}

\pagegarde{Contraction de texte : reformulation des idées d’un texte à résumer}{Français}{Tle Tech}{Français – classe de Terminale technique}{N09}{15 séances (fiche de progression harmonisée)}{Cette leçon outille l'élève de Terminale technique pour maîtriser la contraction de texte, épreuve phare du Baccalauréat : repérage de la structure argumentative, techniques de réduction et reformulation des idées essentielles. Elle s'appuie sur la lecture méthodique du Vieux nègre et la médaille de Ferdinand Oyono et mobilise des notions de stylistique (ironie, paradoxe), de sémantique (variétés lexicales du français) et de communication (énoncés constatifs et performatifs).
\vspace{4pt}
\begin{multicols}{2}\small
\begin{itemize}
\item À la fin de cette leçon, je sais identifier la structure argumentative d'un texte (thèse, arguments, exemples)
\item Je sais dégager les idées essentielles d'un texte et les hiérarchiser
\item Je sais appliquer les techniques de réduction (suppression, généralisation, condensation) pour respecter une limite de mots imposée
\item Je sais reformuler les idées d'un texte avec mes propres mots sans en trahir le sens
\item Je sais construire le plan d'un texte argumentatif et m'en servir pour rédiger un résumé
\item Je sais repérer et interpréter l'ironie et le paradoxe dans un texte littéraire
\end{itemize}\end{multicols}}

\begin{sommaire}
\begin{multicols}{2}
\begin{enumerate}
\item Lecture méthodique du Vieux nègre et la médaille (séances 1-2)
\item Stylistique : ironie et paradoxe
\item Sémantique et communication : variétés du français, énoncés constatifs et performatifs
\item La structure argumentative du texte
\item Techniques de réduction et reformulation des idées essentielles
\item Le plan du texte et la production du résumé
\item Activité d'intégration
\item Méthodes et exercices résolus
\item Exercices
\item Corrigés des exercices
\item Fiche bilan
\end{enumerate}
\end{multicols}
\end{sommaire}

\begin{objectifs}
\begin{itemize}
\item À la fin de cette leçon, je sais identifier la structure argumentative d'un texte (thèse, arguments, exemples)
\item Je sais dégager les idées essentielles d'un texte et les hiérarchiser
\item Je sais appliquer les techniques de réduction (suppression, généralisation, condensation) pour respecter une limite de mots imposée
\item Je sais reformuler les idées d'un texte avec mes propres mots sans en trahir le sens
\item Je sais construire le plan d'un texte argumentatif et m'en servir pour rédiger un résumé
\item Je sais repérer et interpréter l'ironie et le paradoxe dans un texte littéraire
\item Je sais distinguer les variétés lexicales du français (français standard, familier, régional) et leur effet dans un texte
\item Je sais distinguer un énoncé constatif d'un énoncé performatif dans un échange verbal
\item Je sais rédiger, justifier et améliorer une contraction de texte conforme aux exigences de l'examen
\end{itemize}
\end{objectifs}

\begin{prerequis}
\begin{itemize}
\item Notion de résumé vue en classe de Première (dégager l'essentiel d'un texte narratif)
\item Types de textes et registres (narratif, descriptif, argumentatif)
\item Connecteurs logiques et organisation du paragraphe
\item Lecture préalable du Vieux nègre et la médaille (parcours de l'œuvre en début d'année)
\item Notions de base sur l'énonciation (qui parle, à qui, dans quelle situation)
\end{itemize}
\end{prerequis}

\newpage

\coursec{Lecture méthodique du \emph{Vieux nègre et la médaille} (séances 1-2)}

\courssub{Situation d'accroche et problématique}

À la veille du Baccalauréat technique, un élève de Yaoundé reçoit un sujet d'épreuve : « Résumez ce texte de 600 mots en 150 mots ». Paniqué, il recopie des phrases entières, supprime quelques virgules, et dépasse largement la limite imposée. Le correcteur, lui, n'attendait qu'une chose : les \cle{idées essentielles}, reformulées avec les mots du candidat. Or, dans la vie professionnelle, cette compétence est quotidienne : compte-rendu de réunion dans une entreprise de Douala, synthèse d'un rapport ministériel à Yaoundé, note de service dans un atelier, résumé d'un dossier technique. Partout, il faut condenser sans trahir.

Mais avant de condenser, il faut \cle{comprendre}. On ne peut pas réduire un texte dont on n'a pas saisi la structure, le sens et les intentions. C'est pourquoi notre parcours commence par une lecture méthodique approfondie d'une œuvre majeure de la littérature camerounaise : \emph{Le Vieux nègre et la médaille} de Ferdinand Oyono. Ce roman nous servira de terrain d'entraînement tout au long de l'unité.

\textbf{Problématique :} comment lire méthodiquement un récit pour en dégager la structure et le sens, conditions indispensables à une contraction de texte réussie ?

\courssub{Présentation de l'œuvre et de son auteur}

\begin{definition}[La lecture méthodique]
La \cle{lecture méthodique} est une lecture guidée qui analyse un texte selon des étapes précises et ordonnées : la \textbf{situation d'énonciation} (qui parle, à qui, quand, où), la \textbf{structure} (schéma narratif, organisation du texte), le \textbf{thème} (ce dont parle le texte) et le \textbf{style} (comment l'auteur le dit : ton, figures, point de vue). Elle transforme le lecteur passif en analyste actif.
\end{definition}

\emph{Le Vieux nègre et la médaille} est un roman de \cle{Ferdinand Oyono}, écrivain camerounais, publié en \textbf{1956}, quatre ans avant l'indépendance du Cameroun. L'action se déroule sous la \cle{colonisation française}, dans un pays où l'administration coloniale, l'Église missionnaire et le système des chefs imposés bouleversent les sociétés traditionnelles.

Le roman raconte l'histoire de \cle{Meka}, un vieil homme du village de Doum, qui a tout donné à la cause coloniale : ses terres pour la construction de la mission, ses deux fils morts « pour la France » pendant la Seconde Guerre mondiale. Pour le récompenser, l'administration décide de lui remettre une \cle{médaille} de l'amitié franco-africaine, lors de la fête du 14 juillet. Mais cette médaille, symbole de reconnaissance, va se révéler un symbole d'illusion : la nuit qui suit la cérémonie, perdu sous la pluie, Meka est arrêté et brutalisé par les mêmes Blancs qui l'avaient honoré quelques heures plus tôt. Il comprend alors que la médaille ne change rien à sa condition d'homme colonisé.

\begin{savaistu}[Ferdinand Oyono, de l'écriture à l'État]
Ferdinand Oyono, né en 1929 à Ngoulemakong, dans la région du Sud du Cameroun, écrivit \emph{Une vie de boy} et \emph{Le Vieux nègre et la médaille} alors qu'il était encore étudiant en France. Ces deux romans, publiés en 1956, font de lui l'une des voix majeures de la littérature anticoloniale africaine. Après l'indépendance, il devint diplomate puis homme d'État camerounais, occupant notamment les fonctions d'ambassadeur et de ministre. Son œuvre littéraire, bien que courte, est étudiée dans le monde entier.
\end{savaistu}

\courssub{Lecture méthodique n° 5 : situation d'énonciation, personnages et espace}

\textbf{Situation d'énonciation.} Le récit est mené par un \cle{narrateur hétérodiégétique} (extérieur à l'histoire) : ce n'est pas Meka qui raconte sa propre histoire, mais une voix qui observe les personnages de l'extérieur. Le récit est à la troisième personne, au passé. Ce narrateur adopte une \cle{focalisation variable} : le plus souvent \textbf{zéro} (vision panoramique, « caméra »), il glisse fréquemment vers une \textbf{focalisation interne} sur Meka (accès à ses pensées, ses sensations : la douleur des souliers, l'angoisse de la nuit). Cette distance maîtrisée permet au narrateur de porter sur les événements un regard critique et ironique.

\textbf{Les personnages principaux.}

\begin{center}
\begin{tabularx}{\linewidth}{|l|X|}\hline
\rowcolor{popblueL} \textbf{Personnage} & \textbf{Rôle dans le récit} \\ \hline
Meka & Vieil homme de Doum, héros du roman ; il a sacrifié ses terres et ses deux fils ; il reçoit la médaille puis découvre l'illusion de la reconnaissance coloniale. \\ \hline
Le commandant & Chef du poste administratif de Mfoudi ; il organise la cérémonie ; il incarne l'administration coloniale, condescendante et méprisante. \\ \hline
Le père Vandermayer & Missionnaire installé à Doum ; il a bénéficié des terres de Meka ; il représente l'Église coloniale, complice du système. \\ \hline
Kelara et Engamba & Épouse et frère de Meka ; ils incarnent la communauté villageoise, témoin de la déchéance du vieil homme. \\ \hline
\end{tabularx}
\end{center}

\textbf{L'espace.} Deux lieux structurent le roman : \cle{Doum}, le village de Meka, espace africain traditionnel dominé par la mission, et \cle{Mfoudi}, le poste administratif où se déroule la cérémonie de remise de la médaille. Le trajet entre Doum et Mfoudi — le passage du monde villageois au monde colonial — est aussi un passage symbolique : c'est là que Meka perdra ses illusions.

\begin{popfigure}
\begin{center}
\begin{tikzpicture}[scale=0.95]
% Contour très simplifié du Cameroun (triangle pointe au nord)
\draw[thick, popdark, fill=popgreen!15] (0.5,0.2) -- (1.5,2.0) -- (3.5,4.5) -- (6.5,4.8) -- (7.5,3.5) -- (7.0,1.5) -- (5.5,0.5) -- (3.0,0.1) -- cycle;
% Fleuve (indicatif)
\draw[popblue, thick] (1.2,0.8) .. controls (2.5,1.8) and (4.0,3.0) .. (5.5,4.0);
% Villes
\fill[poporange] (3.8,2.8) circle (0.09);
\node[anchor=south west, popdark, font=\footnotesize] at (3.85,2.85) {\textbf{Doum (région du Centre, cadre fictif)}};
\fill[popink] (4.2,2.4) circle (0.09);
\node[anchor=north west, popdark, font=\footnotesize] at (4.25,2.35) {Yaoundé};
\fill[poppurple] (1.8,1.0) circle (0.09);
\node[anchor=north, popdark, font=\footnotesize] at (1.8,0.9) {Douala};
% Nord
\draw[->, thick, popdark] (7.2,0.8) -- (7.2,1.5) node[above]{\footnotesize N};
\node[popmuted, font=\footnotesize] at (6.0,0.3) {Océan Atlantique};
\node[popmuted, font=\footnotesize] at (4.5,4.6) {Nord-Cameroun};
\end{tikzpicture}
\end{center}
\legende{Carte très simplifiée du Cameroun : le récit se situe dans la région du Centre, entre le village de Doum et le poste administratif de Mfoudi (localisation indicative).}
\end{popfigure}

\courssub{Repérage du schéma narratif}

Tout récit obéit à une armature appelée \cle{schéma narratif} : \textbf{situation initiale}, \textbf{élément perturbateur}, \textbf{péripéties}, \textbf{dénouement}, \textbf{situation finale}. Repérer cette armature est la première étape de toute contraction : c'est elle qui fournira le squelette du résumé.

\begin{center}
\begin{tabularx}{\linewidth}{|l|X|}\hline
\rowcolor{popblueL} \textbf{Étape} & \textbf{Application au roman} \\ \hline
Situation initiale & Meka, vieil homme respecté de Doum, a tout sacrifié aux Blancs : ses terres à la mission, ses deux fils à la guerre. Il attend une reconnaissance. \\ \hline
Élément perturbateur & Le commandant annonce que Meka recevra une médaille le 14 juillet à Mfoudi. L'espoir naît ; le village se prépare. \\ \hline
Péripéties & Le voyage à Mfoudi, la cérémonie humiliante (Meka, en souliers neufs qui le blessent, exhibé comme un trophée), les beuveries, la nuit de pluie, l'errance, l'arrestation et la brutalisation par les Blancs. \\ \hline
Dénouement & De retour à Doum, Meka est un homme brisé : la médaille n'a rien changé. La désillusion est totale ; le regard sur les Blancs est transformé. \\ \hline
Situation finale & Meka, lucide, ne rit plus. La médaille est jetée ou oubliée ; le vieil homme a compris la réalité de la domination. \\ \hline
\end{tabularx}
\end{center}

\begin{popfigure}
\begin{center}
\begin{tikzpicture}[font=\footnotesize]
\draw[->, very thick, popdark] (0,0) -- (14,0);
\foreach \x/\titre/\desc in {
 0.8/{Situation initiale}/{Meka a tout donné aux Blancs},
 3.5/{Perturbateur}/{Annonce de la médaille},
 6.5/{Péripéties}/{Cérémonie, pluie, arrestation},
 10.0/{Dénouement}/{Retour brisé à Doum},
 12.8/{Situation finale}/{Désillusion, lucidité}}
{
 \fill[poporange] (\x,0) circle (0.12);
 \node[anchor=south, popdark, align=center] at (\x,0.3) {\textbf{\titre}};
 \node[anchor=north, popmuted, align=center, text width=2.6cm] at (\x,-0.3) {\desc};
}
\end{tikzpicture}
\end{center}
\legende{Frise chronologique du schéma narratif du \emph{Vieux nègre et la médaille} : les grandes étapes du parcours de Meka.}
\end{popfigure}

\courssub{Lecture méthodique n° 6 : les temps forts du récit}

Au-delà du squelette narratif, quatre \cle{temps forts} concentrent le sens du roman. Ce sont eux qu'il faudra absolument conserver dans un résumé, car ils portent l'évolution du personnage.

\begin{enumerate}
\item \textbf{La cérémonie de remise de la médaille} : Meka, chaussé de souliers neufs qui le blessent, est exhibé devant la foule. Le commandant lui épingle la médaille comme on décore un objet. La scène, apparemment glorieuse, est en réalité profondément humiliante : Meka est un instrument de propagande coloniale.
\item \textbf{La nuit de pluie} : après les réjouissances, Meka s'égare dans la nuit et sous l'orage. Cette errance physique traduit son égarement intérieur : l'homme qui croyait avoir trouvé sa place chez les Blancs est littéralement perdu.
\item \textbf{L'arrestation} : arrêté comme un vagabond par les forces coloniales, Meka est maltraité. La médaille sur sa poitrine ne le protège de rien. C'est le moment de bascule du roman.
\item \textbf{La désillusion} : de retour au village, Meka ne rit plus, ne croit plus. Sa lucidité nouvelle est une défaite, mais aussi une naissance : celle d'un regard décolonisé.
\end{enumerate}

\courssub{Point de vue, ton et figures de pensée : ironie et paradoxe}

Le narrateur adopte un \cle{point de vue externe} (focalisation zéro) mais se glisse fréquemment dans la conscience de Meka (focalisation interne). Surtout, il porte sur les événements un regard \cle{ironique} : il décrit la cérémonie avec les mots pompeux de l'administration, tout en laissant voir le ridicule et la cruauté de la scène. L'écart entre ce qui est dit (la « reconnaissance » de la France) et ce qui est montré (l'humiliation d'un vieil homme) crée l'ironie, arme majeure d'Oyono.

\begin{definition}[Ironie et Paradoxe]
L'\cle{ironie} est une figure de pensée qui consiste à dire le contraire de ce que l'on pense, ou à présenter une situation où il y a un décalage flagrant entre l'apparence et la réalité, dans un but critique ou moqueur.\\
Le \cle{paradoxe} est une affirmation en apparence contradictoire ou absurde, mais qui révèle une vérité profonde. Exemple dans le roman : « Meka est honoré par la médaille » (apparaît comme un honneur) mais « Meka est avili par la médaille » (réalité vécue). L'honneur est une humiliation.
\end{definition}

\begin{exemplebox}[Ironie et paradoxe dans la cérémonie]
\begin{itemize}
\item \textbf{Ironie narrative} : Le narrateur rapporte les discours officiels louant la « France généreuse » et l'« amitié franco-africaine » tandis que Meka souffre le martyre dans ses souliers neufs et que le commandant le traite en enfant.
\item \textbf{Paradoxe central} : La médaille, symbole de \emph{distinction} et d'\emph{égalité}, devient le marqueur le plus visible de la \emph{soumission} et de l'\emph{infériorisation}. « Il était décoré, donc il était libre » (pensée de Meka) $\rightarrow$ « Il était décoré, donc il était nu » (réalité de l'arrestation).
\end{itemize}
\end{exemplebox}

\begin{methode}[Dégager le thème d'un récit]
Pour dégager le thème d'un récit, procéder en trois étapes :
\begin{enumerate}
\item \textbf{Relever les récurrences} : quels mots, quelles images reviennent ? Ici : la médaille, les souliers qui blessent, la pluie, le mépris des Blancs, le rire de Meka (qui disparaît).
\item \textbf{Identifier les oppositions} : quels contraires structurent le texte ? Ici : honneur annoncé / humiliation réelle ; gratitude attendue / violence subie ; apparence de l'assimilation / réalité de la domination ; jour (cérémonie) / nuit (arrestation).
\item \textbf{Examiner la chute} : comment le récit se termine-t-il ? Ici : la médaille, censée couronner une vie de dévouement, se révèle vide. La chute révèle le sens : tout le roman dénonce l'illusion de l'assimilation.
\end{enumerate}
\end{methode}

Le \cle{thème central} est donc la \textbf{critique de l'aliénation coloniale} : Meka a intériorisé les valeurs du colonisateur au point de sacrifier les siennes, et la médaille est le symbole de cette illusion. Oyono dénonce l'\cle{assimilation} comme un leurre : le colonisé aura beau se conformer au modèle du colonisateur (porter le costume, parler français, donner ses fils), il ne sera jamais reconnu comme son égal.

\begin{attention}[Résumé ou analyse ?]
Erreur fréquente : confondre le \textbf{résumé du récit} (raconter l'histoire : « Meka reçoit une médaille, puis il est arrêté ») avec l'\textbf{analyse du sens} (ce que l'auteur dénonce : l'illusion de l'assimilation coloniale). Dans une contraction de texte, on ne raconte pas : on restitue les idées. La question à se poser n'est pas « que se passe-t-il ? » mais « que veut dire l'auteur à travers ce qui se passe ? ».
\end{attention}

\courssub{Relevé des passages clés pour la contraction}

Pour préparer les exercices de contraction des prochaines séances, retenons les passages qui condensent les idées essentielles du roman :

\begin{exemplebox}[Passages clés à retenir]
\begin{itemize}
\item \textbf{Le sacrifice de Meka} : il a donné ses terres à la mission et perdu ses deux fils à la guerre — idée : le dévouement total du colonisé.
\item \textbf{La cérémonie} : « On lui épingla la médaille » — Meka est un objet de propagande, non un homme honoré.
\item \textbf{Les souliers neufs} : ils le blessent pendant toute la journée — image concrète de l'assimilation qui fait mal.
\item \textbf{L'arrestation} : la médaille ne le protège pas des coups — idée : l'égalité proclamée est un mensonge.
\item \textbf{Le retour à Doum} : Meka ne rit plus — idée : la lucidité naît de la désillusion.
\end{itemize}
\end{exemplebox}

\begin{exoresolu}[Repérer les étapes du schéma narratif]
\textbf{Énoncé.} Voici un extrait résumé : « Le commandant fit savoir à Meka que le gouvernement avait décidé de lui décerner une médaille. Tout le village de Doum fut saisi d'effervescence. » À quelle étape du schéma narratif cet extrait correspond-il ? Justifiez.
\tcblower
\textbf{Solution.} Cet extrait correspond à l'\textbf{élément perturbateur}. En effet, avant cette annonce, la vie de Meka suivait son cours ordinaire (situation initiale : le vieil homme a tout donné aux Blancs et vit au village). L'annonce de la médaille vient rompre cet équilibre : elle déclenche l'action, suscite l'espoir de Meka et l'agitation du village, et lance le personnage dans une série d'événements nouveaux (le voyage à Mfoudi, la cérémonie, l'arrestation). C'est précisément la fonction d'un élément perturbateur : introduire un changement qui met le récit en mouvement.
\end{exoresolu}

\begin{exercice}[À faire en classe]
\begin{enumerate}
\item Relevez dans le roman deux passages où le narrateur utilise l'ironie pour décrire la cérémonie du 14 juillet. Expliquez l'écart entre ce qui est dit et ce qui est montré.
\item Identifiez un paradoxe dans le sort réservé à Meka le jour de la cérémonie. Formulez-le par une phrase du type : « C'est en étant honoré que Meka est le plus humilié. »
\item Rédigez en cinq phrases maximum le schéma narratif complet du roman, sans raconter les détails.
\item Pourquoi peut-on dire que la médaille est un « symbole d'illusion » ? Appuyez votre réponse sur le dénouement.
\end{enumerate}
\end{exercice}

\begin{corrige}[Exercice « À faire en classe »]
\begin{enumerate}
\item On peut relever la description pompeuse de la cérémonie (discours officiels, drapeaux, musique) alors que Meka souffre dans ses souliers neufs et est exhibé comme une bête curieuse : le narrateur épouse le vocabulaire officiel de la « reconnaissance » pour mieux laisser voir l'humiliation réelle. De même, le geste du commandant qui épingle la médaille est décrit avec solennité, alors qu'il s'agit d'un geste mécanique et condescendant.
\item Paradoxe : « C'est au moment où il reçoit la plus haute distinction que Meka est traité avec le plus de mépris » ou « La médaille qui devrait le rapprocher des Blancs scelle définitivement son exclusion. »
\item Meka, vieil homme de Doum, a tout sacrifié à la cause coloniale (situation initiale). L'annonce de la remise d'une médaille bouleverse sa vie (élément perturbateur). Le voyage à Mfoudi, la cérémonie humiliante, la nuit de pluie et l'arrestation brutale constituent les péripéties. Le retour au village, marqué par la désillusion totale, forme le dénouement. Meka est désormais lucide sur la réalité de la colonisation (situation finale).
\item La médaille est un symbole d'illusion parce qu'elle promet une reconnaissance et une égalité que le dénouement dément : la nuit même de la cérémonie, Meka est arrêté et maltraité par les Blancs. La distinction honorifique ne change rien à sa condition de colonisé ; elle récompense sa soumission, non sa dignité.
\end{enumerate}
\end{corrige}

\begin{aretenir}[L'essentiel de la séance]
\begin{itemize}
\item \emph{Le Vieux nègre et la médaille} (1956) de Ferdinand Oyono dénonce l'aliénation coloniale et l'illusion de l'assimilation.
\item La lecture méthodique suit quatre étapes : situation d'énonciation (narrateur hétérodiégétique, focalisation variable), structure (schéma narratif en 5 étapes), thème, style (ironie, paradoxe).
\item Le schéma narratif (situation initiale, élément perturbateur, péripéties, dénouement, situation finale) est le squelette du futur résumé.
\item Le ton ironique et les paradoxes permettent à Oyono de critiquer la colonisation sans discours explicite.
\item Comprendre le sens du texte est la condition première de toute contraction réussie : on ne résume pas une histoire, on restitue des idées.
\end{itemize}
\end{aretenir}

\coursec{Stylistique : ironie et paradoxe}

Dans la section précédente, notre lecture méthodique du \emph{Vieux nègre et la médaille} de Ferdinand Oyono nous a permis de suivre le parcours de Meka, ce vieillard qui reçoit une médaille en récompense de sa fidélité à la France, et qui découvre, au fil des événements, l'amertume cachée derrière cet « honneur ». Mais comprendre ce roman ne consiste pas seulement à suivre l'intrigue : il faut aussi saisir \emph{comment} Oyono dit les choses. Car l'auteur camerounais dit rarement ce qu'il pense de manière directe. Il feint, il exagère, il fait semblant d'admirer pour mieux condamner. C'est tout l'art de l'\cle{ironie} et du \cle{paradoxe}, deux figures de pensée que nous allons maintenant étudier en profondeur. Cette compétence sera décisive pour la suite de l'unité : on ne peut pas résumer correctement un texte ironique si l'on ne sait pas distinguer ce qui est dit de ce qui est pensé.

\courssub{L'ironie : dire le contraire pour mieux dire}

\textbf{L'intuition : une expérience de la vie quotidienne}

Commençons par une situation simple. Un camarade arrive en classe avec une heure de retard, essoufflé, les chaussures pleines de boue. Le professeur le regarde et déclare : « Bravo, quelle ponctualité exemplaire ! » Toute la classe rit. Pourquoi ? Parce que chacun comprend immédiatement que le professeur pense exactement le \emph{contraire} de ce qu'il dit. Le mot « bravo » et l'adjectif « exemplaire » sont des éloges, mais le contexte (le retard, la boue, le ton de la voix) les transforme en reproche. C'est cela, l'ironie : un écart entre les mots prononcés et la pensée réelle de celui qui parle, écart que le destinataire est invité à percevoir.

\begin{definition}[L'ironie]
L'\cle{ironie} est une figure de pensée qui consiste à dire le contraire de ce que l'on pense, ou à dire davantage ou moins que ce que l'on pense, avec une intention critique, railleuse ou comique. Celui qui emploie l'ironie fait semblant d'adhérer à ce qu'il dit, tout en donnant à son lecteur ou à son auditeur les moyens de comprendre qu'il s'en moque.
\end{definition}

\textbf{Les marques de l'ironie}

L'ironie n'est jamais signalée explicitement : c'est au lecteur de la débusquer. Elle laisse cependant des traces reconnaissables dans le texte.

\begin{itemize}
\item L'\cle{antiphrase} : c'est le procédé de base. On affirme le contraire de ce que l'on pense (« Quel courage ! » dit à un fuyard). L'éloge cache un blâme, le blâme peut cacher un éloge.
\item L'\cle{exagération} (ou hyperbole ironique) : on pousse l'éloge si loin qu'il en devient invraisemblable et donc suspect. Quand tout est « magnifique », « sublime », « inoubliable » à propos d'un objet dérisoire, le lecteur comprend que l'admiration est feinte.
\item Le \cle{ton faussement admiratif} ou naïf : le narrateur adopte le point de vue de celui qui croit, qui admire, qui s'émerveille, pour mieux laisser apparaître l'absurdité de la situation.
\item Le \cle{contraste avec le contexte} : les faits racontés démentent les paroles prononcées. C'est souvent l'indice le plus sûr.
\end{itemize}

\begin{methode}[Repérer l'ironie en quatre questions]
Face à un passage suspect, pose-toi successivement les questions suivantes :
\begin{enumerate}
\item \textbf{Que dit le texte ?} Releve les mots d'éloge ou de blâme, les superlatifs, les exclamations admiratives.
\item \textbf{Que montre le contexte ?} Qui souffre dans la scène ? Qui profite de la situation ? Les faits confirment-ils les paroles ?
\item \textbf{Y a-t-il un écart ?} Si l'éloge porte sur quelque chose de dérisoire, de douloureux ou d'injuste, l'écart est le signal de l'ironie.
\item \textbf{Quelle est l'intention de l'auteur ?} Demande-toi ce qu'il dénonce réellement : c'est le sens véritable du passage, celui qu'il faudra restituer dans un résumé.
\end{enumerate}
\end{methode}

\courssub{L'ironie dans \emph{Le Vieux nègre et la médaille}}

Oyono est un maître de l'ironie. Tout le roman est construit sur un décalage permanent entre le discours officiel colonial et la réalité vécue par les villageois.

\textbf{L'éloge exagéré de la médaille}

La médaille décernée à Meka est présentée par l'administration coloniale comme une récompense immense, la consécration d'une vie entière. Or que représente-t-elle objectivement ? Un petit bout de métal, remis à un vieillard qui a donné ses terres aux missionnaires, perdu ses deux fils à la guerre pour la France, et vécu toute sa vie dans la soumission. Plus le discours officiel enfle l'« honneur », plus le lecteur mesure la disproportion entre la récompense et le sacrifice. L'exagération elle-même devient accusatrice.

\begin{exemplebox}[« Quel honneur pour Meka ! »]
Lorsque le narrateur rapporte l'annonce de la décoration, il adopte le ton de l'émerveillement collectif : quel honneur pour Meka, quelle fierté pour le village, quelle reconnaissance de la France ! Mais le lecteur attentif remarque aussitôt l'écart : cet « honneur » récompense une vie de renoncements — la terre donnée, les fils morts, la soumission acceptée. Le narrateur feint l'admiration alors que la médaille est une récompense \emph{dérisoire} pour une vie de soumission. L'ironie transforme l'éloge en réquisitoire : plus on célèbre la médaille, plus on dénonce ce qu'elle prétend compenser.
\end{exemplebox}

\textbf{Le contraste entre discours colonial et réalité vécue}

Le discours des administrateurs et des missionnaires est plein de mots généreux : « amitié », « civilisation », « gratitude », « fraternité ». Mais la réalité que vit Meka après la cérémonie — l'humiliation, l'arrestation, la nuit passée sous la pluie, le mépris des mêmes qui l'applaudissaient — dément chacun de ces mots. L'ironie d'Oyono consiste à laisser les deux registres se télescoper : les belles paroles d'un côté, les faits brutaux de l'autre. Le lecteur n'a pas besoin de commentaire explicite : le contraste suffit à dénoncer l'\cle{hypocrisie coloniale}.

\begin{popfigure}
\begin{center}
\begin{tabular}{|p{0.30\linewidth}|p{0.30\linewidth}|p{0.30\linewidth}|}
\hline
\rowcolor{popblueL} \textbf{Énoncé ironique} & \textbf{Sens littéral (premier degré)} & \textbf{Sens réel (intention de l'auteur)} \\
\hline
« Quel honneur pour Meka de recevoir cette médaille ! » & La médaille est une récompense magnifique qui comble Meka de fierté. & La médaille est dérisoire face aux sacrifices de Meka (terre, fils, soumission) ; l'éloge dénonce l'injustice. \\
\hline
« La France sait reconnaître ses fidèles serviteurs. » & La France est généreuse et reconnaissante envers ceux qui l'ont servie. & La France use des hommes puis les jette ; la « reconnaissance » est un mot vide qui masque l'exploitation. \\
\hline
« Les Blancs sont les amis des vieillards comme Meka. » & Les colons éprouvent une véritable amitié pour les villageois. & Cette amitié est de façade : dès que Meka dérange, il est humilié et maltraité par ceux-là mêmes qui le célébraient. \\
\hline
\end{tabular}
\end{center}
\legende{Tableau comparatif de trois énoncés ironiques du roman : pour chacun, le sens littéral (ce que les mots disent) s'oppose au sens réel (ce que l'auteur pense). C'est cet écart qu'il faut restituer dans un résumé.}
\end{popfigure}

\courssub{Le paradoxe : la vérité qui choque}

\textbf{L'intuition}

Revenons à notre expérience de classe. Supposons qu'un élève déclare : « Ce retard m'a rendu service. » Le sens commun proteste : un retard est une faute, il ne peut pas être un service. Pourtant, si le retard a permis à l'élève d'éviter un accident sur la route, l'affirmation, d'abord choquante, révèle une vérité. C'est le mécanisme du paradoxe : une affirmation contraire à l'opinion commune, qui déroute d'abord, puis oblige à réfléchir, et finit par révéler une vérité profonde que l'opinion commune cachait.

\begin{definition}[Le paradoxe]
Le \cle{paradoxe} est une affirmation contraire à l'opinion commune (au \emph{sens commun}), qui choque par son apparente absurdité, mais qui, à la réflexion, exprime une vérité profonde ou inattendue. Il ne faut pas le confondre avec la contradiction simple : le paradoxe n'est pas une erreur de raisonnement, c'est une provocation à penser.
\end{definition}

\begin{attention}[Ironie et paradoxe : ne pas confondre]
Les deux figures sont proches mais distinctes. L'\cle{ironie} dit le \emph{contraire} de ce qu'on pense (le sens réel est l'inverse du sens littéral). Le \cle{paradoxe} dit quelque chose qui \emph{semble} faux ou absurde mais qui est vrai (le sens réel est le sens littéral, à condition de le comprendre en profondeur). Dans l'ironie, on renverse l'énoncé ; dans le paradoxe, on le creuse.
\end{attention}

\textbf{Le paradoxe central du roman : l'honneur qui humilie}

Toute l'architecture du \emph{Vieux nègre et la médaille} repose sur un paradoxe : la médaille, censée être le plus grand honneur de la vie de Meka, devient la source de son humiliation la plus profonde. Le jour de gloire se transforme en jour de honte ; la récompense révèle la servitude au lieu de l'effacer. En recevant la médaille, Meka ne s'élève pas : il découvre qu'il n'a jamais été considéré comme un homme, mais comme un objet décoratif de la propagande coloniale. L'« honneur » est donc, en vérité, une insulte déguisée.

Ce paradoxe produit un double effet de sens. D'une part, il dénonce l'hypocrisie du système colonial, qui pare de mots nobles une relation de domination. D'autre part, il installe une \cle{distance critique} entre le narrateur et les événements : le narrateur ne prend jamais parti ouvertement, il laisse le paradoxe agir comme une évidence qui s'impose d'elle-même au lecteur. C'est une stratégie plus efficace qu'un discours militant direct, car c'est le lecteur lui-même qui conclut.

\begin{aretenir}[L'essentiel]
\begin{itemize}
\item L'\cle{ironie} : dire le contraire de ce que l'on pense, pour railler (antiphrase, exagération, ton faussement admiratif).
\item Le \cle{paradoxe} : affirmer ce qui choque le sens commun mais révèle une vérité.
\item Dans le roman, l'ironie porte sur l'éloge exagéré de la médaille ; le paradoxe structure toute l'œuvre : l'honneur devient humiliation.
\item Effet de sens commun : dénonciation de l'hypocrisie coloniale et distance critique du narrateur.
\end{itemize}
\end{aretenir}

\courssub{Le lien avec la contraction de texte}

Pourquoi étudier ces figures dans une unité consacrée au résumé ? Parce que la contraction de texte exige de \emph{reformuler} les idées de l'auteur avec ses propres mots. Or, si l'on reformule un énoncé ironique au premier degré, on dit exactement l'inverse de ce que l'auteur pensait : on commet un contresens.

\begin{attention}[Le piège majeur du résumé]
Dans un résumé, reformuler une ironie au premier degré constitue un \cle{contresens} majeur, plus grave encore qu'une omission. Écrire « Oyono montre que la médaille est un grand honneur pour Meka », c'est trahir le texte : Oyono montre précisément le contraire. La règle d'or est donc la suivante : avant de reformuler, toujours se demander si le passage est au premier ou au second degré. En cas d'ironie, on restitue la \emph{pensée réelle} de l'auteur, éventuellement en signalant le procédé : « l'auteur dénonce ironiquement... », « sous des éloges exagérés, le narrateur critique... ».
\end{attention}

\begin{exoresolu}[Reformuler sans trahir]
\textbf{Énoncé.} Voici un passage reformulé par un élève dans le cadre d'un résumé : « Meka reçoit une médaille, ce qui constitue pour lui et pour son village un honneur exceptionnel dont tous sont fiers. » Explique pourquoi cette reformulation est fautive, puis propose une reformulation correcte.
\tcblower
\textbf{Solution détaillée.} La reformulation de l'élève prend l'épisode au premier degré : elle retient l'éloge (honneur exceptionnel, fierté) sans percevoir l'ironie qui le traverse. Or le contexte du roman — les sacrifices de Meka, le caractère dérisoire de la récompense, les humiliations qui suivent la cérémonie — montre que l'auteur pense le contraire de ce que les mots d'éloge affirment. L'élève commet donc un contresens : il attribue à Oyono une admiration qu'Oyono moque. Une reformulation correcte doit restituer la pensée réelle de l'auteur et, si possible, signaler le procédé. Par exemple : « En présentant avec une admiration exagérée la médaille décernée à Meka, Oyono dénonce ironiquement la dérision d'une récompense qui prétend payer une vie entière de soumission et de sacrifices. » Cette version respecte le sens profond du texte : c'est elle qui pourra figurer dans le résumé.
\end{exoresolu}

\begin{exercice}[À toi de jouer]
Releve dans ta lecture du roman deux passages où le narrateur semble louer les colons ou la médaille. Pour chacun : (1) identifie la marque d'ironie (antiphrase, exagération, ton admiratif) ; (2) formule le sens réel de l'auteur ; (3) rédige une phrase de résumé qui restitue ce sens sans contresens. Tu pourras t'appuyer sur la méthode en quatre questions étudiée plus haut.
\end{exercice}

\begin{corrige}[Exercice : indications]
Les passages les plus riches sont ceux de l'annonce de la décoration et du discours officiel lors de la cérémonie. Pour l'annonce : la marque est l'exagération collective (tout le village s'émerveille d'un petit ruban) ; le sens réel est la dénonciation de la disproportion entre le sacrifice et la récompense ; phrase de résumé possible : « L'auteur souligne, par l'enthousiasme excessif qu'il prête aux villageois, le contraste entre l'immensité des sacrifices de Meka et l'insignifiance de la récompense. » Pour le discours officiel : la marque est le contraste entre les mots généreux (« amitié », « reconnaissance ») et les faits (mépris, humiliation) ; le sens réel est la dénonciation de l'hypocrisie coloniale ; phrase de résumé possible : « Sous les belles paroles du discours officiel, le narrateur laisse apparaître le mépris réel des colons pour ceux qu'ils prétendent honorer. »
\end{corrige}

\begin{savaistu}[L'ironie, une arme camerounaise]
Ferdinand Oyono, né en 1929 à Ngoulemakong, près d'Ebolowa, dans la région du Sud, a publié \emph{Le Vieux nègre et la médaille} en 1956, quelques années avant l'indépendance du Cameroun. À une époque où la critique frontale du colonialisme était dangereuse, l'ironie offrait une arme précieuse : elle permettait de dénoncer sans jamais avoir l'air de le faire, en laissant le lecteur tirer lui-même les conclusions. Cette tradition de la satire ironique reste vivante dans la littérature camerounaise et africaine francophone, et Oyono est devenu, après sa carrière littéraire, un diplomate et un homme d'État de premier plan : ambassadeur, puis ministre dans les gouvernements camerounais.
\end{savaistu}

En résumé, l'ironie et le paradoxe ne sont pas des ornements du style d'Oyono : ils sont le cœur de son propos. Savoir les repérer, c'est accéder au sens véritable du roman ; savoir les restituer sans les aplatir, c'est la condition d'un résumé fidèle. Nous allons maintenant poursuivre notre étude de la langue en examinant, dans la section suivante, les variétés du français et la distinction entre énoncés constatifs et performatifs, avant d'aborder la structure argumentative du texte.

\coursec{Sémantique et communication : variétés du français, énoncés constatifs et performatifs}

Dans la section précédente, nous avons vu comment Ferdinand Oyono joue avec les figures de pensée — l'ironie et le paradoxe — pour dénoncer l'absurdité de la colonisation. Mais l'ironie ne travaille pas seulement sur les idées : elle travaille aussi sur les \emph{mots}. Pourquoi le commandant parle-t-il un français si solennel, si « précieux », alors que Meka s'exprime dans une langue simple, proche de la terre et du village ? Pourquoi certaines phrases, prononcées au bon moment par la bonne personne, ont-elles le pouvoir de \emph{changer la réalité} ? C'est à ces deux questions que répond la présente section : d'abord en étudiant les \cle{variétés lexicales du français} et leurs effets de sens dans le roman, ensuite en découvrant la distinction fondamentale entre \cle{énoncés constatifs} et \cle{énoncés performatifs}.

\courssub{Les variétés lexicales du français}

\courssub{Intuition : une même langue, plusieurs façons de la parler}

Imaginez la scène suivante. Le matin, un élève dit à son camarade : « Eh, t'as piqué mon stylo, rends-moi ça vite ! » Le soir, devant le proviseur, le même élève déclare : « Monsieur le Proviseur, je vous prie de bien vouloir m'excuser de ce retard. » Le message de fond est proche (il s'agit chaque fois d'objets, de fautes, d'échanges), mais les \emph{mots choisis} ne sont pas les mêmes. Cette observation simple montre que le français n'est pas un bloc uniforme : c'est une langue à \emph{plusieurs couches}, que le locuteur choisit selon la situation, l'interlocuteur et l'effet recherché.

\begin{definition}[Les registres de langue]
On appelle \cle{registre de langue} (ou niveau de langue) l'ensemble des choix lexicaux et grammaticaux adaptés à une situation de communication. On distingue principalement :
\begin{itemize}
\item le \cle{registre soutenu} : langue soignée, précise, recherchée (discours officiels, écrits administratifs, littérature classique) ;
\item le \cle{registre courant} ou \cle{français standard} : langue de la communication ordinaire, correcte et neutre (école, presse, conversation polie) ;
\item le \cle{registre familier} : langue décontractée de la vie quotidienne (« boulot » pour travail, « bagnole » pour voiture) ;
\item le \cle{registre populaire} : langue des milieux modestes, souvent marquée par des tournures non conformes à la norme (« y a pas » pour « il n'y a pas »).
\end{itemize}
À ces registres s'ajoutent les \cle{variétés régionales} : le français parlé au Cameroun, en Côte d'Ivoire ou au Québec possède des mots et des tournures propres à chaque pays.
\end{definition}

\begin{aretenir}[Ce qu'il faut retenir]
Le choix d'un mot n'est jamais neutre : il signale l'origine sociale du locuteur, son degré d'éducation, la distance qu'il met entre lui et son interlocuteur, et l'effet qu'il veut produire. En analyse de texte, relever le registre lexical, c'est déjà interpréter.
\end{aretenir}

\courssub{Le français du Cameroun : une variété régionale vivante}

Le français parlé au Cameroun s'est enrichi, au contact des langues nationales et de la réalité locale, de mots et d'emplois spécifiques. On parle de \cle{français camerounais} : ce n'est pas un « mauvais français », c'est une variété légitime, reconnue par les linguistes, qui exprime des réalités que le français de France ne nomme pas.

\begin{center}
\begin{tabularx}{\linewidth}{|l|X|X|}
\hline
\rowcolor{popblueL} \textbf{Registre / variété} & \textbf{Exemple lexical} & \textbf{Équivalent standard / effet} \\
\hline
Soutenu (discours officiel) & « Votre bravoure honore la France » & Langue de cérémonie, solennelle et abstraite \\
\hline
Courant / standard & « Le vieux Meka reçoit une médaille » & Langue neutre du récit \\
\hline
Familier & « bouquin », « flic », « bouffer » & Proximité, relâchement \\
\hline
Populaire & « y a », « c'est-y pas vrai ? » & Marque sociale des milieux modestes \\
\hline
Français du Cameroun & « ndolé », « makossa », « le paysan » (le villageois), « être au village » (vivre en brousse) & Réalités locales, authenticité africaine \\
\hline
\end{tabularx}
\end{center}

\begin{savaistu}[Le français du Cameroun, une richesse]
Le français du Cameroun possède des mots propres qui enrichissent la langue : \emph{ndolé} (plat de feuilles et d'arachides), \emph{makossa} (musique et danse de Douala), \emph{le paysan} employé au sens de « villageois, homme de la brousse », ou encore des expressions comme « être au village » pour « vivre en brousse ». Ces mots entrent peu à peu dans les dictionnaires : une langue vivante est une langue qui emprunte et qui s'adapte. Les écrivains africains, comme Oyono, en profitent pour donner à leurs textes une couleur locale immédiatement reconnaissable.
\end{savaistu}

\courssub{Repérage dans \emph{Le Vieux Nègre et la médaille}}

Oyono organise dans son roman un véritable \emph{contraste linguistique} entre deux mondes.

\textbf{1. Le registre soutenu des discours officiels.} Le commandant, le père Vandermayer et les autorités coloniales emploient une langue cérémonieuse, pleine d'abstractions : « la France éternelle », « la reconnaissance de la patrie », « vos mérites éminents ». Ce vocabulaire grandiloquent crée une \cle{distance sociale} : celui qui parle ainsi se place \emph{au-dessus} de son auditoire. Mais chez Oyono, cette langue pompeuse devient aussi un piège : plus le discours est solennel, plus le ridicule de la situation (décorer un vieillard épuisé pour le récompenser d'avoir tout perdu) éclate au grand jour. Le soutenu devient un instrument de \cle{dérision}.

\textbf{2. Le langage des villageois.} Meka, Kelara et les gens de Doum s'expriment dans un français simple, parfois approximatif, mêlé de réalités villageoises : on parle de la terre, des champs, des ancêtres, de la « brousse ». Ce langage direct produit un effet d'\cle{authenticité} : le lecteur entend la voix réelle du village africain, avec sa sagesse proverbiale et son bon sens.

\textbf{3. Les traces du français camerounais.} Le roman porte la marque du lexique local : noms de plats, de danses, emplois particuliers de mots comme « paysan » ou « frère » (au sens d'« homme du même village »). Ces traces ancrent le récit dans le sol camerounais.

\begin{exemplebox}[Effet de sens d'un choix lexical]
Lorsque le commandant déclare que la médaille est « le témoignage éclatant de la reconnaissance de la France », le mot \emph{éclatant} et la périphrase \emph{témoignage de la reconnaissance} relèvent du soutenu. Effet produit : le discours sonne creux, car cette « reconnaissance » ne se traduit par rien de concret — ni terre rendue, ni fils ressuscités. Le contraste entre la noblesse des mots et la pauvreté de la réalité nourrit l'ironie étudiée dans la section précédente.
\end{exemplebox}

\courssub{Énoncés constatifs et énoncés performatifs}

\courssub{Intuition : parler, est-ce seulement décrire ?}

Considérons deux phrases prononcées le jour de la fête à Doum :
\begin{itemize}
\item « La cérémonie a lieu à Doum. »
\item « Je déclare la cérémonie ouverte. »
\end{itemize}
La première phrase \emph{décrit} un fait : on peut vérifier si elle est vraie ou fausse (la cérémonie a-t-elle bien lieu à Doum ?). La seconde ne décrit rien : au moment où le commandant la prononce, la cérémonie \emph{devient} ouverte. Les mots, ici, ne constatent pas : ils \emph{agissent}. Telle est l'intuition fondatrice de la distinction entre constatifs et performatifs.

\begin{definition}[L'énoncé constatif]
Un \cle{énoncé constatif} décrit une réalité, rapporte un fait, formule une information. Il est soumis à la question de la \cle{vérité} : il est soit vrai, soit faux. Exemples : « Meka a donné ses terres à la mission », « Il pleut sur Doum », « La médaille est en argent ».
\end{definition}

\begin{definition}[L'énoncé performatif]
Un \cle{énoncé performatif} accomplit, par le seul fait d'être prononcé dans de bonnes conditions, l'action qu'il énonce. Dire, c'est faire. Les verbes performatifs typiques sont : \emph{promettre, déclarer, baptiser, condamner, ordonner, jurer, décorer, pardonner, nommer}. Exemples : « Je te promets de revenir », « Je déclare la séance ouverte », « Je vous décore ».
\end{definition}

\begin{attention}[Un performatif n'est ni vrai ni faux]
On ne peut pas répondre « c'est faux ! » à quelqu'un qui dit « je te promets de venir » : la promesse est faite dès qu'elle est prononcée. Un performatif n'est donc pas vrai ou faux : il \cle{réussit} ou il \cle{échoue}. Il échoue si les conditions ne sont pas remplies : un élève qui crie dans la cour « Je déclare le lycée fermé ! » n'accomplit rien du tout, car il n'a pas l'autorité pour le faire. Le performatif exige une \emph{situation rituelle} et un \emph{locuteur habilité}.
\end{attention}

\courssub{Les conditions de réussite du performatif}

Pour qu'un performatif « prenne », c'est-à-dire transforme réellement la situation, plusieurs conditions doivent être réunies :

\begin{enumerate}
\item \textbf{L'autorité du locuteur} : celui qui parle doit être habilité à accomplir l'acte (un juge peut condamner, un commandant peut décorer, un élève ne peut pas fermer le lycée).
\item \textbf{La situation rituelle} : l'énonciation doit se dérouler dans un cadre reconnu (cérémonie officielle, tribunal, mariage), souvent avec des formules consacrées (« Au nom de... »).
\item \textbf{L'acceptation par l'auditoire} : les témoins doivent reconnaître la validité de l'acte.
\item \textbf{La sincérité minimale} : pour les actes comme promettre, le locuteur est supposé avoir l'intention de tenir sa parole.
\end{enumerate}

\begin{methode}[Identifier un énoncé performatif]
Pour savoir si un énoncé est performatif, procédez en trois étapes :
\begin{enumerate}
\item Repérez un verbe à la première personne du présent (« je promets », « je déclare », « j'ordonne »).
\item Insérez mentalement « \cle{par ces mots, je...} » devant le verbe : si la phrase reste acceptable (« par ces mots, je te promets de venir »), c'est un performatif.
\item Vérifiez que l'énoncé ne peut pas être qualifié de vrai ou de faux, mais seulement de réussi ou de raté.
\end{enumerate}
Testez avec « je cours » : « par ces mots, je cours » est absurde — « courir » n'est pas un verbe performatif, « je cours » est un constatif.
\end{methode}

\begin{exemplebox}[La remise de la médaille : un performatif en action]
Quand le commandant dit à Meka : « Au nom de la France, je vous remets cette médaille », il ne décrit pas une remise de médaille : il \emph{accomplit réellement} la remise de la médaille. Avant la phrase, Meka n'était pas décoré ; après la phrase, il l'est. Toutes les conditions de réussite sont réunies : le commandant a l'autorité (il représente la puissance coloniale), la situation est rituelle (cérémonie publique du 14 juillet, en présence des notables), la formule consacrée est employée (« Au nom de la France »). Le performatif est donc l'instrument même du \cle{pouvoir colonial} : par la parole, l'administrateur transforme le statut des hommes.
\end{exemplebox}

\begin{popfigure}
\begin{center}
\begin{tikzpicture}[
  grow=down,
  level distance=1.5cm,
  level 1/.style={sibling distance=7.2cm},
  level 2/.style={sibling distance=3.6cm},
  every node/.style={draw, rounded corners, align=center, font=\small, fill=popcream},
  edge from parent/.style={draw=popmuted, -{Stealth[length=2mm]}}
]
\node[fill=popblueL, font=\small\bfseries] {Énoncé}
  child { node[fill=popgreenL, align=center]{\textbf{Constatif}\\ décrit une réalité}
    child { node[align=center]{\textbf{Vrai}\\ « La cérémonie\\ a lieu à Doum »} }
    child { node[align=center]{\textbf{Faux}\\ « Meka est\\ un jeune homme »} }
  }
  child { node[fill=poporangeL, align=center]{\textbf{Performatif}\\ accomplit un acte}
    child { node[align=center]{\textbf{Réussi}\\ « Je vous décore »\\ (commandant, cérémonie)} }
    child { node[align=center]{\textbf{Raté}\\ « Je déclare Doum\\ capitale » (Meka)} }
  };
\end{tikzpicture}
\end{center}
\legende{Arbre de classification des énoncés. Un constatif se juge selon la vérité (vrai / faux) ; un performatif se juge selon la réussite (réussi / raté), qui dépend de l'autorité du locuteur et de la situation rituelle. Exemples tirés de l'univers du roman.}
\end{popfigure}

\courssub{Application : le discours du commandant, ou le pouvoir par les mots}

Analysons maintenant le discours de remise de la médaille comme un entrelacement de constatifs et de performatifs.

\textbf{Les constatifs du discours} (« Vous avez donné vos terres à la mission », « Vos fils sont morts pour la France ») présentent les sacrifices de Meka comme des faits. Ils sont vrais — et c'est précisément leur vérité qui rend la suite insoutenable : le lecteur mesure l'écart entre l'immensité des pertes et la maigreur de la récompense.

\textbf{Les performatifs du discours} (« Au nom de la France, je vous remets cette médaille », « je vous félicite ») accomplissent des actes officiels. Ils réussissent parfaitement sur le plan formel : autorité, rituel, formules consacrées, tout y est. Mais Oyono révèle un paradoxe (nous retrouvons la figure de la section précédente) : le performatif réussit \emph{matériellement} tout en échouant \emph{moralement}. La médaille est bien remise, mais elle ne répare rien ; l'acte de parole colonial, si puissant soit-il, ne peut pas ressusciter les fils ni rendre les terres. Le performatif apparaît ainsi comme le symbole d'un pouvoir qui croit que les mots suffisent à effacer les injustices.

\begin{exoresolu}[Constatif ou performatif ?]
\textbf{Énoncé.} Classez les énoncés suivants en constatifs ou performatifs, en justifiant : (a) « La médaille de Meka brille au soleil » ; (b) « Je déclare la fête terminée », dit le commandant ; (c) « Je vous promets que la France n'oubliera jamais vos services » ; (d) « Meka est rentré chez lui sous la pluie ».
\tcblower
\textbf{Solution détaillée.}
\begin{itemize}
\item (a) \textbf{Constatif} : il décrit une réalité observable ; on peut vérifier s'il est vrai ou faux (la médaille brille-t-elle ?).
\item (b) \textbf{Performatif} : le verbe « déclarer » est à la première personne du présent ; le test « par ces mots, je déclare la fête terminée » fonctionne ; le commandant a l'autorité et la situation est rituelle : l'énoncé réussit, la fête est réellement terminée.
\item (c) \textbf{Performatif} : « promettre » est un verbe performatif (« par ces mots, je vous promets... »). La promesse est faite dès qu'elle est prononcée ; on ne peut pas la traiter de fausse, seulement la soupçonner d'être insincère — ce que l'ironie d'Oyono suggère fortement.
\item (d) \textbf{Constatif} : il rapporte un fait du récit, soumis à la vérité. Le test échoue : « par ces mots, Meka rentre chez lui » n'a pas de sens.
\end{itemize}
\end{exoresolu}

\begin{exercice}[À vous de jouer]
\begin{enumerate}
\item Relevez dans le discours du commandant deux énoncés constatifs et deux énoncés performatifs. Justifiez chaque classement à l'aide du test « par ces mots, je... ».
\item Un villageois de Doum s'écrie : « Je déclare que la médaille ne vaut rien ! » Ce performatif réussit-il ? Analysez les conditions d'énonciation.
\item Identifiez dans le roman trois mots ou expressions relevant du français du Cameroun ou du langage des villageois, et expliquez l'effet de sens produit (authenticité, distance, dérision).
\end{enumerate}
\end{exercice}

\begin{corrige}[Exercice : éléments de correction]
\begin{enumerate}
\item Exemples attendus : constatifs — « Vous avez donné vos terres », « Vos fils sont morts pour la France » (vérifiables, vrais ou faux) ; performatifs — « Je vous remets cette médaille », « Je vous félicite » (le test « par ces mots » fonctionne, l'acte s'accomplit en le disant).
\item Le performatif échoue : le villageois n'a pas l'autorité requise, il n'existe pas de rituel reconnu pour « dévaluer » une médaille par la parole. Son énoncé reste une opinion déguisée en acte ; en réalité, « la médaille ne vaut rien » fonctionne comme un constatif (un jugement de valeur), précédé d'une formule sans efficacité.
\item Exemples : « paysan » au sens de villageois (ancrage local), les noms de mets et de danses (authenticité), les proverbes de Kelara (sagesse orale africaine). Effet : le lecteur entend la voix du village, ce qui rend d'autant plus criant le contraste avec la langue officielle du commandant.
\end{enumerate}
\end{corrige}

\begin{aretenir}[Bilan de la section]
\begin{itemize}
\item Le français possède plusieurs \cle{registres} (soutenu, courant, familier, populaire) et des \cle{variétés régionales} comme le français du Cameroun ; chaque choix lexical produit un effet de sens (distance sociale, dérision, authenticité).
\item Un \cle{constatif} décrit le monde : il est vrai ou faux. Un \cle{performatif} transforme le monde : il réussit ou échoue selon l'autorité du locuteur et la situation rituelle.
\item Dans \emph{Le Vieux Nègre et la médaille}, le discours du commandant montre que le performatif est l'outil du pouvoir colonial : il accomplit des actes réels (décorer, nommer) tout en masquant l'injustice — ce que l'ironie d'Oyono dénonce.
\end{itemize}
Ces compétences d'analyse du lexique et des énoncés nous seront précieuses dans la suite de la leçon : pour résumer un texte argumentatif, il faudra en effet distinguer ce que le texte \emph{décrit} de ce qu'il \emph{fait} à son lecteur, et reformuler les idées essentielles dans un français standard, débarrassé des effets de registre.
\end{aretenir}

\coursec{La structure argumentative du texte}

Après avoir étudié les variétés du français et la distinction entre énoncés constatifs et performatifs, nous disposons d'outils pour analyser \emph{comment} un texte construit son sens et son efficacité. Un texte argumentatif ne se contente pas de constater ou d'accomplir des actes de langage : il organise une démonstration pour faire adhérer le lecteur à une thèse. C'est cette architecture interne — le squelette logique du texte — que nous allons disséquer, car la maîtriser est la condition \emph{sine qua non} d'une contraction réussie : on ne peut résumer que ce qu'on a compris dans sa structure profonde.

\courssub{Qu'est-ce qu'un texte argumentatif ?}

\begin{definition}[Texte argumentatif]
Un \textbf{texte argumentatif} est un texte dont la finalité principale est de \textbf{convaincre} (par la raison, le logos) ou de \textbf{persuader} (par l'émotion, le pathos, et l'éthos) un destinataire d'adhérer à une \textbf{thèse} (une opinion, une valeur, une décision). Il ne se limite pas à informer (fonction informative) ni à donner des ordres (fonction injonctive) : il \emph{justifie} un point de vue.
\end{definition}

\begin{propriete}[Architecture canonique : Thèse – Arguments – Exemples]
Tout texte argumentatif bien construit obéit à une structure fondamentale en trois étages, souvent complétée par une conclusion :
\begin{enumerate}
    \item \textbf{La Thèse} (ou proposition centrale) : l'idée force, le jugement que l'auteur défend. Elle est souvent explicite (affirmée d'emblée ou en conclusion) ou implicite (à reconstruire).
    \item \textbf{Les Arguments} : les raisons logiques, les preuves ou les motifs qui soutiennent la thèse. Ils répondent à la question « Pourquoi ? ».
    \item \textbf{Les Exemples (ou illustrations)} : des faits concrets, des anecdotes, des citations, des statistiques, des références historiques qui \emph{incarnent} les arguments. Ils répondent à la question « Par exemple ? ».
    \item \textbf{La Conclusion} : réaffirmation de la thèse, ouverture, appel à l'action.
\end{enumerate}
\end{propriete}

\begin{popfigure}
\begin{tikzpicture}[scale=0.9, every node/.style={font=\small}]
    % Couleurs
    \colorlet{thesiscol}{popred!80!popdark}
    \colorlet{argcol}{popblue!80!popdark}
    \colorlet{excol}{popgreen!70!popdark}
    \colorlet{conclcol}{poppurple!80!popdark}

    % Pyramide inversée (base large = exemples)
    % Niveau 1 : Thèse (sommet)
    \node[draw=thesiscol, thick, fill=thesiscol!15, rounded corners=6pt, minimum width=8cm, minimum height=1.2cm, align=center, font=\bfseries] (thesis) at (0,4) {THÈSE\\ \textit{(Idée directrice, jugement à faire admettre)}};

    % Niveau 2 : Arguments
    \node[draw=argcol, thick, fill=argcol!15, rounded corners=6pt, minimum width=10cm, minimum height=1.4cm, align=center, font=\bfseries] (arg) at (0,2.2) {ARGUMENTS\\ \textit{(Raisons logiques, preuves, motifs : « Pourquoi ? »)}};

    % Niveau 3 : Exemples (base large)
    \node[draw=excol, thick, fill=excol!15, rounded corners=6pt, minimum width=12cm, minimum height=1.6cm, align=center, font=\bfseries] (ex) at (0,0.2) {EXEMPLES / ILLUSTRATIONS\\ \textit{(Faits concrets, anecdotes, stats, citations : « Par exemple ? »)}};

    % Flèches de soutien
    \draw[thick, -{Stealth[length=3mm]}, popmuted] (thesis.south) -- node[right, font=\tiny, pos=0.5, xshift=0.5cm] {\textit{soutient}} (arg.north);
    \draw[thick, -{Stealth[length=3mm]}, popmuted] (arg.south) -- node[right, font=\tiny, pos=0.5, xshift=0.5cm] {\textit{illustre}} (ex.north);

    % Conclusion (latérale ou bas)
    \node[draw=conclcol, thick, fill=conclcol!15, rounded corners=6pt, minimum width=5cm, minimum height=1cm, align=center, font=\bfseries\itshape] (concl) at (0,-1.8) {CONCLUSION\\ \textit{(Réaffirmation / Ouverture)}};
    \draw[thick, -{Stealth[length=3mm]}, popmuted, dashed] (ex.south) -- (concl.north);

    % Annotation "Contraction"
    \node[draw=poporange, thick, fill=poporange!15, rounded corners=4pt, align=left, font=\small, anchor=north west] at (6.5, 3.5) {\textbf{Contraction :}\\
    \textcolor{thesiscol}{\textbf{Garder}} (reformulée)\\
    \textcolor{argcol}{\textbf{Garder}} (synthétisés)\\
    \textcolor{excol}{\textbf{Supprimer}} ou \textbf{Généraliser}};
\end{tikzpicture}
\legende{Schéma pyramidal de la structure argumentative. La contraction conserve le sommet et le niveau intermédiaire ; elle élimine ou généralise la base (exemples).}
\end{popfigure}

\courssub{Les trois composantes en détail}

\courssub{La Thèse : le « quoi »}
C'est le \textbf{cœur nucléaire} du texte. Elle peut être :
\begin{itemize}
    \item \textbf{Explicite} : formulée clairement, souvent en début (thèse d'attaque) ou en fin (thèse de conclusion) de texte. Repérables via : \emph{« Je pense que... », « Il est évident que... », « Donc... », « C'est pourquoi... »}.
    \item \textbf{Implicite} : non énoncée directement, le lecteur doit l'inférer à partir de l'ensemble des arguments.
\end{itemize}

\begin{methode}[Repérer la thèse]
\begin{enumerate}
    \item Cherchez les \textbf{connecteurs de conclusion} (\cle{donc}, \cle{par conséquent}, \cle{c'est pourquoi}, \cle{ainsi}, \cle{en définitive}).
    \item Posez la question : \textbf{« Que veut me faire penser / croire / faire l'auteur ? »}
    \item Vérifiez si l'énoncé trouvé \textbf{regroupe tous les arguments} du texte (test de la généralité).
\end{enumerate}
\end{methode}

\courssub{Les Arguments : le « pourquoi »}
Ce sont les piliers logiques. On distingue classiquement quatre types (recoupant les catégories rhétoriques) :

\begin{center}
\begin{tabularx}{\linewidth}{|l|X|X|}\hline
\rowcolor{popblueL}
\textbf{Type d'argument} & \textbf{Définition / Mécanisme} & \textbf{Indices de repérage} \\ \hline
\textbf{Argument d'autorité} & S'appuie sur un expert, un témoin reconnu, un texte sacré, la loi, la tradition. & « Selon X... », « Comme l'écrit Y... », « La loi stipule... », « Nos ancêtres disaient... » \\ \hline
\textbf{Argument de l'expérience / fait} & S'appuie sur l'observation directe, des statistiques, des données chiffrées, un vécu. & « J'ai vu... », « Les chiffres montrent... », « Il est prouvé que... », « Dans notre village... » \\ \hline
\textbf{Argument de valeur} & Mobilise des valeurs partagées (justice, honneur, solidarité, progrès, tradition, liberté). & « C'est injuste... », « L'honneur commande... », « Au nom de la solidarité... », « Pour le bien commun... » \\ \hline
\textbf{Argument logique (déductif/inductif)} & Liaison rationnelle cause/effet, analogie, \textit{a fortiori}, \textit{reductio ad absurdum}. & « Puisque... alors... », « Si A alors B », « C'est contradictoire car... », « Par analogie... » \\ \hline
\end{tabularx}
\end{center}

\begin{attention}[Piège : Confondre Argument et Exemple]
Un argument est une \textbf{raison générale} ; un exemple est une \textbf{occurrence particulière}.
\begin{itemize}
    \item \textit{Argument :} « La corruption freine le développement. » (Raison générale)
    \item \textit{Exemple :} « En 2022, le projet de barrage de X a pris 5 ans de retard à cause de détournements. » (Cas singulier)
\end{itemize}
Dans une contraction, on garde l'argument, on supprime l'exemple (ou on le généralise : « Des détournements retardent les grands chantiers »).
\end{attention}

\courssub{Les Exemples : le « concret »}
Ils ancrent l'abstrait dans le réel. Ils sont \textbf{secondaires} par rapport à l'argument qu'ils illustrent. Repérables via : \cle{par exemple}, \cle{ainsi}, \cle{notamment}, \cle{tel est le cas de}, \cle{citons}, \cle{prenons le cas de}.

\courssub{Les connecteurs logiques : le ciment du raisonnement}

Les \textbf{connecteurs logiques} (ou opérateurs argumentatifs) balisent le cheminement de la pensée. Les identifier, c'est voir le squelette.

\begin{center}
\begin{tabularx}{\linewidth}{|l|X|l|}\hline
\rowcolor{popblueL}
\textbf{Fonction} & \textbf{Connecteurs principaux} & \textbf{Effet} \\ \hline
\textbf{Enchaînement / Ajout} & d'abord, ensuite, de plus, par ailleurs, également, en outre & Accumule les arguments \\ \hline
\textbf{Explication / Cause} & car, parce que, puisque, comme, en effet, du fait que & Justifie l'argument précédent \\ \hline
\textbf{Conséquence / Conclusion} & donc, par conséquent, c'est pourquoi, ainsi, alors, en conclusion & Déduit la thèse ou un argument intermédiaire \\ \hline
\textbf{Opposition / Concession} & mais, cependant, pourtant, nevertheless, bien que, si... toutefois & Nuance, réfute une objection (argument concessif) \\ \hline
\textbf{Illustration} & par exemple, ainsi, notamment, c'est le cas de, prenons le cas de & Introduit un exemple \\ \hline
\textbf{Réforme / Correction} & en réalité, en fait, au contraire, loin de là & Rectifie une idée reçue (ironie/paradoxe) \\ \hline
\end{tabularx}
\end{center}

\begin{exemplebox}[Connecteurs dans \textit{Le Vieux nègre et la médaille}]
Dans le discours du Commandant (extrait étudié en séance 1), on repère :
\begin{itemize}
    \item \textbf{« D'abord »} (énumération des « bienfaits » de la colonisation) $\rightarrow$ \textit{Enchaînement}.
    \item \textbf{« Car »} / \textbf{« Puisque »} (justification de l'ordre établi) $\rightarrow$ \textit{Cause}.
    \item \textbf{« Cependant »} (lorsqu'il concède une souffrance pour mieux la minimiser) $\rightarrow$ \textit{Concession rhétorique}.
    \item \textbf{« C'est pourquoi »} (justification de la médaille) $\rightarrow$ \textit{Conclusion partielle}.
\end{itemize}
Ces connecteurs révèlent une structure \emph{en apparence} logique mais \emph{rhétoirement} manipulatoire (ironie du texte).
\end{exemplebox}

\courssub{Application guidée : Analyse structurelle d'extraits du \textit{Vieux nègre et la médaille}}

Nous appliquons la méthode à deux passages clés : le discours officiel du Commandant (séance 1) et la réflexion finale de Meka (séance 6).

\begin{exoresolu}[Analyse du discours du Commandant (Extrait : remise de la médaille)]
\textbf{Texte (reconstruit) :} « D'abord, mes amis, cette médaille récompense cinquante ans de fidélité. Car vous avez servi la France sans faillir. Puisque vous avez porté l'uniforme avec honneur, c'est pourquoi la République vous honore aujourd'hui. Cependant, certains disent que le colonisateur a apporté la misère. En réalité, il a apporté l'école, la route, la médecine. Prenons le cas de l'hôpital de Yaoundé : il soigne des milliers de Noirs. Donc, cette médaille est le symbole de notre amitié indéfectible. »

\textbf{Analyse structurelle :}

\begin{center}
\begin{tabularx}{\linewidth}{|l|X|}\hline
\rowcolor{popblueL}
\textbf{Composante} & \textbf{Contenu repéré} \\ \hline
\textbf{THÈSE (Explicite, finale)} & « Cette médaille est le symbole de notre amitié indéfectible » / La colonisation est bienfaitrice. \\ \hline
\textbf{ARGUMENT 1 (Autorité / Valeur)} & Fidélité, service, honneur (valeurs militaires/coloniales). \textit{Connecteur : « D'abord », « Car »}. \\ \hline
\textbf{ARGUMENT 2 (Logique / Fait)} & Apports matériels : école, route, médecine. \textit{Connecteur : « En réalité » (réfutation objection)}. \\ \hline
\textbf{EXEMPLE (Illustration Arg 2)} & « L'hôpital de Yaoundé soigne des milliers de Noirs ». \textit{Connecteur : « Prenons le cas de »}. \\ \hline
\textbf{CONCLUSION} & « Donc, cette médaille... » (Réaffirmation thèse). \\ \hline
\end{tabularx}
\end{center}

\textbf{Observation critique :} La structure \emph{simule} une démonstration rationnelle (Thèse $\leftarrow$ Arguments $\leftarrow$ Exemple), mais les arguments sont biaisés (sélection partielle des faits, occultation de la violence coloniale). L'ironie de Mongo Beti réside dans l'écart entre cette structure formelle « parfaite » et la réalité vécue par Meka.
\tcblower
\textbf{Solution détaillée (Modélisation pour la contraction) :}
\begin{itemize}
    \item \textbf{Idée essentielle 1 (Thèse) :} Le Commandant présente la médaille comme symbole d'amitié franco-africaine et récompense de la fidélité.
    \item \textbf{Idée essentielle 2 (Argument synthétisé) :} Il justifie cela par les « bienfaits » de la colonisation (éducation, santé, infrastructures).
    \item \textbf{À éliminer :} L'exemple de l'hôpital de Yaoundé (détail illustratif), les connecteurs de rhétorique orale (« mes amis », « D'abord »).
\end{itemize}
\end{exoresolu}

\begin{exoresolu}[Analyse de la réflexion finale de Meka (Extrait : nuit après la cérémonie)]
\textbf{Texte (reconstruit) :} « Meka regarda la médaille. Elle brillait. Mais son fils était mort. Sa femme pleurait. Ses terres étaient prises. La médaille ne mange pas. Elle ne soigne pas. Elle ne ramène pas les morts. C'est un morceau de fer. Donc, il l'a jetée dans la latrine. »

\textbf{Analyse structurelle :}

\begin{center}
\begin{tabularx}{\linewidth}{|l|X|}\hline
\rowcolor{popblueL}
\textbf{Composante} & \textbf{Contenu repéré} \\ \hline
\textbf{THÈSE (Implicite, finale : action = thèse vécue)} & La médaille est une insulte, un symbole vide de la domination coloniale ; elle doit être rejetée. \\ \hline
\textbf{ARGUMENT 1 (Expérience / Valeur : Vie/Mort)} & La médaille ne résout pas les souffrances concrètes (mort du fils, deuil, spoliation). \textit{Connecteur : « Mais » (opposition forte)}. \\ \hline
\textbf{ARGUMENT 2 (Logique / Utilitariste)} & Inutilité matérielle : « ne mange pas, ne soigne pas, ne ramène pas les morts ». \\ \hline
\textbf{EXEMPLES (Vécus directs = arguments forts ici)} & Mort du fils, larmes de la femme, terres prises. \textit{Note : ici, le « fait vécu » tient lieu d'argument probant, pas de simple illustration.} \\ \hline
\textbf{CONCLUSION (Acte performatif)} & Jeter la médaille dans la latrine (rejet total, performatif). \\ \hline
\end{tabularx}
\end{center}

\textbf{Distinction cruciale :} Chez Meka, les « exemples » (mort du fils, terres) \textbf{sont} les arguments (arguments d'expérience vécue irréfutables). Ils ne sont \textbf{PAS} supprimables en contraction, contrairement à l'exemple de l'hôpital chez le Commandant. La contraction doit garder : « Meka rejette la médaille car elle ne répare pas ses drames personnels (mort fils, spoliation) ».
\tcblower
\textbf{Solution détaillée (Modélisation pour la contraction) :}
\begin{itemize}
    \item \textbf{Idée essentielle 1 (Thèse/Action) :} Meka jette la médaille dans la latrine, la refusant comme symbole mensonger.
    \item \textbf{Idée essentielle 2 (Arguments fondateurs) :} Ce rejet est motivé par l'inanité de la décoration face à ses tragédies personnelles (deuil, spoliation) et l'incapacité de l'objet à satisfaire des besoins vitaux.
    \item \textbf{À garder (généralisé) :} « Face à la mort de son fils et à la perte de ses terres, la médaille lui apparaît comme un morceau de fer inutile. »
\end{itemize}
\end{exoresolu}

\courssub{Distinction Idées Essentielles vs Idées Secondaires : la clé de la contraction}

C'est l'application directe de l'analyse structurelle. La contraction consiste à \textbf{conserver la structure porteuse (Thèse + Arguments)} et à \textbf{traiter les éléments périphériques}.

\begin{aretenir}[Règle d'or pour la contraction]
\begin{center}
\begin{tabular}{|c|c|c|}\hline
\rowcolor{popblueL}
\textbf{Niveau structurel} & \textbf{Statut} & \textbf{Traitement en contraction} \\ \hline
\textbf{Thèse} & \textbf{ESSENTIELLE} (Idée directrice) & \textbf{Reformuler} (synthèse en 1 phrase) \\ \hline
\textbf{Arguments} & \textbf{ESSENTIELS} (Piliers) & \textbf{Regrouper / Synthétiser} (fusionner arguments proches) \\ \hline
\textbf{Exemples / Illustrations} & \textbf{SECONDAIRES} (Détails) & \textbf{Supprimer} OU \textbf{Généraliser} (remplacer par une classe) \\ \hline
\textbf{Répétitions / Reformulations} & \textbf{SECONDAIRES} (Redondance) & \textbf{Supprimer} \\ \hline
\textbf{Digressions / Incises} & \textbf{SECONDAIRES} (Hors sujet) & \textbf{Supprimer} \\ \hline
\textbf{Connecteurs de structure} & \textbf{OUTILS} (Squelette) & \textbf{Conserver / Adapter} pour cohérence \\ \hline
\end{tabular}
\end{center}
\end{aretenir}

\begin{exemplebox}[Exemple : Éditorial \textit{Cameroun Tribune} (Incivisme)]
\textbf{Texte source (résumé) :} « L'incivisme gangrène notre société (Thèse). D'abord, les dégradations de biens publics coûtent des milliards (Arg 1 : Économique). Par exemple, le pont de la Banane a été vandalisé l'an dernier (Ex 1). Ensuite, l'incivilisme routier tue (Arg 2 : Humain/Sécurité). Ainsi, 300 morts sur la route Yaoundé-Douala en 2023 (Ex 2 : Stat). Enfin, l'incurie administrative freine le développement (Arg 3 : Développement). C'est pourquoi chacun doit changer (Conclusion). »

\textbf{Contraction (cible : 30 mots) :} 
\begin{quote}
\textit{« L'éditorial dénonce l'incivisme (dégradations, insécurité routière, frein au développement) qui coûte vies et argent, appelant à un changement de comportement. »}
\end{quote}
\textbf{Analyse des choix :}
\begin{itemize}
    \item \textbf{Gardés :} Thèse (incivisme), 3 arguments regroupés (biens publics, routes, admin), conclusion (appel).
    \item \textbf{Supprimés :} Exemple précis « pont de la Banane », chiffre précis « 300 morts / 2023 / Yaoundé-Douala » $\rightarrow$ généralisés en « dégradations », « insécurité routière », « vies ».
\end{itemize}
\end{exemplebox}

\begin{attention}[Erreurs fatales en contraction]
\begin{enumerate}
    \item \textbf{Garder un exemple, supprimer l'argument.} $\rightarrow$ Le résumé devient une liste de faits sans logique.
    \item \textbf{Confondre thèse et argument.} $\rightarrow$ On répète la même idée au lieu de montrer l'enchaînement.
    \item \textbf{Conserver les connecteurs oraux/ponctuels} (« mes amis », « écoutez bien », « figurez-vous »).
    \item \textbf{Oublier la conclusion} si elle apporte une nuance ou un appel à l'action nouveau.
\end{enumerate}
\end{attention}

\courssub{Méthode opératoire complète : De la lecture au plan de contraction}

\begin{methode}[Analyse structurelle en 5 étapes (Fiche méthode Bac)]
\begin{enumerate}
    \item \textbf{Lecture active :} Souligner (crayon rouge) la \textbf{Thèse} (début ou fin). Se demander : « Quel est le \cle{point de vue} défendu ? »
    \item \textbf{Numérotation des arguments :} Souligner (crayon bleu) les \textbf{arguments} principaux. Numéroter dans la marge : \textcircled{1}, \textcircled{2}, \textcircled{3}... Vérifier les connecteurs (\textit{d'abord, ensuite, de plus, cependant}).
    \item \textbf{Encadrement des exemples :} Encadrer (crayon vert) les \textbf{exemples/illustrations} (\textit{par exemple, cas de, tel que}). Noter « EX » à côté.
    \item \textbf{Évaluation de la valeur informative :} Pour chaque élément vert (exemple) : \textbf{Est-il unique ? Remplaçable par une généralité ?} $\rightarrow$ OUI = Supprimer. NON (fait historique majeur, expérience vécue centrale comme chez Meka) $\rightarrow$ Généraliser et garder comme argument.
    \item \textbf{Construction du plan de contraction :} Rédiger le \textbf{plan détaillé} (Thèse + Arguments synthétisés) avant d'écrire la phrase finale.
\end{enumerate}
\end{methode}

\begin{experience}[Atelier pratique : Annotation collaborative]
\textbf{Consigne :} Sur un extrait de 15 lignes (discours politique ou article d'opinion), en binôme :
\begin{enumerate}
    \item Identifier la thèse (surbrillance \textcolor{popred}{ROUGE}).
    \item Identifier 3 arguments (surbrillance \textcolor{popblue}{BLEU} + numéro).
    \item Identifier 2 exemples (surbrillance \textcolor{popgreen}{VERT} + « EX »).
    \item Rédiger le plan de contraction (3 lignes max).
    \item Échanger avec un autre binôme : le plan permet-il de reconstruire le sens sans le texte original ?
\end{enumerate}
\textbf{Matériel :} Extrait photocopié, 3 surligneurs (rouge, bleu, vert), papier brouillon.
\textbf{Durée :} 20 min.
\end{experience}

\begin{savaistu}[Argumentation et cultures africaines : le \textit{Palaver}]
La structure thèse-arguments-exemples n'est pas une invention occidentale. La tradition du \textbf{Palaver} (ou \textit{Mbar}, \textit{Ngondo}, \textit{Ébiosa} selon les aires culturelles) repose sur une rhétorique codifiée :
\begin{itemize}
    \item \textbf{Thèse} : La plainte ou la proposition initiale de l'aîné ou du plaignant.
    \item \textbf{Arguments} : Les proverbes (\textit{« L'enfant qui ne pleure pas meurt dans le dos »}), les précédents coutumiers (jurisprudence orale = argument d'autorité), l'analyse des faits (argument d'expérience).
    \item \textbf{Exemples} : Récits mythiques, généalogies, événements vécus par le clan.
    \item \textbf{Conclusion} : Verdict du conseil des sages (souvent consensus / réconciliation, non victoire binaire).
\end{itemize}
Comprendre cette structure endogène aide à analyser les textes argumentatifs africains (comme \textit{Le Vieux nègre...} où Meka opère un \textit{contre-palaver} intérieur) et à ne pas plaquer aveuglément des grilles purement aristotéliciennes.
\end{savaistu}

\begin{exercice}[Entraînement : Structure et contraction]
\textbf{Texte :} « Il est urgent de réformer l'enseignement technique au Cameroun. En effet, les programmes sont obsolètes depuis 20 ans. Par exemple, on y enseigne encore la réparation de machines à écrire alors que les bureaux n'utilisent que des ordinateurs. De plus, les ateliers manquent d'équipements. Ainsi, au Lycée Technique de Nkongsamba, 50 élèves partagent un seul tour à bois. C'est pourquoi le chômage des jeunes techniciens dépasse 60\%. Donc, l'État doit investir massivement dans la modernisation des plateaux techniques et la révision des curricula. »

\textbf{Questions :}
\begin{enumerate}
    \item Repérez la thèse, les arguments (type), les exemples. Coloriez (R/B/V) sur votre copie.
    \item Proposez un plan de contraction (Thèse + Arguments synthétisés).
    \item Rédigez la contraction en \textbf{30 mots maximum}.
\end{enumerate}
\end{exercice}

\begin{corrige}[Exercice n°1]
\begin{enumerate}
    \item \textbf{Thèse (Rouge) :} « Il est urgent de réformer l'enseignement technique au Cameroun. » (Début) / « L'État doit investir... » (Fin = thèse d'action).
    \item \textbf{Argument 1 (Bleu) :} Obsolescence programmes (Argument logique/fait). \textit{Exemple (Vert) :} Machines à écrire vs ordinateurs.
    \item \textbf{Argument 2 (Bleu) :} Manque équipements ateliers (Argument d'expérience/fait). \textit{Exemple (Vert) :} Lycée Nkongsamba, 50 élèves / 1 tour.
    \item \textbf{Argument 3 (Bleu) :} Conséquence : chômage jeunes (Argument logique/conséquence). Chiffre 60\% (Donnée stat = illustration argument 3).
    \item \textbf{Plan contraction :} Thèse (Urgence réforme) $\rightarrow$ Arg 1 (Programmes décalés marché) $\rightarrow$ Arg 2 (Équipements insuffisants) $\rightarrow$ Arg 3 (Chômage massif) $\rightarrow$ Conclusion (Investissement État).
    \item \textbf{Contraction (28 mots) :} \textit{« Le texte réclame une réforme urgente de l'enseignement technique : programmes obsolètes, ateliers sous-équipés provoquent un chômage massif des techniciens ; l'État doit moderniser cursus et matériel. »}
\end{enumerate}
\end{corrige}

\courssub{Synthèse : Vers la production du résumé}

L'analyse de la structure argumentative n'est pas un exercice académique gratuit : c'est le \textbf{moteur de la contraction}.
\begin{itemize}
    \item \textbf{Sans thèse identifiée} $\rightarrow$ Résumé sans fil conducteur (liste de phrases).
    \item \textbf{Sans arguments hiérarchisés} $\rightarrow$ Résumé déséquilibré (trop de détails sur un point, oubli d'un pan entier).
    \item \textbf{Sans tri exemples/arguments} $\rightarrow$ Résumé trop long (non contraction) ou anecdotique.
\end{itemize}

La section suivante (« Techniques de réduction et reformulation des idées essentielles ») nous donnera les outils linguistiques (nominalisation, subordination, suppression du superflu, synonymie) pour transformer ce \textbf{plan structurel} en un \textbf{texte contracté fluide, cohérent et respectueux de la limite de mots}.

\begin{aretenir}[Check-list avant d'écrire la contraction]
\renewcommand{\labelitemi}{$\square$}
\begin{itemize}
    \item Ai-je formulé la \textbf{thèse} en une phrase claire ?
    \item Ai-je listé \textbf{tous les arguments distincts} (sans confondre arg 1 et arg 2) ?
    \item Ai-je \textbf{éliminé tous les exemples} (sauf s'ils \textbf{sont} l'argument, cas Meka) ?
    \item Ai-je noté les \textbf{connecteurs logiques} à réutiliser pour la cohérence ?
    \item Le nombre d'idées essentielles est-il compatible avec la \textbf{limite de mots} imposée ?
\end{itemize}
\end{aretenir}

\coursec{Techniques de réduction et reformulation des idées essentielles}

Dans la section précédente, nous avons appris à repérer la \cle{structure argumentative} d'un texte : la thèse de l'auteur, les arguments qui la soutiennent, les exemples qui les illustrent. Ce travail d'analyse n'est pas une fin en soi : il prépare directement l'exercice que nous abordons maintenant, la \cle{contraction de texte}. En effet, on ne peut réduire un texte que si l'on sait distinguer l'essentiel (thèse et arguments) de l'accessoire (exemples, détails, répétitions). La contraction est un exercice classique du Probatoire et du Baccalauréat technique : on vous donne un texte et une consigne du type « Résumez ce texte en 100 mots ». Ce chapitre vous donne les outils techniques pour y parvenir.

\courssub{Le principe de la contraction}

\begin{definition}[La contraction de texte]
La \cle{contraction de texte} (ou résumé) est un exercice qui consiste à réduire un texte à une fraction imposée de sa longueur (souvent le quart, noté 1/4), en conservant toutes les idées essentielles de l'auteur, sans rien ajouter (ni avis personnel, ni exemple nouveau) et sans rien trahir (ni déformation du sens, ni oubli d'une idée importante).
\end{definition}

L'intuition est simple : imaginez que vous devez raconter à un ami pressé un match de football de 90 minutes en deux minutes. Vous ne pouvez pas tout dire ; vous direz le score, les buteurs, le moment décisif. Vous supprimerez les descriptions des passes, du public, de la météo. Mais attention : vous ne devez ni inventer un but qui n'a pas eu lieu, ni donner votre avis sur l'arbitre. La contraction fonctionne exactement ainsi, mais avec des règles précises et des techniques codifiées.

\begin{aretenir}[Les trois lois de la contraction]
\begin{enumerate}
\item \textbf{Fidélité} : le résumé respecte le sens, la thèse et l'ordre des idées du texte original.
\item \textbf{Neutralité} : le résumé ne contient aucune idée, aucun jugement, aucun exemple qui ne figure pas dans le texte.
\item \textbf{Reformulation} : le résumé est rédigé avec les mots du candidat, non avec les phrases de l'auteur.
\end{enumerate}
\end{aretenir}

\courssub{Les trois techniques de réduction}

Pour passer d'un texte long à un texte court, le candidat dispose de trois outils complémentaires, qu'il utilise presque toujours ensemble.

\textbf{Technique n° 1 : la suppression.}\par
La \cle{suppression} consiste à éliminer purement et simplement tout ce qui n'est pas indispensable à la compréhension des idées essentielles. On supprime :
\begin{itemize}
\item les \textbf{exemples} et les illustrations concrètes (on garde l'idée qu'ils illustrent) ;
\item les \textbf{répétitions} et les reformulations d'une même idée par l'auteur ;
\item les \textbf{détails} descriptifs, anecdotiques, les précisions de lieu, de date, de couleur ;
\item les \textbf{figures de style} (comparaisons, métaphores, hyperboles) qui ornent le propos ;
\item les \textbf{énumérations} longues, que l'on remplace par un terme générique (voir technique n° 2) ;
\item les incises, les parenthèses, les adverbes d'intensité (« extrêmement », « vraiment »).
\end{itemize}

\textbf{Technique n° 2 : la généralisation.}\par
La \cle{généralisation} consiste à remplacer un terme spécifique (ou une liste de termes spécifiques) par un terme générique qui les englobe tous. C'est le passage du particulier au général.

\begin{exemplebox}[La généralisation en action]
\begin{itemize}
\item « manguiers, goyaviers et palmiers » $\rightarrow$ « arbres fruitiers » ;
\item « son père, sa mère, ses oncles et ses cousins » $\rightarrow$ « sa famille » ;
\item « il marcha, courut, sauta par-dessus la clôture » $\rightarrow$ « il s'enfuit » ;
\item « la houe, la machette et la pioche » $\rightarrow$ « les outils agricoles ».
\end{itemize}
\end{exemplebox}

\textbf{Technique n° 3 : la condensation.}\par
La \cle{condensation} consiste à dire la même chose en moins de mots, en jouant sur la grammaire : on fusionne plusieurs phrases en une seule, on remplace une proposition subordonnée par un nom ou un groupe nominal, un verbe par un nom d'action (nominalisation).

\begin{center}
\begin{tabularx}{\linewidth}{|l|X|X|}\hline
\rowcolor{popblueL} \textbf{Transformation} & \textbf{Forme longue} & \textbf{Forme condensée} \\ \hline
Subordonnée de cause & parce qu'il était fatigué & par fatigue \\ \hline
Subordonnée de temps & lorsqu'il arriva au village & à son arrivée au village \\ \hline
Nominalisation & il résista aux colons & sa résistance aux colons \\ \hline
Deux phrases fusionnées & Meka était vieux. Il était épuisé par le travail. & Meka, vieillard épuisé, \\ \hline
Relative & la médaille qui lui fut décernée & la médaille reçue \\ \hline
\end{tabularx}
\end{center}

\begin{attention}[Piège de la condensation]
Condenser ne signifie pas supprimer une idée. « Parce qu'il était fatigué » devient « par fatigue » : l'idée de cause est conservée. Si vous supprimez totalement la cause, vous appauvrissez le raisonnement de l'auteur et vous perdez des points de contenu.
\end{attention}

\courssub{La reformulation : dire autrement les idées essentielles}

\begin{definition}[La reformulation]
La \cle{reformulation} est la réexpression d'une idée avec d'autres mots et d'autres constructions syntaxiques, sans aucun changement de sens. Elle prouve au correcteur que le candidat a réellement compris le texte.
\end{definition}

La reformulation peut porter sur trois plans, que l'on combine :
\begin{enumerate}
\item \textbf{Le vocabulaire} : remplacer les mots par des synonymes ou des équivalents (« vieillard » $\rightarrow$ « vieil homme » ; « labeur » $\rightarrow$ « travail » ; « décerna » $\rightarrow$ « remit »).
\item \textbf{La construction syntaxique} : transformer la phrase (une phrase complexe devient simple ; une subordonnée devient un nom ; deux phrases n'en font plus qu'une).
\item \textbf{La voix} : passer de la voix passive à la voix active ou inversement (« la médaille lui fut décernée par le commandant » $\rightarrow$ « le commandant lui décerna la médaille »).
\end{enumerate}

\begin{exemplebox}[Exemple chiffré tiré du \emph{Vieux nègre et la médaille}]
Phrase originale (23 mots) : « Meka, vieillard épuisé par des années de labeur et de soumission, reçut avec une fierté naïve la médaille que lui décerna le commandant. »\par
Version contractée (5 mots) : « Meka reçut fièrement une médaille. »\par
Techniques utilisées : suppression des expansions (« vieillard épuisé par des années de labeur et de soumission », « que lui décerna le commandant »), condensation de « avec une fierté naïve » en « fièrement », reformulation du verbe.
\end{exemplebox}

\begin{attention}[Le copiage est sanctionné]
L'erreur la plus fréquente et la plus grave consiste à recopier des phrases entières du texte. Au Probatoire et au Baccalauréat technique, le copiage est lourdement sanctionné : le correcteur souligne les passages recopiés et ne les compte pas. Seuls peuvent être conservés les \textbf{mots-outils} (articles, prépositions, conjonctions) et les \textbf{termes irremplaçables} (noms propres comme « Meka », notions techniques sans synonyme). Tout le reste doit être reformulé.
\end{attention}

\courssub{Illustration : le passage de l'original au contracté}

\begin{popfigure}
\begin{center}
\begin{tikzpicture}[
  boite/.style={draw=popdark, rounded corners=3pt, fill=popcream, text width=0.62\linewidth, align=left, font=\small, inner sep=6pt},
  tech/.style={font=\footnotesize\itshape, text=popblue, midway, sloped, above}
]
\node[boite] (avant) at (0,0) {\textbf{Avant (31 mots)} : « Le vieux Meka, qui avait travaillé toute sa vie comme un esclave pour les colons, éprouva une joie immense et sincère lorsque le commandant du district lui remit solennellement sa médaille. »};
\node[boite, fill=popgreenL] (apres) at (0,-4.2) {\textbf{Après (8 mots)} : « Meka, serviteur des colons, reçut joyeusement sa médaille. »};
\draw[-{Stealth[length=3mm]}, very thick, poporange] (avant.south west) ++(0.5,0) -- ++(0,-1.1) node[tech]{suppression des détails};
\draw[-{Stealth[length=3mm]}, very thick, popblue] (avant.south) -- ++(0,-1.4) node[tech]{condensation (relative $\rightarrow$ nom)};
\draw[-{Stealth[length=3mm]}, very thick, poppurple] (avant.south east) ++(-0.5,0) -- ++(0,-1.1) node[tech]{reformulation};
\end{tikzpicture}
\end{center}
\legende{Schéma « avant / après » : une phrase de 31 mots est réduite à 8 mots grâce à la combinaison des trois techniques de réduction et de la reformulation.}
\end{popfigure}

\courssub{La méthode complète de la contraction}

\begin{methode}[Les 4 étapes de la contraction]
\begin{enumerate}
\item \textbf{Lire et comprendre} : lire le texte deux fois ; repérer le thème, la thèse de l'auteur et le ton ; deviner le sens des mots inconnus grâce au contexte.
\item \textbf{Dégager le plan et les idées essentielles} : souligner les articulations logiques (d'abord, ensuite, mais, donc) ; noter au brouillon une idée essentielle par partie du texte.
\item \textbf{Réduire} : appliquer à chaque phrase les trois techniques — suppression de l'accessoire, généralisation des énumérations, condensation des constructions lourdes.
\item \textbf{Reformuler et vérifier} : rédiger le résumé avec ses propres mots, compter les mots, ajuster si nécessaire, puis relire pour vérifier la fidélité et corriger l'orthographe.
\end{enumerate}
\end{methode}

\courssub{Le comptage des mots et la marge de tolérance}

La consigne impose toujours un nombre de mots. Il faut savoir compter vite et juste.

\begin{methode}[Compter les mots efficacement]
\begin{enumerate}
\item Comptez les mots d'une ligne complète et moyenne de votre résumé (par exemple 11 mots).
\item Multipliez par le nombre de lignes pleines (11 mots $\times$ 13 lignes = 143 mots).
\item Ajoutez les mots de la dernière ligne incomplète.
\item Rappel : un mot composé à trait d'union (« peut-être ») compte pour \textbf{un} mot ; les noms propres composés se comptent normalement.
\end{enumerate}
\end{methode}

\begin{attention}[Respecter la limite à $\pm$ 10 \%]
Le nombre de mots imposé admet une marge de tolérance de 10 \% en plus ou en moins. Ainsi, un résumé de 150 mots demandé doit compter entre \textbf{135 et 165 mots} ; un résumé de 100 mots, entre 90 et 110 mots. En dessous, vous avez supprimé des idées essentielles ; au-dessus, vous avez gardé de l'accessoire. Dans les deux cas, vous êtes sanctionné.
\end{attention}

\courssub{Entraînement guidé : contraction d'un extrait de 200 mots en 50 mots}

Travaillons ensemble, étape par étape, comme au tableau, sur un passage inspiré du \emph{Vieux nègre et la médaille} de Ferdinand Oyono.

\begin{exoresolu}[Contraction guidée d'un extrait]
\textbf{Énoncé.} Résumez en 50 mots environ le passage suivant (200 mots environ) : « Meka avait grandi dans le village de Mfoudi, un petit hameau perdu au milieu des collines, des forêts épaisses et des rivières impétueuses. Toute sa vie, il avait travaillé la terre : il avait défriché, semé, sarclé, récolté, saison après saison, année après année, sans jamais se plaindre. Ses mains étaient devenues dures comme l'écorce des vieux arbres, son dos s'était courbé comme un arc sous le poids des années et des charges portées. Il avait donné ses terres aux missionnaires, parce qu'il croyait sincèrement que les Blancs apportaient la civilisation et le salut. Il avait même perdu ses deux fils, partis à la guerre pour défendre la France, cette patrie lointaine qu'il n'avait jamais vue et qu'il imaginait comme un paradis. Quand le commandant du district arriva au village, un matin de juillet, accompagné de ses gardes, de ses secrétaires et de tout son cortège, pour lui remettre la médaille de l'amitié franco-camerounaise, Meka ressentit une fierté immense : il se sentait enfin reconnu, enfin considéré comme un ami des Blancs, lui le vieux nègre que personne n'avait jamais regardé. »
\tcblower
\textbf{Solution détaillée.}\par
\textbf{Étape 1 — Lecture.} Thème : la vie de Meka et la remise de la médaille. Ton : récit ironique (la « reconnaissance » des Blancs est illusoire).\par
\textbf{Étape 2 — Idées essentielles.} (a) Meka a travaillé toute sa vie ; (b) il a tout donné aux Blancs (terres, fils) par confiance ; (c) la médaille le remplit de fierté car il se croit enfin reconnu.\par
\textbf{Étape 3 — Réduction.} Suppression des détails (Mfoudi, collines, juillet, le cortège), généralisation (« défriché, semé, sarclé, récolté » $\rightarrow$ « cultivé » ; « gardes, secrétaires, cortège » supprimé), condensation (« parce qu'il croyait » $\rightarrow$ « croyant » ; « ses mains dures comme l'écorce » supprimée comme image).\par
\textbf{Étape 4 — Résumé rédigé (51 mots)} : « Meka, paysan camerounais, a peiné toute sa vie sans se plaindre. Confiant dans les Blancs, il leur a tout sacrifié : ses terres aux missionnaires et ses deux fils morts pour la France. La médaille reçue du commandant le comble de fierté : il se croit enfin reconnu et accepté par les colons. » (51 mots : dans la marge de 45 à 55 mots.)
\end{exoresolu}

\begin{popfigure}
\begin{center}
\begin{tikzpicture}
\begin{axis}[
  ybar, bar width=0.9cm,
  width=0.85\linewidth, height=6.5cm,
  ymin=0, ymax=230,
  ylabel={Nombre de mots},
  symbolic x coords={Texte original, Suppression, G\'en\'eralisation, Condensation, R\'esum\'e final},
  xtick=data,
  x tick label style={font=\small, rotate=15, anchor=east},
  nodes near coords,
  every node near coord/.append style={font=\small\bfseries, color=popdark},
  axis lines=left,
  ymajorgrids=true, grid style={popline!40},
]
\addplot[fill=popblue, draw=popdark] coordinates {(Texte original,200) (Suppression,120) (G\'en\'eralisation,85) (Condensation,62) (R\'esum\'e final,50)};
\end{axis}
\end{tikzpicture}
\end{center}
\legende{Diagramme en barres : diminution progressive du nombre de mots à chaque étape de la réduction, de 200 mots (texte original) à 50 mots (résumé final), soit le quart imposé.}
\end{popfigure}

\courssub{Entraînement personnel}

\begin{exercice}[À vous de contracter]
Résumez en 40 mots environ le passage suivant (160 mots environ) : « Après la cérémonie, Meka rentra chez lui, le cœur léger, la médaille brillante épinglée sur sa poitrine. Tout le village l'accompagnait : les femmes chantaient, les enfants dansaient, les vieillards frappaient des mains. On avait tué une chèvre, préparé du vin de palme, et la fête dura jusqu'au milieu de la nuit. Mais le lendemain, la pluie se mit à tomber, une pluie violente, interminable, qui emporta le toit de la case de Meka. Le vieil homme alla chercher refuge chez les pères de la mission, parce qu'il pensait que ses amis blancs l'accueilleraient. On le chassa comme un chien. Alors seulement, trempé, humilié, tremblant de froid et de colère, Meka comprit que la médaille n'était qu'un morceau de métal sans valeur et que l'amitié des Blancs n'était qu'un mensonge. »
\end{exercice}

\begin{corrige}[Exercice : contraction en 40 mots]
\textbf{Idées essentielles} : (a) fête après la médaille ; (b) la pluie détruit sa case ; (c) les missionnaires le chassent ; (d) Meka comprend la vanité de la médaille et le mensonge des Blancs.\par
\textbf{Résumé possible (38 mots)} : « Fêté par son village après la cérémonie, Meka perd sa case sous la pluie. Chassé par les missionnaires chez qui il cherche refuge, il comprend, humilié, que la médaille est sans valeur et l'amitié des Blancs un mensonge. » (38 mots : dans la marge de 36 à 44 mots.)
\end{corrige}

\begin{savaistu}[La contraction, une compétence professionnelle]
La contraction de texte n'est pas qu'un exercice scolaire : elle entraîne une compétence clé de la vie professionnelle, la rédaction de \cle{notes de synthèse}. Dans la fonction publique camerounaise, un administrateur doit savoir résumer en une page un dossier de cent pages pour permettre à son supérieur de décider rapidement. Les concours administratifs (ENAM, IRIC, concours des douanes et des impôts) comportent d'ailleurs souvent une épreuve de synthèse de documents, directement héritée de la contraction apprise au lycée.
\end{savaistu}

\begin{aretenir}[L'essentiel de la section]
\begin{itemize}
\item La contraction réduit un texte à une fraction imposée (souvent 1/4), avec fidélité, neutralité et reformulation.
\item Trois techniques de réduction : \textbf{suppression} (exemples, répétitions, détails, figures, énumérations), \textbf{généralisation} (terme spécifique $\rightarrow$ terme générique), \textbf{condensation} (subordonnée $\rightarrow$ nom, fusion de phrases).
\item La reformulation réexprime les idées avec d'autres mots, d'autres constructions, une autre voix, sans changer le sens.
\item Interdits : copier des phrases, ajouter son avis, sortir de la marge de $\pm$ 10 \% du nombre de mots imposé.
\item Méthode en 4 étapes : lire et comprendre ; dégager le plan ; réduire ; reformuler et compter.
\end{itemize}
\end{aretenir}

\coursec{Le plan du texte et la production du résumé}

Après avoir maîtrisé les techniques de réduction lexicale et syntaxique — suppression des détails, généralisation des énumérations, fusion des phrases, nomination — nous disposons du « matériau » allégé. Il reste l'étape architecturale : organiser ce matériau selon la logique même du texte source. C'est la fonction du \cle{plan}. Sans plan, la contraction risque de devenir une suite de phrases disjointes, perdant la progression argumentative ou narrative qui fait le sens de l'œuvre. Cette section vous guide de la construction du squelette à la rédaction finale, en passant par l'atelier d'amélioration et l'auto-évaluation, compétences décisives pour l'épreuve du Probatoire et du Baccalauréat technique.

\courssub{Le rôle du plan dans la contraction : garant de la fidélité structurelle}

Le plan n'est pas un exercice scolaire accessoire ; c'est la \cle{colonne vertébrale} du résumé. Un texte original possède une architecture (mouvements, transitions, articulations) que le résumé doit reproduire à l'échelle réduite. Respecter le plan, c'est garantir deux fidélités simultanées : la \cle{fidélité du sens} (les idées principales dans leur ordre logique) et la \cle{fidélité des proportions} (l'espace accordé à chaque partie reflète son importance dans le texte source).

\begin{propriete}[Loi de l'homothétie structurelle]
Un bon résumé reproduit la structure du texte original dans les mêmes proportions. Si le texte consacre 50\,\% de sa longueur à la thèse, 30\,\% à l'antithèse et 20\,\% à la synthèse, le résumé doit respecter approximativement cette répartition. Le plan est l'outil de contrôle de cette homothétie.
\end{propriete}

\begin{definition}[Plan de contraction]
C'est la représentation schématique, linéaire et hiérarchisée des grandes parties (mouvements) du texte, où chaque partie est réduite à son \cle{idée directrice} formulée en une phrase complète. Il sert de feuille de route à la rédaction.
\end{definition}

\begin{attention}[Piège : confondre plan du texte et plan de dissertation]
Le plan de contraction \textbf{suit} l'auteur (plan « chronologique » ou « logique de l'auteur »). Il ne s'agit pas d'imposer un plan dialectique (thèse/antithèse/synthèse) à un texte qui n'en a pas (ex. : un récit, une description, une explication linéaire). Le plan épouse la forme du texte : narratif $\rightarrow$ plan narratif ; argumentatif $\rightarrow$ plan dialectique ou analytique ; explicatif $\rightarrow$ plan thématique.
\end{attention}

\courssub{Construction du plan : du texte aux grandes parties}

La construction du plan est une opération de \cle{segmentation} suivie de \cle{formulation}. Elle demande une relecture active, crayon en main.

\begin{methode}[Construire le plan d'un texte en 4 étapes]
\begin{enumerate}
    \item \textbf{Segmenter (découper en mouvements) :} Repérez les indices de rupture : connecteurs forts (\og or \fg{}, \og cependant \fg{}, \og c'est pourquoi \fg{}), changements de temps, de lieu ou de personnage, sauts de paragraphe, reprises thématiques. Numérotez les mouvements (I, II, III...).
    \item \textbf{Identifier l'idée directrice de chaque mouvement :} Pour chaque bloc, demandez-vous : « Quelle est la \emph{seule} chose que l'auteur veut faire comprendre ici ? ». Éliminez les exemples, les illustrations, les détails.
    \item \textbf{Formuler en phrase complète (verbe opérateur) :} Chaque idée directrice devient un item du plan. Utilisez un verbe de liaison intellectuelle : \emph{L'auteur montre que...}, \emph{L'auteur affirme que...}, \emph{L'auteur illustre...}, \emph{L'auteur démontre que...}, \emph{L'auteur raconte comment...}, \emph{L'auteur explique pourquoi...}. Évitez les titres nominaux vagues (ex. : \og La médaille \fg{}).
    \item \textbf{Vérifier l'enchaînement :} Lisez les items à la suite. Forment-ils un propos cohérent ? La logique (cause/conséquence, opposition, approfondissement) est-elle visible ?
\end{enumerate}
\end{methode}

\begin{exemplebox}[Plan d'un extrait à visée critique inspiré du \emph{Vieux nègre et la médaille} (Ferdinand Oyono)]
Extrait centré sur la remise de la médaille à Meka (texte de travail rédigé pour l'entraînement, dans l'esprit du roman).
\begin{enumerate}
    \item \textbf{L'auteur montre que} la cérémonie officielle est une mise en scène grotesque où l'administration coloniale s'auto-félicite par des discours vides.
    \item \textbf{L'auteur illustre} la réalité vécue par Meka : l'humiliation physique (le costume inadapté, la chaleur), l'incompréhension du sens de la décoration et la solitude du récipiendaire.
    \item \textbf{L'auteur révèle} la prise de conscience lucide de Meka : la médaille n'est pas un honneur mais une chaîne, symbole de son aliénation et de la duperie coloniale.
\end{enumerate}
\emph{Observation :} le plan suit la chronologie narrative (cérémonie $\rightarrow$ vécu intérieur $\rightarrow$ révélation finale) tout en faisant émerger la dimension argumentative (ironie $\rightarrow$ réalité $\rightarrow$ vérité).
\end{exemplebox}

\begin{experience}[Atelier : segmentation et formulation du plan]
\textbf{Matériel :} un texte argumentatif court (environ 200 mots) sur l'exode rural (fourni par l'enseignant ou extrait du manuel). \textbf{Protocole :}
\begin{enumerate}
    \item \textbf{Individuel (10 min) :} tracez des traits verticaux \og | \fg{} aux ruptures. Écrivez l'idée directrice de chaque bloc en marge.
    \item \textbf{Binôme (10 min) :} comparez vos découpages. Un seul plan doit émerger. Négociez la formulation exacte des items (verbe opérateur + proposition complétive).
    \item \textbf{Mise en commun (5 min) :} affichage des plans au tableau. Discussion : les proportions sont-elles respectées ? Les verbes opérateurs sont-ils variés et précis ?
\end{enumerate}
\textbf{Interprétation :} la difficulté réside souvent dans le choix du verbe : \og critique \fg{} vs \og dénonce \fg{} vs \og montre les limites de \fg{}. Le verbe engage l'interprétation de l'intention de l'auteur.
\end{experience}

\courssub{Les types de plans et leur adaptation au résumé}

Selon la nature du texte (argumentatif, narratif, explicatif, descriptif), le plan emprunte des structures types. Les identifier aide à vérifier la cohérence de son propre découpage.

\begin{definition}[Typologie des plans en contraction]
\begin{itemize}
    \item \textbf{Plan thématique (ou par aspects) :} texte explicatif ou descriptif. L'auteur traite un sujet sous plusieurs angles. Ex. : I. Les causes de l'exode. II. Les conséquences sur le village. III. Les solutions envisagées.
    \item \textbf{Plan dialectique (thèse / antithèse / synthèse) :} texte argumentatif classique. L'auteur confronte deux points de vue avant de trancher. Ex. : I. L'éloge de la modernité (thèse). II. La critique de la tradition (antithèse). III. La nécessaire synthèse.
    \item \textbf{Plan analytique (problème / causes / solutions) ou (fait / analyse / portée) :} texte réflexif, essai. L'auteur pose un problème, l'analyse, en tire des conséquences.
    \item \textbf{Plan narratif / chronologique :} récit, nouvelle, témoignage. L'auteur suit le temps : I. Situation initiale. II. Élément déclencheur et péripéties. III. Dénouement et situation finale.
\end{itemize}
\end{definition}

\begin{attention}[Piège : le plan « plaqué »]
Ne forcez pas un texte narratif dans un moule dialectique. Si \emph{Le Vieux nègre et la médaille} suit une chronologie ironique, votre plan doit refléter cette chronologie (I. La préparation. II. La cérémonie. III. Le retour et la révélation), même si l'intention globale est critique. Le respect de la \cle{forme} du texte est une condition de la fidélité.
\end{attention}

\courssub{Rédaction du résumé : de l'introduction à la conclusion}

Une fois le plan validé, la rédaction suit un canevas immuable. Le style doit être \cle{impersonnel}, \cle{concis} et \cle{fluide}.

\begin{methode}[Rédiger le résumé en 3 temps]
\begin{enumerate}
    \item \textbf{L'introduction (1 à 2 phrases) :} elle situe le texte sans le paraphraser.
        \begin{itemize}
            \item Auteur + titre (souligné ou entre guillemets) + nature du texte (extrait de roman, article, discours...).
            \item Thème général + thèse principale (l'idée force du texte tout entier).
        \end{itemize}
        \emph{Modèle :} Dans cet extrait du \emph{Vieux nègre et la médaille} de Ferdinand Oyono, l'auteur dénonce l'absurdité et la cruauté du système colonial à travers la cérémonie ironique de remise d'une médaille à un vieux paysan, Meka.
    \item \textbf{Le développement (cœur du résumé) :} il déroule le plan. Chaque item du plan devient un bloc de phrases.
        \begin{itemize}
            \item Respectez l'ordre du plan.
            \item Utilisez des \cle{connecteurs logiques} pour marquer l'enchaînement : \og dans un premier temps \fg{}, \og ensuite \fg{}, \og cependant \fg{}, \og par conséquent \fg{}, \og finalement \fg{}.
            \item Appliquez les techniques de réduction (fusion, généralisation, nomination) vues précédemment.
            \item \textbf{Style impersonnel obligatoire :} pas de \og je pense que \fg{}, \og nous voyons que \fg{}. Préférez : \og l'auteur montre \fg{}, \og le texte révèle \fg{}, \og il apparaît que \fg{}.
        \end{itemize}
    \item \textbf{La conclusion (1 phrase) :} elle referme le résumé sans ajouter d'idée nouvelle. Elle peut synthétiser la portée du texte ou ouvrir sur sa signification profonde.
        \emph{Modèle :} Ainsi, par le biais de l'ironie, Oyono transforme une décoration officielle en une condamnation sans appel de l'ordre colonial.
\end{enumerate}
\end{methode}

\begin{attention}[Les 4 interdits du résumé (sanctionnés à l'examen)]
\begin{enumerate}
    \item \textbf{Le jugement personnel :} \og c'est triste \fg{}, \og c'est bien fait \fg{}, \og l'auteur a raison \fg{}.
    \item \textbf{La citation longue :} plus de 3 ou 4 mots consécutifs du texte original (sauf titre, nom propre, terme technique incontournable). Reformulez systématiquement.
    \item \textbf{L'exemple détaillé :} \og il cite le cas de Meka qui met son costume... \fg{} $\rightarrow$ \og il illustre l'humiliation par le costume inadapté \fg{}.
    \item \textbf{La paraphrase de proximité :} copier la structure des phrases en changeant deux mots. C'est du plagiat, pas de la reformulation. Changez la syntaxe (passif/actif, nominalisation/verbalisation).
\end{enumerate}
\end{attention}

\courssub{Production et amélioration : atelier d'écriture intégrée}

C'est ici que tout se joue : passer du plan à la phrase, contrôler le nombre de mots, polir la langue. La méthode du \cle{brouillon codé} est recommandée.

\begin{experience}[Atelier complet : de l'extrait au résumé final (30 min)]
\textbf{Support :} extrait de travail inspiré du \emph{Vieux nègre et la médaille} (environ 200 mots — scène de la remise de la médaille, discours du commandant, réaction de Meka). \textbf{Consigne :} résumé en \textbf{60 mots $\pm$ 10\,\%} (soit 54 à 66 mots).

\textbf{Étape 1 — Brouillon codé (10 min) :}
\begin{itemize}
    \item Tracez le plan en marge (3 mouvements).
    \item Soulignez dans le texte les segments à garder (idées principales).
    \item Barrez les segments à supprimer (exemples, adjectifs, connecteurs lourds).
    \item Écrivez au-dessus des segments gardés le mot de substitution (généralisation) ou la fusion prévue.
\end{itemize}

\textbf{Étape 2 — Rédaction du premier jet (10 min) :}
Rédigez sur une feuille à part en suivant le plan. Comptez les mots (1 ligne $\approx$ 8 à 10 mots). Visez 6 à 7 lignes.

\textbf{Étape 3 — Relecture croisée et correction collective (10 min) :}
Échangez votre copie avec un voisin. Grille de relecture (voir ci-dessous). Annotation au crayon rouge : \og HS \fg{} (hors sujet / infidèle), \og LM \fg{} (limite de mots dépassée), \og RL \fg{} (reformulation lâche / citation), \og SY \fg{} (syntaxe fautive). Retour à l'auteur $\rightarrow$ version définitive.
\end{experience}

\begin{exoresolu}[Production d'un résumé corrigé]
\textbf{Extrait source (texte de travail, environ 100 mots) :}
\og Le commandant arriva, majestueux dans son uniforme blanc. Il fit un long discours sur la grandeur de la France, la civilisation apportée aux indigènes, le dévouement de Meka. Il épingla la médaille sur la poitrine du vieux qui tremblait, non de fierté, mais de fièvre. Meka ne comprenait pas les grands mots. Il sentait seulement le poids du métal froid et la brûlure du soleil sur son crâne nu. Quand le commandant eut fini, Meka rentra chez lui et posa la médaille sur sa natte. Elle brillait. Mais Meka vit soudain l'éclat du métal comme le reflet d'une lame : il comprit qu'on venait de lui attacher une chaîne au cou. Il cracha par terre. \fg{}

\textbf{Plan établi :}
\begin{enumerate}
    \item L'auteur montre que la cérémonie officielle glorifie l'œuvre coloniale par un discours pompeux.
    \item L'auteur oppose le ressenti physique et l'incompréhension de Meka à la solennité factice.
    \item L'auteur révèle la lucidité finale de Meka : la médaille est le symbole de son asservissement.
\end{enumerate}

\textbf{Premier jet (72 mots — trop long) :}
Dans cet extrait du \emph{Vieux nègre et la médaille}, Ferdinand Oyono dénonce l'hypocrisie coloniale. Lors d'une cérémonie solennelle, le commandant vante la civilisation française et décore Meka. Mais le vieux ne ressent que fièvre, chaleur et poids du métal, ne comprenant pas les discours. De retour chez lui, il contemple la médaille qui brille. Soudain, il perçoit cet éclat comme une lame : la décoration n'est pas un honneur mais une chaîne symbolisant son aliénation. Il exprime son rejet par un crachat.

\textbf{Version finale (61 mots — conforme) :}
Dans \emph{Le Vieux nègre et la médaille}, Oyono dénonce l'hypocrisie coloniale. Le commandant vante la « civilisation » en décorant Meka, mais le vieux ne perçoit que souffrance physique et incompréhension. Seul, face à la médaille brillante, il réalise soudain que ce métal n'honore pas : il enchaîne. L'éclat devient lame, la décoration symbole d'aliénation. Meka crache, unique réponse à l'outrage.

\textbf{Analyse des gains :} fusion des deux premières phrases (cérémonie + discours), nomination (\og souffrance physique \fg{} pour fièvre/chaleur/poids), élagage des détails (natte, soleil), reformulation syntaxique (\og il exprime son rejet par un crachat \fg{} $\rightarrow$ \og Meka crache, unique réponse \fg{}).
\end{exoresolu}

\courssub{Grille d'auto-évaluation et remédiation : vers l'autonomie}

L'auto-évaluation n'est pas une formalité : c'est l'outil qui transforme l'élève en correcteur de lui-même, condition de la réussite en examen terminal.

\begin{popfigure}
\centering
\begin{tabularx}{\linewidth}{|>{\bfseries}l|X|X|X|}\hline
\rowcolor{popblueL}
\textbf{Critère} & \textbf{Indicateur de réussite} & \textbf{Points de vigilance} & \textbf{Action corrective} \\ \hline
\textbf{1. Fidélité au sens} & Les idées principales sont toutes présentes ; l'ordre est respecté ; aucune idée inventée. & Oubli d'un mouvement ; inversion chronologique ; contresens sur l'ironie. & Re-reporter le plan sur le texte source ; vérifier la correspondance 1 pour 1. \\ \hline
\textbf{2. Économie / limite} & Nombre de mots dans la fourchette $\pm$ 10\,\%. & Trop long (détails gardés) ou trop court (idée majeure coupée). & Compter les mots ; appliquer suppression, généralisation, fusion sur les phrases surnuméraires. \\ \hline
\textbf{3. Reformulation} & Aucune citation longue ; syntaxe changée ; vocabulaire propre (synonymes, hyperonymes). & Paraphrase de proximité ; « copier-coller » déguisé. & Réécrire la phrase à l'envers (fin $\rightarrow$ début) ; nominaliser les verbes, verbaliser les noms. \\ \hline
\textbf{4. Structure / cohérence} & Introduction (auteur/titre/thèse) ; développement (plan + connecteurs) ; conclusion ; style impersonnel. & Introduction absente ; « je » présent ; connecteurs manquants ou erronés. & Vérifier le canevas I/D/C ; remplacer « je pense » par « le texte montre » ; insérer des connecteurs logiques. \\ \hline
\textbf{5. Correction de la langue} & Orthographe (accords, homophones) ; syntaxe (ponctuation, phrases complètes). & Accords du participe passé ; confusion a/à, ou/où, et/est ; phrases sans verbe. & Relire à l'envers (dernière phrase $\rightarrow$ première) pour focaliser sur la forme ; dictionnaire. \\ \hline
\end{tabularx}
\legende{Grille d'auto-évaluation de la contraction de texte (5 critères, barème type sur 20 points).}
\end{popfigure}

\begin{aretenir}[Check-list « Prêt pour l'examen »]
Avant de rendre votre copie, cochez mentalement :
\begin{itemize}
    \item J'ai lu le texte 2 fois (globale + détaillée).
    \item J'ai construit un plan à 3 ou 4 items (verbe opérateur + idée directrice).
    \item Mon introduction cite l'auteur, le titre, le thème et la thèse.
    \item Mon développement suit le plan, utilise des connecteurs, au style impersonnel.
    \item Ma conclusion synthétise sans rien ajouter.
    \item J'ai compté mes mots : je suis dans la fourchette $\pm$ 10\,\%.
    \item J'ai relu à l'envers pour traquer les fautes de langue.
\end{itemize}
\end{aretenir}

\begin{savaistu}[Le contrat du Bac technique : fidélité, économie, langue]
Au Probatoire et au Baccalauréat technique, l'épreuve de contraction (souvent 200 à 250 mots à réduire à 60-80 mots) est notée sur \textbf{20 points}, répartis classiquement ainsi :
\begin{itemize}
    \item \textbf{Fidélité au sens (8 à 10 pts) :} respect des idées, de l'ordre, des nuances (ironie, paradoxe), absence de contresens. C'est le critère roi.
    \item \textbf{Économie / respect de la contrainte (4 à 6 pts) :} fourchette de mots respectée. Hors fourchette = pénalité lourde, souvent note plafonnée.
    \item \textbf{Qualité de la langue / rédaction (6 à 8 pts) :} orthographe, syntaxe, ponctuation, connecteurs, style impersonnel, fluidité. Une copie sans faute mais infidèle perd l'essentiel des points de fidélité.
\end{itemize}
\emph{Astuce :} un résumé de 50 mots au lieu de 60 demandés (hors tolérance) perd des points sur l'économie ; un résumé de 70 mots également. \textbf{Comptez, recomptez, ajustez.}
\end{savaistu}

\courssub{Bilan et remédiation ciblée}

\begin{exercice}[Entraînement autonome — Semaine 1]
Choisissez un article d'opinion (presse ou blog institutionnel) d'environ 200 mots.
\begin{enumerate}
    \item Construisez le plan (3 mouvements).
    \item Rédigez le résumé en 60 mots $\pm$ 10\,\%.
    \item Appliquez la grille d'auto-évaluation. Notez-vous sur 20.
    \item Réécrivez une version corrigée.
\end{enumerate}
\end{exercice}

\begin{corrige}[Exercice n°1 — Pistes de remédiation types]
\begin{itemize}
    \item \textbf{Si fidélité $<$ 6/10 :} revoyez la segmentation : vous coupez dans les idées. Travaillez sur des textes plus courts (100 mots) pour identifier UNE idée par paragraphe.
    \item \textbf{Si économie $<$ 4/6 :} entraînez-vous à la \og nomination \fg{} : transformez plusieurs phrases en un nom + un verbe. Ex. : \og il est arrivé, il a parlé, il est parti \fg{} $\rightarrow$ \og son intervention \fg{}.
    \item \textbf{Si langue $<$ 5/8 :} fiches de révision ciblées : accords du participe passé (avec avoir/être), homophones grammaticaux (10 paires types), connecteurs logiques (classement par relation : cause, conséquence, opposition, concession, chronologie).
    \item \textbf{Si structure faible :} entraînez-vous à écrire des introductions types (modèles à mémoriser) et des conclusions « miroir » (reprendre le verbe de la thèse : \og dénonce \fg{} $\rightarrow$ \og Ainsi, il dénonce... \fg{}).
\end{itemize}
\end{corrige}

\begin{popfigure}
\centering
\begin{tikzpicture}[
    node distance=1.2cm and 0.6cm,
    box/.style={draw=popdark, thick, rounded corners, fill=popcream, text width=2.6cm, align=center, font=\small, minimum height=1.8cm},
    arrow/.style={-Stealth, thick, popblue},
    labelnode/.style={font=\tiny\itshape, text=popmuted, align=center}
]
% Nœuds de l'escalier
\node[box, align=center] (lecture){\textbf{1. LECTURE ACTIVE}\\\emph{Globale + détaillée}\\\og Que dit le texte ? \fg{}};
\node[box, right=of lecture, align=center] (plan){\textbf{2. PLAN STRUCTURÉ}\\\emph{3-4 mouvements}\\\og L'auteur montre que... \fg{}};
\node[box, right=of plan, align=center] (reduction){\textbf{3. RÉDUCTION CODÉE}\\\emph{Sur le texte}\\\og Barrer / souligner / remplacer \fg{}};
\node[box, right=of reduction, align=center] (reformulation){\textbf{4. RÉFORMULATION}\\\emph{Brouillon}\\\og Phrases nouvelles + connecteurs \fg{}};
\node[box, right=of reformulation, align=center] (final){\textbf{5. RÉSUMÉ FINAL}\\\emph{Propre + compté}\\\og Intro / dév. / concl. \fg{}};

% Flèches de progression
\draw[arrow] (lecture) -- (plan) node[labelnode, midway, above] {Segmenter};
\draw[arrow] (plan) -- (reduction) node[labelnode, midway, above] {Formuler};
\draw[arrow] (reduction) -- (reformulation) node[labelnode, midway, above] {Transformer};
\draw[arrow] (reformulation) -- (final) node[labelnode, midway, above] {Polir, compter};

% Flèches de retour (remédiation)
\draw[arrow, poporange, dashed] (final.south) to[bend right=20] node[labelnode, below] {Relecture croisée} (reformulation.south);
\draw[arrow, poporange, dashed] (reformulation.south) to[bend right=20] node[labelnode, below] {Ajustement} (reduction.south);
\draw[arrow, poporange, dashed] (reduction.south) to[bend right=20] node[labelnode, below] {Relecture} (plan.south);
\draw[arrow, poporange, dashed] (plan.south) to[bend right=20] node[labelnode, below] {Recompréhension} (lecture.south);

% Titre global
\node[above=0.3cm of reduction, font=\bfseries\large, text=popdark] {Du texte source au résumé final : la méthode en 5 étapes};
\end{tikzpicture}
\legende{Schéma du processus itératif de la contraction : chaque étape valide la précédente ; les flèches orange en pointillés symbolisent les allers-retours nécessaires (remédiation).}
\end{popfigure}

Cette section achève l'unité « Contraction de texte ». Vous possédez maintenant la chaîne complète : \textbf{Lire $\rightarrow$ Planifier $\rightarrow$ Réduire $\rightarrow$ Reformuler $\rightarrow$ Vérifier}. C'est une compétence transversale, utile bien au-delà de l'examen : synthétiser un rapport technique, rédiger une note de synthèse, préparer une soutenance. Entraînez-vous sur des textes variés (techniques, littéraires, journalistiques) pour automatiser le geste expert.

\newpage

\coursec{Activité d'intégration}

\courssub{Objectifs de l'activité}

À l'issue de cet atelier d'intégration, l'élève sera capable de :

\begin{itemize}
\item analyser la \cle{structure argumentative} d'un texte de presse authentique : repérer la thèse, les arguments et les exemples ;
\item établir le \cle{plan} du texte en distinguant les idées essentielles des idées secondaires ;
\item appliquer les trois \cle{techniques de réduction} : suppression, généralisation et condensation ;
\item \cle{reformuler} les idées essentielles avec ses propres mots, sans copier-coller ni déformer la pensée de l'auteur ;
\item repérer les \cle{énoncés performatifs} et les marques d'\cle{ironie} qui engagent la prise de position de l'auteur ;
\item respecter une contrainte quantitative : produire un résumé d'environ \cle{125 mots}, soit le quart du texte de départ ;
\item évaluer la production d'un pair à l'aide d'une grille et améliorer sa propre production.
\end{itemize}

Cette activité mobilise l'ensemble des savoir-faire de l'unité : appliquer les notions de la leçon à une situation concrète de la vie courante (le débat sur les sachets plastiques dans les marchés camerounais) et rédiger puis justifier une démarche, une réponse ou une conclusion.

\courssub{Situation}

Le journal de ta classe prépare un dossier sur l'environnement urbain au Cameroun. Le rédacteur en chef te confie la tâche suivante : résumer en \textbf{125 mots} (marge tolérée : 112 à 137 mots) l'éditorial ci-dessous, paru dans un quotidien camerounais, afin qu'il puisse être publié dans la rubrique « Vu dans la presse ». Le résumé doit être fidèle à la pensée de l'auteur, rédigé avec tes propres mots, et ne comporter ni exemple, ni citation, ni commentaire personnel.

\begin{exemplebox}[Texte à résumer : éditorial (500 mots)]
\textbf{Faut-il interdire les sachets plastiques dans les marchés de Yaoundé ?}

Depuis des années, nos marchés étouffent sous les sachets plastiques. À Mokolo, à Mvog-Mbi, au marché central, le spectacle est le même : des montagnes de sachets noirs jonchent les allées, bouchent les caniveaux et, à la première pluie, transforment les rues en marécages. Faut-il, pour en finir avec ce fléau, interdire purement et simplement les sachets plastiques dans les marchés de Yaoundé ? Nous répondons : oui, sans hésiter.

D'abord, l'interdiction est une nécessité sanitaire. Les sachets abandonnés stagnent dans les rigoles et deviennent des nids à moustiques. Or le paludisme reste, rappelons-le, la première cause de consultation dans nos hôpitaux. À l'hôpital central de Yaoundé, les services de pédiatrie débordent chaque saison des pluies. Qui oserait prétendre que ces petits sachets noirs, si pratiques en apparence, ne coûtent rien à la collectivité ?

Ensuite, l'interdiction est une nécessité économique. La ville de Yaoundé dépense chaque année des sommes considérables pour curer ses caniveaux obstrués par les plastiques. Cet argent, gaspillé à réparer les dégâts, pourrait servir à construire des salles de classe ou des dispensaires. Pensons-y : chaque sachet offert « gratuitement » au client est en réalité payé, demain, par le contribuable.

Certains objectent que l'interdiction pénaliserait les petits commerçants et les ménages modestes, habitués à ce conditionnement bon marché. L'argument est recevable, mais il ne résiste pas à l'examen. Nos grands-mères se rendaient au marché avec des paniers, des cabas et des feuilles de bananier ; nul n'est mort de cette pratique ! Le retour aux emballages biodégradables créerait même des emplois : vannerie, fabrication de sacs en papier, couture de cabas.

Bien sûr, une loi ne suffit pas. L'interdiction devra s'accompagner d'une campagne de sensibilisation dans les marchés et les écoles, et d'un contrôle réel des importations. Mais à force d'attendre le consensus parfait, nous ne ferons rien. Nous exigeons donc des autorités municipales qu'elles prennent leurs responsabilités : qu'elles interdisent les sachets plastiques, et vite. La santé de nos enfants et l'avenir de notre ville valent bien quelques sacrifices d'habitude.
\end{exemplebox}

\courssub{Consignes}

Travail en groupes de quatre, en cinq étapes, puis mise en commun.

\begin{enumerate}
\item \textbf{Analyse de la structure.} Relevez la question posée par l'éditorial et la thèse défendue par l'auteur. Dressez la liste des arguments en les numérotant, et associez à chacun l'exemple qui l'illustre. Comment l'auteur répond-il à l'objection des adversaires de l'interdiction ?
\item \textbf{Plan du texte.} Regroupez les arguments en grandes parties et rédigez le plan du texte en trois ou quatre étapes (une phrase par étape).
\item \textbf{Repérage des passages supprimables.} Soulignez au crayon les exemples, les répétitions, les énumérations et les formules d'insistance. Relevez aussi deux énoncés performatifs (ordre, exigence, promesse...) et dites ce qu'ils révèlent de la prise de position de l'auteur. Un procédé ironique apparaît dans le texte : lequel ?
\item \textbf{Production de la contraction.} Rédigez collectivement un résumé de 125 mots environ (112 à 137 mots) en appliquant les trois techniques : suppression, généralisation, condensation. Interdictions : copier une phrase du texte, donner un avis personnel, employer « je ».
\item \textbf{Évaluation croisée.} Échangez votre production avec un autre groupe et évaluez-la selon la grille ci-dessous. Rédigez ensuite votre version améliorée en tenant compte des remarques reçues.
\end{enumerate}

\begin{center}
\begin{tabularx}{\linewidth}{|l|X|c|}
\hline
\rowcolor{popblueL} \textbf{Critère} & \textbf{Attendu} & \textbf{Note /5} \\
\hline
Fidélité & La thèse et tous les arguments essentiels sont présents, sans déformation ni avis personnel. & \\
\hline
Limite de mots & Le résumé compte entre 112 et 137 mots. & \\
\hline
Reformulation & Aucune phrase copiée ; techniques de réduction visibles (suppression, généralisation, condensation). & \\
\hline
Langue & Orthographe, grammaire et ponctuation correctes ; style clair et impersonnel. & \\
\hline
\end{tabularx}
\end{center}

\courssub{Ressources à mobiliser}

\begin{aretenir}[Boîte à outils de la contraction]
\begin{itemize}
\item \textbf{Structure argumentative} : thèse (la position défendue), arguments (les raisons), exemples (les illustrations, supprimables), objection et réfutation.
\item \textbf{Techniques de réduction} : \emph{suppression} (exemples, répétitions, détails, énumérations) ; \emph{généralisation} (remplacer une liste par un terme générique : « Mokolo, Mvog-Mbi, marché central » $\rightarrow$ « les marchés de la ville ») ; \emph{condensation} (fondre plusieurs phrases en une seule).
\item \textbf{Reformulation} : changer les mots et la construction, jamais le sens ; style impersonnel ; pas de citation.
\item \textbf{Énoncés performatifs} : énoncés qui accomplissent une action (exiger, ordonner, promettre) ; ils signalent l'engagement de l'auteur, par opposition aux énoncés \emph{constatifs}, qui se contentent de décrire un fait.
\item \textbf{Ironie} : dire le contraire de ce qu'on pense pour faire réfléchir ; à signaler dans l'analyse, à neutraliser dans le résumé.
\end{itemize}
\end{aretenir}

\begin{methode}[Les quatre temps de la contraction]
\begin{enumerate}
\item \textbf{Comprendre} : lire le texte deux fois, repérer la thèse et numéroter les arguments ; distinguer l'essentiel (thèse, arguments) de l'accessoire (exemples, répétitions, détails).
\item \textbf{Organiser} : rédiger le plan du texte en trois ou quatre étapes ; chaque étape deviendra une ou deux phrases du résumé.
\item \textbf{Réduire} : supprimer l'accessoire, généraliser les énumérations, condenser les développements ; viser le quart du texte (ici 125 mots pour 500).
\item \textbf{Reformuler et vérifier} : rédiger avec ses propres mots au style impersonnel, puis compter les mots et contrôler la fidélité à la pensée de l'auteur.
\end{enumerate}
\end{methode}

\begin{attention}[Pièges fréquents à l'épreuve]
\begin{itemize}
\item Recopier des phrases entières du texte : c'est un collage, pas une contraction ; toute phrase reprise est sanctionnée.
\item Conserver un exemple « parce qu'il est frappant » : les exemples sont toujours supprimables.
\item Ajouter son avis (« je pense que l'auteur a raison ») : le résumé est impersonnel et neutre.
\item Oublier la réfutation de l'objection : elle fait partie de la logique argumentative du texte.
\item Négliger le comptage des mots : hors de la fourchette imposée, des points sont perdus.
\item Traduire l'ironie au premier degré : « si pratiques » ne signifie pas que l'auteur trouve les sachets pratiques.
\end{itemize}
\end{attention}

\courssub{Démarche et corrigé commenté}

\begin{exoresolu}[Corrigé commenté de l'atelier]
Énoncé : résumer l'éditorial en 125 mots environ, fidèlement et avec ses propres mots, après analyse complète.
\tcblower
\textbf{Étape 1 — Thèse et structure.} La question : faut-il interdire les sachets plastiques dans les marchés de Yaoundé ? La thèse de l'auteur : oui, il faut les interdire. Les arguments : (1) nécessité sanitaire (nids à moustiques, paludisme) ; (2) nécessité économique (coût du curage des caniveaux payé par le contribuable) ; (3) réfutation de l'objection (les emballages traditionnels existaient et le retour au biodégradable créerait des emplois) ; (4) condition de réussite (sensibilisation et contrôle) et appel à l'action.

\textbf{Étape 2 — Plan.} I. Constat : les marchés étouffent sous les sachets. II. L'interdiction s'impose pour des raisons sanitaires et économiques. III. L'objection économique des adversaires est réfutée. IV. Conclusion : il faut agir vite, avec sensibilisation et contrôle.

\textbf{Étape 3 — Supprimables.} Exemples : noms des marchés, hôpital central, « nos grands-mères ». Répétitions : « nécessité » reprise deux fois. Performatifs : « Nous exigeons... », « qu'elles interdisent... » : ils montrent que l'auteur ne se contente pas de constater, il agit sur le lecteur et sur les autorités. Ironie : « si pratiques en apparence » et « nul n'est mort de cette pratique ! », qui tournent en dérision les arguments adverses ; dans le résumé, on les rend de façon neutre.

\textbf{Étape 4 — Contraction (modèle, 127 mots).}

\emph{L'auteur constate que les marchés de Yaoundé sont envahis par les sachets plastiques, qui bouchent les caniveaux, et défend leur interdiction. Cette mesure s'impose d'abord pour des raisons sanitaires : les sachets abandonnés favorisent la prolifération des moustiques, donc le paludisme. Elle s'impose aussi économiquement, car le curage des caniveaux coûte cher à la collectivité, au détriment d'investissements utiles. À l'objection selon laquelle l'interdiction pénaliserait les petits commerçants, l'auteur répond que les emballages traditionnels biodégradables peuvent la remplacer et même créer des emplois. Il conclut en appelant les autorités municipales à agir rapidement, en accompagnant l'interdiction de sensibilisation et de contrôle.}

\textbf{Vérification.} Thèse et quatre idées essentielles présentes ; aucun exemple, aucune phrase copiée ; ton neutre ; 127 mots, dans la fourchette autorisée (112 à 137 mots).
\end{exoresolu}

\begin{exercice}[Pour s'entraîner]
Reprends le résumé produit par ton groupe et applique-lui la grille d'évaluation. Puis rédige individuellement une seconde version en veillant à : (1) généraliser au moins une énumération restante ; (2) condenser deux phrases en une seule ; (3) vérifier le comptage final des mots. Compare ta version à celle d'un camarade et justifie tes choix de reformulation.
\end{exercice}

\courssub{Bilan de l'activité}

L'atelier a montré que la contraction n'est pas un découpage mécanique mais une démarche en quatre temps inséparables : \textbf{comprendre} (structure argumentative), \textbf{organiser} (plan), \textbf{réduire} (suppression, généralisation, condensation), \textbf{reformuler} (style impersonnel et fidèle). La mise en commun a dégagé les erreurs les plus fréquentes à retravailler en remédiation : recopier des phrases du texte au lieu de reformuler, conserver un exemple « parce qu'il est frappant », ajouter un avis personnel, oublier la réfutation de l'objection — pourtant essentielle à la logique du texte — et négliger le comptage des mots. L'évaluation croisée a enfin appris aux élèves à lire une production avec un regard de correcteur, compétence directement utile à l'épreuve de contraction du Probatoire et du Baccalauréat technique.

\newpage

\coursec{Méthodes et exercices résolus}

\coursec{Contraction de texte : reformulation des idées d'un texte à résumer}

\courssub{Fiche méthode 1 : Appliquer la technique de réduction à une situation concrète}

\begin{methode}[Réduire un texte sans perdre l'essentiel]
La \cle{contraction de texte} consiste à reproduire les idées essentielles d'un texte dans un volume imposé (généralement le quart ou le cinquième du texte d'origine), sans ajouter d'idée personnelle. Démarche en étapes :
\begin{enumerate}
\item \textbf{Lire deux fois} le texte : une première lecture globale pour saisir le thème, une seconde lecture analytique en soulignant les idées principales.
\item \textbf{Compter les mots} du texte et calculer la longueur attendue (par exemple : texte de 400 mots, résumé demandé en 100 mots).
\item \textbf{Découper le texte en unités de sens} (paragraphes ou groupes de phrases) et formuler l'idée essentielle de chaque unité en une phrase.
\item \textbf{Éliminer} les exemples, les répétitions, les citations, les descriptions pittoresques, les digressions.
\item \textbf{Reformuler} avec ses propres mots : remplacer les longues périodes par des propositions courtes, les énumérations par un terme générique, les dialogues par un discours indirect.
\item \textbf{Rédiger} le résumé en un texte continu, cohérent, avec des connecteurs logiques (d'abord, ensuite, enfin ; cependant, donc).
\item \textbf{Vérifier} : compter les mots du résumé (marge de $\pm 10\,\%$ tolérée), contrôler qu'aucune idée personnelle ni aucun jugement n'a été ajouté.
\end{enumerate}
\textbf{Réflexes à retenir} : ne jamais recopier des phrases entières du texte ; ne jamais écrire « l'auteur dit que » à chaque phrase ; respecter l'ordre des idées du texte.
\end{methode}

\begin{exoresolu}[Réduire un paragraphe argumentatif]
\textbf{Énoncé (type Bac)} : Réduisez le paragraphe suivant (environ 120 mots) en 30 mots environ, en reformulant l'idée essentielle.

« L'éducation est la clé du développement d'un pays. En effet, un peuple instruit maîtrise les techniques agricoles modernes, par exemple l'irrigation goutte-à-goutte ou la sélection des semences ; il comprend les enjeux sanitaires, comme la vaccination ou l'hygiène de l'eau ; il participe activement à la vie démocratique, en votant en connaissance de cause, en contrôlant l'action des élus, en réclamant ses droits. Sans éducation, aucun progrès durable n'est possible, comme le montrent les pays où l'analphabétisme maintient les populations dans la pauvreté et la dépendance. »
\tcblower
\textbf{Solution détaillée}

\textbf{Étape 1 — Analyse.} Le paragraphe contient une thèse (« l'éducation est la clé du développement »), trois arguments illustrés par des exemples (agriculture, santé, démocratie) et une conclusion par l'absurde (« sans éducation, aucun progrès »).

\textbf{Étape 2 — Élimination.} Les exemples précis (irrigation goutte-à-goutte, vaccination, vote...) sont des illustrations : on les supprime et on garde les domaines généraux (économie, santé, vie politique).

\textbf{Étape 3 — Reformulation.} On condense la thèse et les trois arguments en une ou deux phrases.

\textbf{Résumé proposé (32 mots)} : « L'éducation est indispensable au développement : elle permet les progrès économiques, sanitaires et démocratiques, tandis que l'analphabétisme enferme les peuples dans la pauvreté et la dépendance. »

\textbf{Vérification} : 32 mots (consigne respectée), aucune phrase recopiée, aucune idée ajoutée, ordre du texte conservé. \textbf{Conclusion} : le résumé restitue fidèlement la thèse et les arguments essentiels du paragraphe.
\end{exoresolu}

\courssub{Fiche méthode 2 : Dégager la structure argumentative d'un texte}

\begin{methode}[Identifier thèse, arguments et exemples]
Un texte argumentatif se construit selon un schéma qu'il faut savoir repérer avant toute contraction :
\begin{enumerate}
\item \textbf{Repérer la thèse} : l'idée générale que l'auteur défend (souvent au début ou à la fin). Formule à retenir : « L'auteur soutient que... »
\item \textbf{Repérer les arguments} : les raisons qui justifient la thèse. Indices : connecteurs (car, en effet, d'abord, ensuite, de plus, enfin).
\item \textbf{Repérer les exemples} : faits précis, anecdotes, chiffres qui illustrent les arguments. Indices : « par exemple », « ainsi », « c'est le cas de ».
\item \textbf{Construire le schéma} : thèse $\rightarrow$ argument 1 + exemple, argument 2 + exemple, etc. $\rightarrow$ conclusion.
\item \textbf{Justifier la démarche} : préciser si le texte procède par déduction (du général au particulier), par induction (des faits à l'idée générale) ou par l'absurde.
\end{enumerate}
\textbf{Réflexe} : dans \emph{Le Vieux Nègre et la Médaille} de Ferdinand Oyono, le narrateur expose les faits (récit) qui font éclater l'absurdité de la colonisation : l'argumentation y est indirecte, portée par l'ironie.
\end{methode}

\begin{exoresolu}[Analyser la structure d'un extrait]
\textbf{Énoncé (type Bac)} : Dans \emph{Le Vieux Nègre et la Médaille}, Meka, vieillard africain, reçoit une médaille pour avoir donné ses terres à la mission et ses deux fils à la guerre. L'auteur écrit que Meka attend cette distinction « avec une joie immense ». Dégagez la structure argumentative de ce passage et justifiez la démarche de l'auteur.
\tcblower
\textbf{Solution détaillée}

\textbf{Étape 1 — Thèse implicite.} Oyono ne déclare jamais « la colonisation est injuste » : la thèse est \emph{implicite}. Elle se déduit de la situation : Meka a tout sacrifié (ses terres, ses deux fils morts à la guerre) et ne reçoit en échange qu'une médaille sans valeur.

\textbf{Étape 2 — Arguments.} Le texte accumule les faits qui servent d'arguments : (1) le don des terres aux Pères blancs ; (2) la perte des deux fils ; (3) la misère et la naïveté de Meka ; (4) l'indifférence des colons après la cérémonie.

\textbf{Étape 3 — Exemples.} La cérémonie de remise de la médaille, la nuit de pluie passée dehors, l'arrestation finale de Meka sont des exemples concrets qui illustrent l'ingratitude du système colonial.

\textbf{Étape 4 — Démarche.} L'auteur procède par \cle{induction} et par \emph{l'absurde} : il part de faits concrets (le récit) pour conduire le lecteur à conclure lui-même que la colonisation est une exploitation. La « joie immense » de Meka est ironique : elle souligne son aveuglement et rend la critique plus mordante.

\textbf{Conclusion} : la structure argumentative repose sur une démonstration indirecte où le récit lui-même, par accumulation de faits et ironie, dénonce la colonisation sans discours explicite.
\end{exoresolu}

\courssub{Fiche méthode 3 : Reformuler les idées essentielles}

\begin{methode}[Reformuler sans trahir]
La \cle{reformulation} est le cœur de la contraction : il faut dire autrement la même chose. Techniques à appliquer :
\begin{enumerate}
\item \textbf{Substitution lexicale} : remplacer un mot par un synonyme ou un terme générique (« voitures, motos, camions » $\rightarrow$ « véhicules »).
\item \textbf{Changement de construction} : passer de l'actif au passif, du nom au verbe (« la destruction de la forêt » $\rightarrow$ « on détruit la forêt »), du discours direct au discours indirect.
\item \textbf{Fusion} : condenser plusieurs phrases en une seule grâce à une subordonnée.
\item \textbf{Suppression des modalités} : éliminer les adverbes d'intensité, les adjectifs élogieux, les interjections qui n'ajoutent pas d'idée.
\item \textbf{Contrôle de fidélité} : relire et vérifier que le sens est identique (même idée, même nuance, même point de vue de l'auteur).
\end{enumerate}
\begin{center}
\begin{tabularx}{\linewidth}{|X|X|}
\hline
\rowcolor{popblueL} \textbf{Texte original} & \textbf{Reformulation} \\
\hline
« Je suis absolument ravi, mes chers amis, de cette magnifique médaille ! » & Il se réjouit de sa médaille. \\
\hline
« Les pluies, le soleil brûlant et les insectes détruisent les récoltes. » & Les intempéries et les parasites nuisent à l'agriculture. \\
\hline
\end{tabularx}
\end{center}
\end{methode}

\begin{exoresolu}[Reformuler un passage de dialogue]
\textbf{Énoncé (type Bac)} : Reformulez en style indirect et en deux phrases maximum le passage suivant, tiré d'un dialogue de \emph{Le Vieux Nègre et la Médaille} (adapté) :

« — Meka, dit le commandant, tu as été un homme bon et fidèle. Tu as donné tes terres à la mission, tu as donné tes fils à la France. Demain, à la fête du 14 juillet, tu recevras la médaille de l'amitié franco-africaine !
— Ah ! s'écria Meka, les Blancs sont vraiment bons ! Quel bonheur pour moi, quel honneur pour ma famille ! »
\tcblower
\textbf{Solution détaillée}

\textbf{Étape 1 — Idées essentielles.} (1) Le commandant rappelle les sacrifices de Meka (terres, fils) ; (2) il annonce la remise d'une médaille ; (3) Meka exprime sa joie et sa gratitude naïve envers les Blancs.

\textbf{Étape 2 — Passage au discours indirect.} « dit » $\rightarrow$ « déclare que » ; « tu as donné » $\rightarrow$ « il avait donné » (concordance des temps : plus-que-parfait) ; « tu recevras » $\rightarrow$ « il recevrait » (conditionnel présent pour le futur dans le passé) ; « s'écria » $\rightarrow$ « s'exclama que ».

\textbf{Étape 3 — Suppression des modalités.} Les exclamations (« Ah ! », « Quel bonheur ! ») sont remplacées par une formulation neutre : « il exprima sa joie ».

\textbf{Reformulation proposée} : « Le commandant déclara à Meka que, en récompense de sa fidélité et du sacrifice de ses terres et de ses fils, il recevrait le lendemain une médaille. Meka, transporté de joie, s'exclama que les Blancs étaient bons et que cette distinction honorait toute sa famille. »

\textbf{Conclusion} : la reformulation conserve les trois idées essentielles, respecte la concordance des temps et neutralise le ton exclamatif, comme l'exige un résumé objectif.
\end{exoresolu}

\courssub{Fiche méthode 4 : Construire le plan du texte avant de résumer}

\begin{methode}[Établir le plan d'un texte]
Le \cle{plan du texte} est le squelette sur lequel repose le résumé. Démarche :
\begin{enumerate}
\item \textbf{Numéroter les paragraphes} et résumer chacun par un mot-clé ou une courte phrase au brouillon.
\item \textbf{Regrouper} les paragraphes qui développent la même idée : ce sont les parties du plan.
\item \textbf{Intituler chaque partie} par une phrase complète (et non un simple mot) qui exprime une idée.
\item \textbf{Ordonner} : vérifier que les parties suivent la progression du texte (chronologique, logique, dialectique).
\item \textbf{Utiliser le plan} comme trame du résumé : une partie du plan = une ou deux phrases du résumé, reliées par des connecteurs.
\end{enumerate}
\textbf{Réflexe} : pour un récit comme \emph{Le Vieux Nègre et la Médaille}, le plan est souvent chronologique : situation initiale (l'annonce de la médaille) $\rightarrow$ préparatifs et cérémonie $\rightarrow$ retournement (l'errance nocturne, l'arrestation) $\rightarrow$ prise de conscience finale.
\end{methode}

\begin{exoresolu}[Construire le plan d'un résumé]
\textbf{Énoncé (type Bac)} : On vous demande de résumer en 150 mots un texte de 600 mots relatant le parcours de Meka dans \emph{Le Vieux Nègre et la Médaille}. Établissez le plan du texte, puis rédigez la première phrase de chaque partie du résumé.
\tcblower
\textbf{Solution détaillée}

\textbf{Étape 1 — Découpage en unités de sens.} Le récit se divise en quatre moments : (1) Meka apprend qu'il va recevoir une médaille ; (2) le village prépare la cérémonie ; (3) la cérémonie a lieu et Meka, oublié de tous, erre sous la pluie ; (4) arrêté par la police, Meka comprend la duperie.

\textbf{Étape 2 — Plan en phrases complètes.}
\begin{itemize}
\item \textbf{Partie 1} : Meka, vieillard naïf, se réjouit de la médaille promise par l'administration coloniale.
\item \textbf{Partie 2} : Tout le village partage son enthousiasme et prépare la fête.
\item \textbf{Partie 3} : Après la cérémonie, Meka est abandonné et erre toute la nuit sous la pluie.
\item \textbf{Partie 4} : Arrêté comme un vulgaire délinquant, il prend conscience de l'illusion coloniale.
\end{itemize}

\textbf{Étape 3 — Amorces du résumé.} Chaque intitulé devient la première phrase d'une partie du résumé, reliée par un connecteur : « D'abord... », « Ensuite... », « Mais après la cérémonie... », « Finalement... ».

\textbf{Vérification} : 4 parties pour 150 mots $\rightarrow$ environ 35 à 40 mots par partie, répartition équilibrée.

\textbf{Conclusion} : le plan chronologique en quatre parties garantit un résumé complet, ordonné et fidèle à la progression du roman.
\end{exoresolu}

\courssub{Fiche méthode 5 : Identifier ironie et paradoxe pour mieux résumer}

\begin{methode}[Reconnaître les figures de pensée]
Avant de reformuler, il faut comprendre le sens réel du texte, qui peut être caché par des figures :
\begin{enumerate}
\item \textbf{L'\cle{ironie}} : dire le contraire de ce qu'on pense pour critiquer. Indices : exagération, éloge excessif et invraisemblable, contraste entre les paroles et la situation. Exemple : qualifier de « généreux » un colon qui prend les terres.
\item \textbf{Le \cle{paradoxe}} : affirmation contraire à l'opinion commune mais qui contient une vérité. Exemple : « Meka était riche de sa médaille et pauvre de tout le reste. »
\item \textbf{Procédure d'analyse} : repérer l'écart entre le sens littéral et le sens voulu $\rightarrow$ nommer la figure $\rightarrow$ expliquer l'effet (critique, dérision, mise en garde).
\item \textbf{Dans le résumé} : ne pas reformuler l'ironie au premier degré ! Il faut restituer le sens réel (« l'auteur dénonce » et non « l'auteur félicite »).
\end{enumerate}
\end{methode}

\begin{exoresolu}[Analyser une phrase ironique]
\textbf{Énoncé (type Bac)} : Dans \emph{Le Vieux Nègre et la Médaille}, on lit que la médaille remise à Meka est « la juste récompense de tant de sacrifices ». (1) Identifiez la figure de pensée employée. (2) Expliquez son effet. (3) Reformulez cette idée pour un résumé.
\tcblower
\textbf{Solution détaillée}

\textbf{(1) Identification.} Il s'agit d'une \cle{ironie}. En effet, l'expression « juste récompense » affirme le contraire de la réalité : Meka a perdu ses terres et ses deux fils, et il ne reçoit qu'un bout de métal sans valeur. L'écart entre l'énormité du sacrifice et l'insignifiance de la récompense rend l'éloge invraisemblable : c'est l'indice de l'ironie.

\textbf{(2) Effet.} L'ironie fait apparaître l'injustice de la colonisation avec plus de force qu'une accusation directe : le lecteur mesure lui-même l'écart entre les mots et les faits, ce qui provoque à la fois le rire amer et la révolte. Oyono dénonce ainsi l'hypocrisie du discours colonial qui masque l'exploitation par des cérémonies et des honneurs vides.

\textbf{(3) Reformulation pour un résumé.} Au premier degré, la phrase serait mensongère ; le résumé doit restituer le sens réel : « L'auteur souligne, par l'ironie, l'injustice d'une récompense dérisoire comparée aux sacrifices consentis par Meka. »

\textbf{Conclusion} : reconnaître l'ironie évite l'erreur majeure du résumé : trahir la pensée de l'auteur en prenant ses mots au pied de la lettre.
\end{exoresolu}

\courssub{Fiche méthode 6 : Distinguer énoncés constatifs et performatifs, variétés du français}

\begin{methode}[Analyser les énoncés et le lexique]
\begin{enumerate}
\item \textbf{Énoncé \cle{constatif}} : il décrit un état de fait ; on peut le juger vrai ou faux. Exemple : « La cérémonie a lieu le 14 juillet. »
\item \textbf{Énoncé \cle{performatif}} : il accomplit l'action qu'il énonce (promettre, ordonner, déclarer, baptiser). Exemple : « Je te décore de la médaille de l'amitié franco-africaine » : en le disant, le commandant décore réellement Meka. Test : on peut ajouter « par la présente ».
\item \textbf{Variétés lexicales du français} : repérer le registre (soutenu, courant, familier, populaire) et les régionalismes (français du Cameroun : « le ndolé », « le makossa », « être au village »). Dans un résumé, on remplace les termes trop familiers ou régionaux par des équivalents du français courant.
\item \textbf{Justifier} : toute réponse doit nommer le phénomène, citer l'exemple et expliquer son rôle dans le texte.
\end{enumerate}
\end{methode}

\begin{exoresolu}[Classer des énoncés et adapter le registre]
\textbf{Énoncé (type Bac)} : (1) Classez les énoncés suivants en constatifs ou performatifs, en justifiant : a) « Je déclare la cérémonie ouverte. » b) « La médaille brillait au soleil. » c) « Je promets de récompenser les fidèles. » (2) Réécrivez en français courant pour un résumé : « Meka était vraiment dans la galère après la fête, le pauvre vieux ! »
\tcblower
\textbf{Solution détaillée}

\textbf{(1) Classification.}
\begin{itemize}
\item a) « Je déclare la cérémonie ouverte » : énoncé \cle{performatif}. Le verbe « déclarer » accomplit l'action : en prononçant cette phrase, l'autorité ouvre effectivement la cérémonie. Le test « par la présente » fonctionne.
\item b) « La médaille brillait au soleil » : énoncé \cle{constatif}. Il décrit un fait observable, que l'on peut vérifier : il est vrai ou faux.
\item c) « Je promets de récompenser les fidèles » : énoncé \cle{performatif}. Dire « je promets », c'est promettre : l'énonciation réalise l'acte de promesse.
\end{itemize}

\textbf{(2) Réécriture.} « Être dans la galère » est familier ; « le pauvre vieux » est une marque de jugement personnel et de registre familier, à bannir d'un résumé neutre. Reformulation : « Après la cérémonie, Meka se retrouva dans une situation très difficile. »

\textbf{Conclusion} : distinguer constatif et performatif permet de comprendre la force des paroles officielles dans le roman, et neutraliser le registre garantit l'objectivité exigée par la contraction de texte.
\end{exoresolu}

\begin{savaistu}[Ferdinand Oyono et \emph{Le Vieux Nègre et la Médaille}]
Ferdinand Oyono (1929--2010), écrivain et diplomate camerounais, publie \emph{Le Vieux Nègre et la Médaille} en 1956. Ce roman court, écrit dans un français teinté de structures bantoues, utilise l'ironie et le récit à la troisième personne pour dénoncer l'aliénation coloniale. Le personnage de Meka, vieillard naïf qui croit en la reconnaissance française, incarne le drame des « évolués » trahis par le système qu'ils ont servi. L'œuvre est au programme du Probatoire / Baccalauréat technique comme support privilégié pour l'étude de l'argumentation indirecte et de la contraction de texte littéraire.
\end{savaistu}

\begin{attention}[Les 5 erreurs qui coûtent des points au Bac]
\begin{enumerate}
\item \textbf{Recopier des phrases du texte.} Un résumé qui juxtapose des phrases de l'auteur est sanctionné : il faut reformuler avec ses propres mots. Seuls quelques termes techniques irremplaçables peuvent être repris.
\item \textbf{Ajouter son avis personnel.} « Je trouve que l'auteur a raison » ou « cette histoire est touchante » : toute appréciation personnelle est hors sujet dans une contraction. Le résumé doit être neutre et objectif.
\item \textbf{Prendre l'ironie au premier degré.} Écrire « Oyono félicite l'administration coloniale » trahit complètement le texte. Toujours restituer le sens réel : l'auteur dénonce, critique, ridiculise.
\item \textbf{Négliger le comptage des mots.} Un résumé de 200 mots pour 100 demandés, ou de 40 mots, est pénalisé. Compter les mots, indiquer le nombre à la fin et respecter la marge de $\pm 10\,\%$.
\item \textbf{Oublier les connecteurs logiques.} Une suite de phrases juxtaposées sans liens (« d'abord, ensuite, cependant, enfin ») donne un texte haché et incohérent. Le résumé est un texte rédigé, continu, pas une liste d'idées.
\end{enumerate}
\end{attention}

\newpage

\coursec{Exercices}

\coursec{Leçon 09 : Contraction de texte : reformulation des idées d'un texte à résumer}

\courssub{Je vérifie mes connaissances}

\begin{exercice}[QCM : techniques de contraction]
Parmi les propositions suivantes, cochez celle qui correspond à la technique de réduction décrite.
\begin{enumerate}
    \item Remplacer une énumération détaillée par un terme générique (ex. : « lions, éléphants, girafes, zèbres » $\to$ « la faune sauvage »).
    \item Supprimer les adjectifs, les compléments circonstanciels jugés accessoires ou les répétitions.
    \item Regrouper plusieurs phrases en une seule en conservant l'information essentielle.
    \item Reformuler une idée avec ses propres mots sans en altérer le sens.
\end{enumerate}
\begin{itemize}
    \item[A.] Condensation \quad [B.] Suppression \quad [C.] Généralisation \quad [D.] Reformulation
\end{itemize}
Associez chaque item (1 à 4) à la bonne lettre.
\end{exercice}

\begin{exercice}[Vrai / Faux justifié : ironie et paradoxe]
Déterminez si les affirmations suivantes sont vraies (V) ou fausses (F). Justifiez brièvement votre réponse en vous appuyant sur les définitions du cours.
\begin{enumerate}
    \item L'ironie consiste toujours à dire le contraire de ce que l'on pense, dans un but moqueur.
    \item Le paradoxe est une figure de style qui affirme une chose et son contraire dans la même phrase.
    \item Dans \emph{Le Vieux nègre et la médaille}, l'ironie permet à Oyono de dénoncer l'absurdité du système colonial sans le nommer frontalement.
    \item Une reformulation au premier degré d'un énoncé ironique conserve le sens profond de l'auteur.
\end{enumerate}
\end{exercice}

\begin{exercice}[Définitions : énoncés constatifs et performatifs]
Donnez une définition précise et un exemple original (non vu en cours) pour chacun des concepts suivants :
\begin{enumerate}
    \item Énoncé constatif.
    \item Énoncé performatif (explicite).
    \item Variété diastratique du français (donner un exemple camerounais).
\end{enumerate}
\end{exercice}

\begin{exercice}[Repérage : structure argumentative]
Lisez l'extrait suivant (simulé, inspiré du roman) :
\begin{quote}
« Meka regarda la médaille. Elle brillait au soleil, petite chose ridicule pour un si grand sacrifice. \textbf{Car} il avait donné ses terres, \textbf{donc} il avait perdu son autonomie, \textbf{c'est pourquoi} il se sentait nu devant cet objet métallique. \textbf{Mais} le commandant avait souri, \textbf{alors} tout était scellé. »
\end{quote}
Identifiez dans ce texte : la \cle{thèse}, un \cle{argument}, un \cle{connecteur logique} marquant la cause, un \cle{connecteur} marquant la conséquence et un \cle{connecteur} marquant l'opposition.
\end{exercice}

\courssub{Je m'entraîne}

\begin{exercice}[Repérage et explication de l'ironie]
Voici un extrait du discours du commandant lors de la remise de la médaille (chapitre 6, adapté) :
\begin{quote}
« Mes amis, vous voyez devant vous un homme qui a tout compris. Meka a compris que la vraie richesse, ce n'est pas la terre qui nourrit ses enfants, c'est ce petit bout de métal qui brille sur sa poitrine. Il a eu le courage de se débarrasser de l'inutile pour garder l'essentiel. La France vous remercie, Meka, d'avoir si bien servi ses intérêts. »
\end{quote}
\begin{enumerate}
    \item Relevez \textbf{trois} énoncés ironiques dans ce passage.
    \item Pour chacun, expliquez le \textbf{sens littéral} (ce qui est dit) et le \textbf{sens réel} (ce qui est visé).
    \item Montrez en quoi une reformulation « au premier degré » (qui prendrait les mots pour ce qu'ils disent) ferait courir le risque d'un \textbf{contresens total} dans un résumé.
\end{enumerate}
\end{exercice}

\begin{exercice}[Classement : constatifs vs performatifs]
Voici dix énoncés tirés d'une scène officielle de remise de décorations. Classez-les dans le tableau ci-dessous en justifiant chaque choix par l'analyse du verbe et de la situation d'énonciation.
\begin{enumerate}
    \item « Je déclare Meka chevalier de la Légion d'honneur. »
    \item « Le soleil brille sur la place du village. »
    \item « Je vous prie de bien vouloir vous asseoir. »
    \item « Meka a vendu ses terres l'an dernier. »
    \item « Je vous nomme chef de canton. »
    \item « La médaille est en bronze doré. »
    \item « Je promets de veiller sur mes administrés. »
    \item « Le commandant parle français. »
    \item « Je baptise ce pont "Pont de l'Union". »
    \item « Les anciens combattent pour la France. »
\end{enumerate}
\begin{center}
\begin{tabularx}{\linewidth}{|X|X|X|}\hline
\textbf{Énoncé} & \textbf{Type (Constatif / Performatif)} & \textbf{Justification (verbe, intention, effet)} \\ \hline
1 & & \\ \hline
2 & & \\ \hline
3 & & \\ \hline
4 & & \\ \hline
5 & & \\ \hline
6 & & \\ \hline
7 & & \\ \hline
8 & & \\ \hline
9 & & \\ \hline
10 & & \\ \hline
\end{tabularx}
\end{center}
\end{exercice}

\begin{exercice}[Réduction ciblée : suppression, généralisation, condensation]
Condensez chacune des phrases suivantes en \textbf{8 mots maximum}. Indiquez entre parenthèses la technique principale utilisée : \textbf{Suppression (S)}, \textbf{Généralisation (G)} ou \textbf{Condensation (C)}.
\begin{enumerate}
    \item Le vieux Meka, courbé sous le poids des ans, marcha péniblement vers la case du commandant.
    \item Les villageois, les femmes, les enfants et les vieillards observaient la scène en silence total.
    \item La médaille, petite, brillante, ronde et dorée, fut épinglée sur la poitrine usée du récipiendaire.
    \item Le commandant prononça un long discours plein de mots difficiles que personne ne comprenait vraiment.
    \item Après la cérémonie, Meka rentra chez lui, le cœur lourd, l'esprit troublé par l'ironie du sort.
\end{enumerate}
\end{exercice}

\begin{exercice}[Construction du plan et rédaction du résumé]
Texte support (vie scolaire au Cameroun, $\approx$ 250 mots) :
\begin{quote}
« L'année scolaire 2024-2025 au Cameroun s'ouvre sous le signe de la résilience. Malgré les grèves récurrentes des enseignants réclamant de meilleures conditions de vie et de travail, les élèves ont massivement regagné les bancs. \textbf{Premièrement}, l'effort budgétaire de l'État a permis la construction de nouvelles salles de classe dans les zones rurales, réduisant la surcharge. \textbf{Deuxièmement}, l'introduction du numérique dans certains lycées techniques offre des perspectives nouvelles, bien que l'accès à l'électricité et à internet reste inégal. \textbf{Troisièmement}, la réforme des programmes techniques vise à mieux adapter la formation aux besoins du marché de l'emploi local. \textbf{Cependant}, le manque d'enseignants qualifiés dans les filières industrielles et le coût élevé des fournitures techniques freinent l'égalité des chances. \textbf{En définitive}, si la volonté politique est affichée, la concrétisation sur le terrain exige un suivi rigoureux et un dialogue social apaisé pour garantir une éducation de qualité pour tous. »
\end{quote}
\begin{enumerate}
    \item Dégagez la \textbf{thèse} de ce texte.
    \item Construisez le \textbf{plan détaillé en trois parties} (avec arguments et exemples) qui structure l'argumentation.
    \item Rédigez le \textbf{résumé} de ce texte en \textbf{60 mots $\pm$ 10\%} (soit 54 à 66 mots). Respectez la neutralité et l'ordre des idées.
\end{enumerate}
\end{exercice}

\courssub{Je me prépare au Bac}

\begin{exercice}[Sujet type Bac : Contraction de texte (Extrait \emph{Le Vieux nègre et la médaille})]
\textbf{Consigne :} Résumez le texte suivant en \textbf{150 mots $\pm$ 10\%} (soit \textbf{135 à 165 mots}). Vous respecterez l'ordre des idées, la neutralité du ton et la cohérence syntaxique. Tout dépassement ou insuffisance sera pénalisé.

\textbf{Texte : Réflexion de Meka après son arrestation (Chapitre 8, adapté, $\approx$ 600 mots).}

« La nuit était tombée sur le village. Assis sur son lit de fortune, dans la cellule humide du poste de police, Meka repassait le film de sa vie. La médaille, cette petite chose froide qu'il avait tant désirée, gisait maintenant au fond de sa poche, lourde comme un secret honteux. Il revoyait le sourire du commandant, ce sourire qui ne touchait pas les yeux, et les mots mielleux : "La France reconnaît vos loyaux services." Loyaux services... Il avait donné sa meilleure terre, celle qui nourrissait sa famille, celle où reposaient les placenta de ses enfants. Il avait donné la terre de ses pères pour un bout de métal. Et maintenant ? Maintenant, on l'arrêtait pour "incitation à la rébellion" parce qu'il avait osé demander un reçu pour la vente de sa terre. Un reçu ! Comme si la parole d'un blanc avait besoin de papier pour être sacrée.

Il pensait à Kelara, sa femme, à ses yeux secs qui ne pleuraient plus, à ses fils partis en brousse pour ne pas porter l'uniforme. Il pensait à l'Église, à l'école, à tous ces pièges tendus par l'homme blanc pour lui voler son âme sous prétexte de la sauver. Le catéchiste lui avait dit : "Rends à César ce qui est à César." Mais César, aujourd'hui, avait pris la terre, la médaille, la liberté et voulait encore la dignité.

Meka sentit une colère sourde, ancienne, monter en lui. Ce n'était pas la peur de la prison, non. C'était la nausée de sa propre complicité. Il avait cru jouer le jeu du blanc, il s'était fait jouer. La médaille n'était pas une récompense, c'était un collier de chien. Et lui, Meka, le "vieux nègre", avait aboyé de joie en le recevant. Soudain, il comprit que la vraie médaille, celle qui ne s'oxyde pas, c'était la terre qu'il avait perdue. Et qu'il ne la reverrait jamais. Le geôlier claqua la porte. Une nouvelle journée commençait. »

\begin{enumerate}
    \item Effectuez le \textbf{dépouillement des idées essentielles} (liste numérotée).
    \item Proposez un \textbf{plan en trois parties} pour organiser le résumé.
    \item Rédigez le \textbf{résumé final} (comptez vos mots et indiquez le total en fin de copie).
\end{enumerate}
\end{exercice}

\begin{exercice}[Sujet type Bac : Reformulation et auto-évaluation]
\textbf{Consigne :} Réécrivez le paragraphe ci-dessous avec \textbf{vos propres mots} (reformulation), sans en changer le sens, ni ajouter d'interprétation personnelle. Votre production doit faire \textbf{environ 60 mots}.

\textbf{Paragraphe source (82 mots) :}
« La colonisation n'a pas seulement été une occupation militaire ou administrative. Elle a constitué une entreprise systématique de déstructuration des sociétés africaines : accaparement des terres fertiles, imposition d'une économie de rente au profit de la métropole, destruction des chefferies traditionnelles pour les remplacer par des auxiliaires dociles, et surtout, tentative d'effacement des cultures et des langues au profit d'une "civilisation" présentée comme supérieure. L'école coloniale fut l'outil principal de cette aliénation mentale, formant des élites coupées de leurs racines, prêtes à administrer l'exploitation de leurs propres frères. »

\begin{enumerate}
    \item Produisez votre \textbf{reformulation}.
    \item Comptez les mots de votre production.
    \item En vous mettant dans la peau d'un correcteur, listez \textbf{trois critères} sur lesquels vous évaluerez votre propre texte (ex. : fidélité au sens, absence de copiage, correction linguistique, concision).
\end{enumerate}
\end{exercice}

\begin{exercice}[Sujet de remédiation : Correction d'un résumé d'élève]
\textbf{Consigne :} Voici une copie d'élève répondant à l'exercice de contraction précédent (résumé de la réflexion de Meka, consigne : 150 mots $\pm$ 10\%). Ce résumé comporte \textbf{cinq erreurs types} (copiage, dépassement de limite, avis personnel, contresens sur ironie, désorganisation). Identifiez chaque erreur, citez le passage fautif, expliquez pourquoi c'est une erreur et proposez la correction.

\textbf{Proposition d'élève (198 mots) :}
« Meka est en prison et il regrette tout. Il regarde sa médaille qui est dans sa poche. Il pense au commandant qui a menti avec son sourire faux. Il a donné sa terre pour la médaille et maintenant il est arrêté pour avoir demandé un reçu. C'est injuste. Il pense à sa femme Kelara et à ses fils qui sont partis. L'école et l'église ont menti aussi. Le catéchiste a dit "Rends à César" mais César a tout pris. Meka a de la colère. Il comprend qu'il a été bête de croire le blanc. La médaille c'est un collier de chien. Il a perdu sa terre pour toujours. Moi je trouve que c'est triste et que les blancs sont méchants. Meka est un héros tragique. Le geôlier ouvre la porte. »

\begin{enumerate}
    \item Relevez les \textbf{5 erreurs} en complétant le tableau ci-dessous.
    \item Réécrivez une \textbf{version corrigée} respectant la consigne (135-165 mots, neutre, structuré).
\end{enumerate}

\begin{center}
\begin{tabularx}{\linewidth}{|l|X|X|X|}\hline
\textbf{N°} & \textbf{Passage fautif (citation)} & \textbf{Type d'erreur} & \textbf{Explication et correction} \\ \hline
1 & « Meka est en prison et il regrette tout. Il regarde sa médaille qui est dans sa poche. » & \textbf{Copiage / Paraphrase trop proche} & \textbf{Explication :} Reprise quasi littérale du texte original (« Assis sur son lit... la médaille... gisait maintenant au fond de sa poche »). \textbf{Correction :} Reformuler : « En cellule, Meka contemple la médaille, symbole de sa honte. » \\ \hline
2 & « Moi je trouve que c'est triste et que les blancs sont méchants. Meka est un héros tragique. » & \textbf{Avis personnel / Subjectivité} & \textbf{Explication :} Le résumé doit être \textbf{neutre} (zéro focalisation). L'élève juge (« triste », « méchants », « héros »). \textbf{Correction :} Supprimer tout jugement ; restituer la prise de conscience de Meka : « Il réalise sa complicité et la perte définitive de sa terre. » \\ \hline
3 & (Ensemble du texte : 198 mots) & \textbf{Dépassement de la limite} & \textbf{Explication :} Consigne : 150 mots $\pm$ 10\% (135-165). La copie fait 198 mots (+33\%). \textbf{Correction :} Condenser drastiquement (supprimer détails, généraliser, condenser phrases). \\ \hline
4 & « Il comprend qu'il a été bête de croire le blanc. La médaille c'est un collier de chien. » & \textbf{Contresens sur l'ironie / Style familier} & \textbf{Explication :} « Bête » est familier et subjectif. « Collier de chien » est une métaphore du texte qu'il faut intégrer dans une phrase neutre, pas citer brut. \textbf{Correction :} « Il perçoit l'ironie de son sort : la médaille, loin d'être un honneur, symbolise son asservissement (collier de chien). » \\ \hline
5 & (Structure générale : phrases courtes, juxtaposées, sans connecteurs logiques) & \textbf{Désorganisation / Manque de cohérence} & \textbf{Explication :} Pas de paragraphes, pas de connecteurs (cause, conséquence, opposition), ordre chronologique mais sans articulation logique. \textbf{Correction :} Structurer en 3 paragraphes selon le plan (Symbole retourné / Engrenage injuste / Lucidité finale) avec connecteurs. \\ \hline
\end{tabularx}
\end{center}
\end{exercice}




\newpage

\coursec{Corrigés des exercices}

\begin{corrige}[Exercice 1 — QCM : techniques de contraction]
\begin{enumerate}
    \item \textbf{C (Généralisation)} : on passe d'une série d'\cle{hyponymes} (lions, éléphants, girafes, zèbres) à un \cle{hyperonyme} (« la faune sauvage »). C'est le passage du particulier au général.
    \item \textbf{B (Suppression)} : on élimine les éléments secondaires (adjectifs qualificatifs, compléments circonstanciels, répétitions) sans toucher au noyau du message.
    \item \textbf{A (Condensation)} : plusieurs phrases sont fondues en une seule, plus dense, grâce à des subordonnées ou des appositions.
    \item \textbf{D (Reformulation)} : on réécrit l'idée avec ses propres mots (paraphrase) en conservant rigoureusement le sens.
\end{enumerate}
\end{corrige}

\begin{attention}[Point de vigilance]
Ne pas confondre \textbf{condensation} et \textbf{reformulation} : la condensation réduit le \emph{nombre de phrases}, la reformulation change les \emph{mots} à information égale.
\end{attention}

\begin{corrige}[Exercice 2 — Vrai / Faux : ironie et paradoxe]
\begin{enumerate}
    \item \textbf{Faux.} L'ironie consiste le plus souvent à dire le contraire de ce que l'on pense, mais pas \emph{toujours} dans un but moqueur : il existe l'ironie dramatique (le lecteur en sait plus que le personnage) et l'ironie du sort. Le mot « toujours » rend l'affirmation fausse.
    \item \textbf{Faux.} Affirmer une chose et son contraire dans la même phrase définit l'\cle{oxymore} (ex. : « cette obscure clarté »). Le \cle{paradoxe} est une affirmation \emph{en apparence} contradictoire ou absurde qui, à la réflexion, révèle une vérité profonde (ex. : « L'enfant est père de l'homme »).
    \item \textbf{Vrai.} Chez Oyono, l'ironie crée un décalage entre le discours officiel (honneur, civilisation, reconnaissance) et la réalité vécue (spoliation, mépris, arrestation arbitraire). Le système colonial est dénoncé par le rire amer, sans être nommé frontalement.
    \item \textbf{Faux.} Prendre un énoncé ironique au premier degré, c'est prendre pour argent comptant ce que l'auteur moque : la critique devient un éloge. C'est un \textbf{contresens total}, faute grave dans un résumé.
\end{enumerate}
\end{corrige}

\begin{attention}[Point de vigilance]
Au Bac, justifier un Vrai/Faux exige de citer la \textbf{définition exacte} du cours et de montrer en quoi l'affirmation s'en écarte (mot trop absolu : « toujours », « jamais », « seulement »).
\end{attention}

\begin{corrige}[Exercice 3 — Définitions : constatifs, performatifs, variétés]
\begin{enumerate}
    \item \textbf{Énoncé constatif :} énoncé qui \emph{décrit} un état de fait et qui peut être jugé \textbf{vrai} ou \textbf{faux}. \emph{Exemple original :} « Le fleuve Wouri se jette dans l'océan Atlantique. »
    \item \textbf{Énoncé performatif explicite :} énoncé qui, par sa seule prononciation dans un contexte conventionnel approprié, \emph{accomplit} l'action qu'il énonce. Il contient un verbe performatif à la première personne du présent de l'indicatif. Il n'est ni vrai ni faux : il est \textbf{réussi} ou \textbf{raté}. \emph{Exemple original :} « Je vous déclare mari et femme » (prononcé par le maire).
    \item \textbf{Variété diastratique :} variété de langue liée à l'appartenance du locuteur à un groupe social ou professionnel. \emph{Exemples camerounais :} le camfranglais (parler urbain des jeunes), le jargon des chauffeurs de taxi (« le volant », « la recette »), le français administratif des fonctionnaires.
\end{enumerate}
\end{corrige}

\begin{attention}[Point de vigilance]
Un performatif n'est réussi que si les \textbf{conditions d'application} sont remplies : bon locuteur (autorité compétente), bon contexte (cérémonie), bonne formule. « Je vous nomme chef de canton » dit par un passant est un performatif \emph{raté}.
\end{attention}

\begin{corrige}[Exercice 4 — Repérage : structure argumentative]
\begin{itemize}
    \item \textbf{Thèse (implicite) :} la médaille est une récompense dérisoire qui symbolise la spoliation et la perte d'autonomie de Meka.
    \item \textbf{Argument :} « il avait donné ses terres » (le fait fondateur qui justifie le sentiment de dépossession).
    \item \textbf{Connecteur de cause :} \textbf{car}.
    \item \textbf{Connecteurs de conséquence :} \textbf{donc}, \textbf{c'est pourquoi}, \textbf{alors}.
    \item \textbf{Connecteur d'opposition :} \textbf{mais}.
\end{itemize}
\end{corrige}

\begin{attention}[Point de vigilance]
Distinguer \textbf{thèse} (l'idée défendue, souvent implicite dans un texte littéraire) et \textbf{argument} (la raison qui la soutient). Un connecteur n'est pas un argument : il articule les idées.
\end{attention}

\begin{corrige}[Exercice 5 — Repérage et explication de l'ironie]
\begin{enumerate}
    \item \textbf{Trois énoncés ironiques :} (a) « la vraie richesse, ce n'est pas la terre qui nourrit ses enfants, c'est ce petit bout de métal » ; (b) « Il a eu le courage de se débarrasser de l'inutile pour garder l'essentiel » ; (c) « La France vous remercie, Meka, d'avoir si bien servi ses intérêts. »
    \item \textbf{Sens littéral / sens réel :}
    \begin{center}
    \begin{tabularx}{\linewidth}{|X|X|X|}\hline
    \rowcolor{popblueL}\textbf{Énoncé} & \textbf{Sens littéral (dit)} & \textbf{Sens réel (visé)} \\ \hline
    (a) & La médaille vaut plus que la terre. & La terre est vitale ; la médaille est dérisoire et trompeuse. \\ \hline
    (b) & Vendre sa terre est courageux et lucide. & C'est une aliénation, une folie imposée par le colon. \\ \hline
    (c) & La France est reconnaissante. & La France exploite Meka : le « service » est une spoliation. \\ \hline
    \end{tabularx}
    \end{center}
    \item \textbf{Risque de contresens :} un résumé au premier degré écrirait : « Meka a sagement échangé sa terre inutile contre une médaille précieuse et la France l'en a remercié. » Cela \textbf{inverse le sens de l'œuvre} : on valide la propagande coloniale au lieu de restituer la dénonciation d'Oyono. Le résumé doit transmettre la \emph{vision critique de l'auteur}, non le mensonge du personnage.
\end{enumerate}
\end{corrige}

\begin{attention}[Point de vigilance]
Dans un résumé, toute ironie doit être \textbf{explicitée} (ex. : « Meka perçoit l'ironie de son sort : la médaille, loin d'être un honneur, symbolise son asservissement ») et jamais reprise au premier degré.
\end{attention}

\begin{corrige}[Exercice 6 — Classement : constatifs / performatifs]
\begin{center}
\begin{tabularx}{\linewidth}{|l|X|X|}\hline
\rowcolor{popblueL}\textbf{N°} & \textbf{Type} & \textbf{Justification} \\ \hline
1 & Performatif & « Je déclare » : dire, c'est faire — l'énoncé crée la qualité de chevalier. \\ \hline
2 & Constatif & Description météorologique vérifiable (vrai/faux). \\ \hline
3 & Performatif & « Je vous prie » : l'énoncé est lui-même l'acte de prier. \\ \hline
4 & Constatif & Récit d'un événement passé, vérifiable. \\ \hline
5 & Performatif & « Je vous nomme » : l'énoncé installe dans la fonction. \\ \hline
6 & Constatif & Description matérielle de l'objet, vérifiable. \\ \hline
7 & Performatif & « Je promets » : l'énoncé engage le locuteur. \\ \hline
8 & Constatif & Description d'une compétence, vérifiable. \\ \hline
9 & Performatif & « Je baptise » : l'énoncé donne le nom au pont. \\ \hline
10 & Constatif & Affirmation sur une action en cours, vérifiable. \\ \hline
\end{tabularx}
\end{center}
\textbf{Critère de décision :} verbe à la 1\textsuperscript{re} personne du présent + l'énonciation accomplit l'action $\Rightarrow$ performatif ; sinon, énoncé descriptible en vrai/faux $\Rightarrow$ constatif.
\end{corrige}

\begin{attention}[Point de vigilance]
Certains performatifs sont \emph{implicites} (ex. : « Je serai là demain » $\Rightarrow$ promesse). Au Bac, on privilégie le performatif \emph{explicite} (verbe performatif apparent) pour l'identification sûre.
\end{attention}

\begin{corrige}[Exercice 7 — Réduction ciblée (8 mots maximum)]
\begin{enumerate}
    \item « Meka, vieux, marcha péniblement vers le commandant. » (\textbf{S} : suppression de « courbé sous le poids des ans », de « la case du »). — 7 mots.
    \item « La foule villageoise observait la scène en silence. » (\textbf{G} : « villageois, femmes, enfants, vieillards » $\to$ « la foule villageoise » ; « silence total » $\to$ « silence »). — 8 mots.
    \item « La petite médaille dorée fut épinglée sur Meka. » (\textbf{S} : suppression de « brillante, ronde » et de « la poitrine usée du récipiendaire »). — 8 mots.
    \item « Le long discours du commandant fut incompréhensible. » (\textbf{C} : « plein de mots difficiles que personne ne comprenait vraiment » condensé en « incompréhensible »). — 7 mots.
    \item « Meka rentra chez lui, le cœur lourd. » (\textbf{C} : fusion de « le cœur lourd, l'esprit troublé par l'ironie du sort »). — 7 mots.
\end{enumerate}
\end{corrige}

\begin{attention}[Point de vigilance]
Vérifier le \textbf{comptage des mots} et conserver l'\textbf{information essentielle} (qui ? quoi ? où ?). Une réduction qui fait disparaître l'agent ou l'action principale est une mutilation, pas une contraction.
\end{attention}

\begin{corrige}[Exercice 8 — Plan et résumé (vie scolaire)]
\begin{enumerate}
    \item \textbf{Thèse :} malgré une volonté politique affichée et des avancées concrètes, l'école technique camerounaise reste freinée par des obstacles structurels ; sa réussite exige suivi rigoureux et dialogue social apaisé.
    \item \textbf{Plan en trois parties :}
    \begin{enumerate}
        \item \emph{Une rentrée résiliente portée par l'effort de l'État} : retour massif des élèves malgré les grèves ; construction de salles en zone rurale ; numérique dans les lycées techniques ; réforme des programmes.
        \item \emph{Des obstacles persistants à l'égalité des chances} : pénurie d'enseignants qualifiés dans les filières industrielles ; coût élevé des fournitures techniques ; inégalités d'accès à l'électricité et à internet.
        \item \emph{Les conditions de la réussite} : la volonté affichée ne suffit pas ; il faut un suivi rigoureux et un dialogue social apaisé.
    \end{enumerate}
    \item \textbf{Résumé (58 mots) :}
    \begin{quote}
    L'année scolaire 2024-2025 s'ouvre dans la résilience malgré les grèves. L'État construit des salles rurales, introduit le numérique et réforme les programmes techniques. Cependant, le manque d'enseignants qualifiés et le coût des fournitures freinent l'égalité des chances. La volonté politique exige un suivi rigoureux et un dialogue social apaisé pour garantir une éducation de qualité pour tous.
    \end{quote}
    \textbf{Comptage :} 58 mots $\Rightarrow$ dans la fourchette 54--66.
\end{enumerate}
\end{corrige}

\begin{attention}[Point de vigilance]
Respecter scrupuleusement la fourchette de mots ($\pm$ 10\,\%). Compter \emph{chaque} mot (mots-outils inclus). Indiquer le total en fin de résumé.
\end{attention}

\begin{corrige}[Exercice 9 — Sujet type Bac : contraction (réflexion de Meka)]
\begin{enumerate}
    \item \textbf{Dépouillement des idées essentielles :}
    \begin{enumerate}
        \item En cellule, Meka contemple la médaille devenue symbole de honte.
        \item Retour sur la cérémonie : sourire hypocrite du commandant, « loyaux services ».
        \item La cession de la terre ancestrale (nourricière et sacrée) contre un bout de métal.
        \item L'arrestation arbitraire pour avoir demandé un reçu ; la parole du blanc prime sur l'écrit.
        \item Pensées vers Kelara et ses fils réfugiés en brousse.
        \item L'Église et l'école dénoncées comme instruments d'aliénation (« Rends à César... »).
        \item La colère née de sa propre complicité : il a cru jouer, il a été joué.
        \item La médaille = collier de chien ; la vraie médaille, la terre, est perdue à jamais ; la porte se referme.
    \end{enumerate}
    \item \textbf{Plan en trois parties :}
    \begin{enumerate}
        \item \emph{Le symbole retourné} : de la fierté à la honte (cellule, médaille lourde, cérémonie mensongère, spoliation).
        \item \emph{L'engrenage de l'injustice coloniale} : arrestation arbitraire, arbitraire de la parole du colon, complicité de l'Église et de l'école.
        \item \emph{La lucidité tragique} : colère contre soi, médaille-collier, terre perdue définitivement, porte claquée.
    \end{enumerate}
    \item \textbf{Résumé (148 mots) :}
    \begin{quote}
    Enfermé au poste de police, Meka contemple la médaille qui gît dans sa poche, lourde d'un secret honteux. Il revit la cérémonie : le sourire faux du commandant, les louanges sur ses « loyaux services » masquant la spoliation de sa terre ancestrale, unique patrimoine familial. Arrêté pour avoir exigé un reçu de vente, il mesure l'absurdité : la parole du colon n'a pas besoin d'écrit. Ses pensées vont vers Kelara, ses fils fuyant l'uniforme, et vers l'école et l'Église, instruments d'une aliénation consentie. Une colère sourde naît de sa propre complicité : il a cru jouer le jeu de l'occupant, il en a été la victime. La médaille n'est pas un honneur, mais un collier de chien. La vraie médaille, la terre des ancêtres, est définitivement perdue. Le claquement de la porte du geôlier scelle cette lucidité amère.
    \end{quote}
    \textbf{Total : 148 mots} (fourchette 135--165 respectée).
\end{enumerate}
\end{corrige}

\begin{attention}[Barème indicatif (sur 20) et points de vigilance]
Respect de la limite de mots (4 pts) ; fidélité aux idées essentielles et à leur ordre (6 pts) ; neutralité (aucun avis personnel, aucune citation longue) (4 pts) ; reformulation sans copiage (3 pts) ; correction linguistique et connecteurs (3 pts). \textbf{Vigilance :} compter et \emph{indiquer} le nombre de mots ; ne pas reprendre les formules ironiques au premier degré ; ne pas écrire « je » ni juger les personnages.
\end{attention}

\begin{corrige}[Exercice 10 — Reformulation et auto-évaluation]
\begin{enumerate}
    \item \textbf{Reformulation (59 mots) :}
    \begin{quote}
    La colonisation ne se limita pas à l'occupation militaire ; elle déstructura systématiquement les sociétés africaines : terres confisquées, économie organisée au profit de la métropole, chefferies remplacées par des auxiliaires dociles, cultures et langues étouffées. L'école coloniale, principal outil de cette aliénation, forma des élites déracinées, chargées d'administrer l'exploitation de leurs propres frères.
    \end{quote}
    \item \textbf{Nombre de mots :} 59 (consigne « environ 60 mots » respectée).
    \item \textbf{Trois critères d'auto-évaluation :}
    \begin{enumerate}
        \item \textbf{Fidélité au sens :} toutes les idées forces (déstructuration, spoliation foncière, économie de rente, pouvoir, effacement culturel, rôle de l'école, élites) sont présentes, sans ajout ni omission.
        \item \textbf{Reformulation effective :} vocabulaire et syntaxe personnels (« ne se limita pas » pour « n'a pas seulement été » ; « étouffées » pour « tentative d'effacement » ; « déracinées » pour « coupées de leurs racines »).
        \item \textbf{Concision et correction :} contrainte quantitative respectée, phrases grammaticalement correctes, enchaînement logique assuré par la ponctuation et les connecteurs.
    \end{enumerate}
\end{enumerate}
\end{corrige}

\begin{attention}[Point de vigilance]
La reformulation n'est ni un résumé (on garde \emph{toutes} les idées) ni un commentaire (aucune interprétation personnelle). Toute reprise de plus de trois mots consécutifs du texte source s'apparente à du copiage.
\end{attention}

\begin{corrige}[Exercice 11 — Remédiation : correction d'un résumé d'élève]
\begin{enumerate}
    \item \textbf{Les cinq erreurs :}
    \begin{center}
    \begin{tabularx}{\linewidth}{|l|X|X|X|}\hline
    \rowcolor{popblueL}\textbf{N°} & \textbf{Passage fautif} & \textbf{Type d'erreur} & \textbf{Explication et correction} \\ \hline
    1 & « Meka est en prison et il regrette tout. Il regarde sa médaille qui est dans sa poche. » & Copiage / paraphrase trop proche & Reprise quasi littérale du texte (« la médaille... gisait au fond de sa poche »). \emph{Correction :} « En cellule, Meka contemple la médaille, devenue symbole de honte. » \\ \hline
    2 & « Moi je trouve que c'est triste et que les blancs sont méchants. Meka est un héros tragique. » & Avis personnel / subjectivité & Le résumé doit être \textbf{neutre} : aucun jugement (« triste », « méchants », « héros »). \emph{Correction :} supprimer ; restituer la prise de conscience du personnage. \\ \hline
    3 & Ensemble de la copie : 198 mots & Dépassement de la limite & Consigne : 135--165 mots ; 198 mots, soit +33\,\%. \emph{Correction :} supprimer les redites, généraliser, condenser. \\ \hline
    4 & « Il comprend qu'il a été bête de croire le blanc. La médaille c'est un collier de chien. » & Contresens sur l'ironie / registre familier & « Bête » est familier et subjectif ; la métaphore est plaquée sans recul. \emph{Correction :} « Il perçoit l'ironie de son sort : la médaille, loin d'être un honneur, symbolise son asservissement. » \\ \hline
    5 & Phrases courtes juxtaposées, sans paragraphes ni connecteurs & Désorganisation / incohérence & Aucune articulation logique. \emph{Correction :} structurer en trois moments (symbole retourné / engrenage injuste / lucidité finale) reliés par des connecteurs. \\ \hline
    \end{tabularx}
    \end{center}
    \item \textbf{Version corrigée (152 mots) :}
    \begin{quote}
    Enfermé au poste de police, Meka contemple la médaille, jadis désirée, devenue le symbole honteux de sa duperie. Il revit la cérémonie : le sourire hypocrite du commandant et les louanges sur ses « loyaux services » masquaient la spoliation de sa terre ancestrale, nourricière et sacrée. Arrêté pour avoir réclamé un reçu de vente, il découvre l'arbitraire colonial : la parole du maître prime sur le droit écrit. Son esprit rejoint sa femme Kelara, ses fils réfugiés en brousse, et dénonce la complicité de l'Église et de l'école, instruments d'aliénation. Une colère sourde monte, non contre la prison, mais contre sa propre complicité : il a cru jouer le jeu du colon, il en a été la proie. La médaille n'est pas une récompense, c'est un collier de chien. La vraie richesse, la terre des pères, est perdue à jamais. Le claquement de la porte du geôlier scelle cette lucidité définitive.
    \end{quote}
    \textbf{Total : 152 mots} (fourchette respectée).
\end{enumerate}
\end{corrige}

\begin{attention}[Barème indicatif (sur 20) et points de vigilance]
Identification des 5 erreurs (5 $\times$ 2 pts) ; version corrigée : limite de mots et comptage indiqué (3 pts), neutralité (2 pts), ordre des idées et connecteurs (3 pts), correction linguistique (2 pts). \textbf{Vigilance :} citer \emph{précisément} le passage fautif avant de le corriger ; une erreur relevée sans correction proposée ne vaut que la moitié des points.
\end{attention}

\newpage

\coursec{Fiche bilan}

\courssub{Carte mentale de la leçon}
\begin{popfigure}
\begin{tikzpicture}[
    mindmap,
    grow cyclic,
    every node/.style={concept, execute at begin node=\hskip0pt},
    concept color=popblue,
    level 1/.append style={level distance=4.5cm, sibling angle=360/7},
    level 2/.append style={level distance=3cm, sibling angle=45},
    text=white,
    font=\small\bfseries
]
\node[concept, scale=1.2, align=center]{Contraction de texte\\Le Vieux Nègre et la Médaille}
    child[concept color=poporange] { node[concept, align=center]{1. Lecture\\méthodique\\de l'œuvre} }
    child[concept color=popgreen] { node[concept, align=center]{2. Figures\\de pensée\\Ironie \& Paradoxe} }
    child[concept color=poppurple] { node[concept, align=center]{3. Langue :\\Variétés lexicales\\Constatif / Performatif} }
    child[concept color=poppink] { node[concept, align=center]{4. Structure\\argumentative\\Thèse, arguments, conclusion} }
    child[concept color=popgold] { node[concept, align=center]{5. Techniques\\de réduction\\Suppression, généralisation} }
    child[concept color=popdark] { node[concept, align=center]{6. Reformulation\\et rédaction\\Synthèse cohérente} }
    child[concept color=popblue] { node[concept, align=center]{7. Évaluation\\Compte-rendu\\Remédiation} };
\end{tikzpicture}
\legende{Carte mentale récapitulative de la leçon 09 : axes de travail sur l'œuvre et la contraction.}
\end{popfigure}

\courssub{Bilan des savoirs et savoir-faire}
\begin{bilan}

\textbf{1. Réception des textes : \emph{Le Vieux Nègre et la Médaille} (Ferdinand Oyono)}
\begin{itemize}
    \item \cle{Lecture méthodique} : Contextualisation (colonialisme, assimilation), analyse du personnage (Meka), structure narrative (quête de la médaille $\rightarrow$ désillusion).
    \item \cle{Exposés} : Thèmes majeurs (aliénation, dignité, choc des cultures), rôle de l'ironie dramatique, portée satirique.
\end{itemize}

\textbf{2. Langue française : Stylistique, Sémantique, Communication}
\begin{center}
\begin{tabularx}{\linewidth}{|l|X|X|}\hline
\rowcolor{popblueL}
\textbf{Notion} & \textbf{Définition / Repérage} & \textbf{Effet / Fonction} \\ \hline
\textbf{Ironie} & Dire le contraire de ce que l'on pense (ton moqueur). & Dénonciation, distance critique, humour amer. \\ \hline
\textbf{Paradoxe} & Affirmation en apparence contradictoire mais vraie en profondeur. & Provocation intellectuelle, révélation d'une vérité cachée. \\ \hline
\textbf{Variétés lexicales} & Français standard, familier, régional (camfranglais), technique. & Caractérisation sociale, réalisme, expressivité. \\ \hline
\textbf{Énoncé constatif} & Décrit un état de choses, vérifiable (vrai/faux). & Information, description, narration. \\ \hline
\textbf{Énoncé performatif} & \textbf{Agit} en se disant (ex: « Je vous déclare mari et femme »). & Action sociale, engagement, changement d'état. \\ \hline
\end{tabularx}
\end{center}

\textbf{3. Production des textes : La contraction de texte (Résumé / Contraction)}
\begin{methode}[Étapes de la contraction]
\begin{enumerate}
    \item \textbf{Lecture active} : Identifier le sujet, la thèse, le plan, les arguments clés, les connecteurs logiques.
    \item \cle{Découpage} : Segmenter le texte en séquences logiques (macro-structure).
    \item \cle{Sélection} : Ne garder que les idées \textbf{essentielles} (thèse + arguments principaux). Éliminer exemples, détails, répétitions, digressions.
    \item \cle{Réduction linguistique} :
        \begin{itemize}
            \item \textbf{Suppression} : Adjectifs qualificatifs, compléments circonstanciels secondaires, phrases illustratifs.
            \item \textbf{Généralisation} : Remplacer une énumération par un hyperonyme (ex: « chiens, chats, oiseaux » $\rightarrow$ « animaux »).
            \item \textbf{Reformulation} : Recrire avec ses mots (synonymes, changement de voix, nominalisation/verbalisation) sans trahir le sens.
        \end{itemize}
    \item \textbf{Assemblage} : Relier les idées par des connecteurs logiques (donc, cependant, par conséquent, en outre).
    \item \textbf{Vérification} : Respect du nombre de mots ($\pm 10\%$), fidélité au sens, cohérence, correction linguistique.
\end{enumerate}
\end{methode}

\textbf{4. Structure argumentative type (à repérer pour bien résumer)}
\begin{itemize}
    \item \cle{Thèse} : Position défendue (souvent en début ou fin de texte).
    \item \cle{Arguments} : Raisons logiques, faits, autorités, analogies.
    \item \cle{Exemples/Illustrations} : Concrets, à \textbf{supprimer} ou \textbf{condenser} dans le résumé.
    \item \cle{Conclusion} : Ouverture, conséquence, recommandation.
\end{itemize}

\end{bilan}

\courssub{Checklist avant le Bac}
\begin{aretenir}[Checklist Contraction de texte -- Terminale Tech]
    $\square$ Je sais définir la contraction de texte et la distinguer du résumé simple (taux de réduction imposé).
    $\square$ Je repère l'ironie et le paradoxe dans \emph{Le Vieux Nègre et la Médaille} et j'explique leur effet.
    $\square$ Je distingue une variété lexicale d'une autre (standard / familier / régional) dans un extrait.
    $\square$ Je différencie un énoncé constatif (description) d'un énoncé performatif (action).
    $\square$ J'identifie la structure argumentative : thèse, arguments, connecteurs, conclusion.
    $\square$ J'applique les 3 techniques de réduction : suppression, généralisation, reformulation.
    $\square$ Je reformule sans copier-coller : changement de syntaxe, vocabulaire propre, voix passive/active.
    $\square$ Je respecte la consigne de longueur (nombre de mots $\pm 10\%$).
    $\square$ Mon texte contracté est cohérent, connecté logiquement et exempt de fautes de français.
    $\square$ Je gère mon temps : Lecture (15\%), Analyse/Plan (25\%), Rédaction (50\%), Relecture (10\%).
    $\square$ Je connais les critères d'évaluation : Fidélité, Concision, Cohérence, Correction.
\end{aretenir}

\findecours

\end{document}
